\section{Kodaira--Spencer I: Harris--Taylor families and Hodge data}\label{sec:pf:ks1}

\providecommand{\ksoneAfp}{\A_f^{p}}
\providecommand{\ksoneH}[3]{H^{#3}(#1,\Omega^{#2})}
\providecommand{\ksoneSp}{\operatorname{Sp}}
\providecommand{\ksonefs}{\textstyle\sum\nolimits^{\mathrm f}}

\subsection{Introduction}\label{pf:ks1:intro}

This section contains the proofs of two results which are used, together with a fixed point
formula and the holomorphic Lefschetz formula, in the proof of the geometric package
\lref{UpperTheory.Gal.Int.Geo.kodaira_spencer_of_weight_one}. The name of the section refers to
the Kodaira--Spencer computation of the relative dimension of the moduli spaces (\cref{pf:ks1:sec-smooth}) and to
the Kuga--Sato towers whose Hodge cohomology is studied in the second half.

The first result (\cref{pf:ks1:thm-models}) is the existence of integral models of unitary
PEL moduli problems of the type considered by Kottwitz \cite{Kottwitz1992} and Harris--Taylor
\cite{HarrisTaylor2001}: for every Harris--Taylor datum $D$ at a prime $p$, every Noetherian ring $R$
carrying base data for $D$ (in particular every integer prime to $p$ is a unit in $R$) and every
level $N\ge3$ prime to $p$, the moduli problem of abelian schemes with $D$-structure and level-$N$
structure is represented by a quasi-projective $R$-scheme which is smooth of relative dimension
$a(n-1)$. This is classical; we give the precise formal statement, the architecture of its formal
proof, and complete proofs of its steps following the formal proof: the choice of the auxiliary level,
the reduction to relative representability over a Siegel moduli space, the representability of actions and
level maps, the closedness of Kottwitz's condition, the descent of the auxiliary level, and the deformation
theory which gives smoothness and the relative dimension. The representability of Siegel moduli spaces and
Grothendieck--Messing theory are cited from the literature.

The second result (\cref{pf:ks1:thm-hodge}) is new in its formulation. It says that the
holomorphic Lefschetz formula (\cref{pf:ks1:def-hlf}) and a fixed point formula for twisted Hecke
correspondences on Kuga--Sato varieties (\cref{pf:ks1:def-fixk}, proved in
\lref{UpperTheory.Gal.FixZ.FixK_holds}) imply the proposition $\sfname{KSHodge}$
(\cref{pf:ks1:def-kshodge}): if a system $\theta$ of Hecke eigenvalues occurs in algebraic modular
forms of weight $\lambda$ on the definite inner forms of a unitary similitude group of signature
$(1,n-1)$ (detected by traces of good test functions), then for a suitable exponent $k$ depending
only on $\lambda$, and for each Hodge exponent $h_i$ of $\lambda$, the eigensystem $\sigma\circ\theta$
occurs in the coherent cohomology $\ksoneH{Z}{h_i}{b}$ of the $k$-fold Kuga--Sato variety $Z$ over the
corresponding unitary Shimura variety, for some $b$ and every $\sigma\in\Aut(\C)$. The argument has five
steps.
\begin{enumerate}[label=(\roman*)]
\item (Geometric side.) For a Hecke operator $T_h$ composed with the pull-back by a fibrewise
  multiplication $[\alpha]$, $\alpha\in\OO_F^k$, the alternating trace on Hodge cohomology is computed by
  the holomorphic Lefschetz formula, and the fixed point formula identifies the result with a weighted
  sum $\Sigma(c_\alpha;\dots)$ over complex multiplication points (\cref{pf:ks1:prop-geostep}).
\item (Counting identities.) Such weighted sums are equal to sums over conjugacy classes of definite
  similitude pairs of root sums of the weight (\cref{pf:ks1:prop-h1}, a comparison in the style of
  Kottwitz \cite{Kottwitz1992}); and sums over definite pairs are traces of Hecke operators on algebraic
  modular forms in the sense of Gross \cite{Gross1999} (\cref{pf:ks1:prop-h2}, the trace formula for a group
  which is compact at infinity).
\item (Lefschetz weights.) A finite linear combination $\sum_\alpha c_\alpha[\alpha]$, in the spirit of
  Scholl's projectors \cite{Scholl1990}, turns the local Lefschetz weights into the Schur polynomial of
  $\lambda$ times $\sum_iy^{h_i}$ (\cref{pf:ks1:prop-rootw}); this rests on an interpolation
  argument and on two identities for Schur polynomials (\cref{pf:ks1:sec-schur}).
\item (Hodge traces.) Combining (i)--(iii), the alternating trace of a good Hecke operator composed with
  $\sum_\alpha c_\alpha[\alpha]^*$ on $\bigoplus_{a,b}\ksoneH{Z}{a}{b}$, as a polynomial in the
  variable $y$ marking the degree $a$, equals the definite trace times $\sum_iy^{h_i}$
  (\cref{pf:ks1:prop-hodgetr}).
\item (Eigenvectors.) Comparing coefficients of $y^{h_i}$ and an elementary lemma on commuting operators give
  simultaneous eigenvectors in $\ksoneH{Z}{h_i}{b}$ (\cref{pf:ks1:prop-hodgeeig}); base change and Galois
  conjugation give $\sfname{KSHodge}$ (\cref{pf:ks1:prop-kshodge}).
\end{enumerate}

\paragraph{Organisation.} \Cref{pf:ks1:sec-data,pf:ks1:sec-rep,pf:ks1:sec-repsteps,pf:ks1:sec-smooth,pf:ks1:sec-deform} treat the integral models.
\Cref{pf:ks1:sec-towers} recalls Hecke towers and their Hodge cohomology and states the Lefschetz
formula. \Cref{pf:ks1:sec-groups} introduces the group-theoretic objects (forms, fixed point sums,
pairs, definite traces) and states $\sfname{FixK}$, $\sfname{KSHodge}$ and \cref{pf:ks1:thm-hodge}.
\Cref{pf:ks1:sec-counting} proves the two counting identities, \cref{pf:ks1:sec-schur} the identities for
Lefschetz weights, and \cref{pf:ks1:sec-proof} assembles the proof of \cref{pf:ks1:thm-hodge}.
\Cref{pf:ks1:sec-remarks} records peculiarities of the formal statements.

\paragraph{Conventions.}
Rings are commutative unless they are explicitly algebras such as $C$ or $O$ below; $\ksoneSp R$ denotes $\Spec R$.
For a function $u$ on a set $X$ we write $\ksonefs_{x\in X}u(x)$ for the sum of $u$ over its support if
the support is finite, and $0$ otherwise; this is the convention of the formal development (Mathlib's
\texttt{finsum}) and it is used in several definitions below, which is why finiteness statements appear
as separate hypotheses or lemmas. Similarly $|S|$ denotes the cardinality of a finite set $S$ and
$|S|:=0$ if $S$ is infinite, and $0^{-1}:=0$. For subgroups $H,K$ of a group, $[K:H\cap K]$ denotes the index,
taken to be $0$ if it is infinite. For a subgroup $K$ of a group $H$ and
$g\in H$, $\mathbf 1_{KgK}$ is the indicator function of the double coset $KgK$, and the convolution of
$K$-invariant functions is $(\phi_1*\phi_2)(g)=\ksonefs_{x\in H/K}\phi_1(x)\phi_2(x^{-1}g)$.

% ---------------------------------------------------------------------------------------------
\subsection{Harris--Taylor data and their moduli problem}\label{pf:ks1:sec-data}

The following definitions are those of the formal development; they are faithful versions of the
PEL data and moduli problems of \cite[\S5]{Kottwitz1992} and \cite[Ch.~III]{HarrisTaylor2001}.

\begin{definition}[Harris--Taylor datum]\label{pf:ks1:def-datum}
Let $p$ be a prime. A \emph{Harris--Taylor datum} $D$ for $p$ consists of:
\begin{enumerate}[label=(D\arabic*)]
\item a CM field $F$, with complex conjugation $c\mapsto\bar c$ on $F$ and on $\OO_F$;
\item $\omega\in\OO_F$ and integers $\mathrm{tr},\mathrm{nm}$ with $\omega+\bar\omega=\mathrm{tr}$,
  $\omega\bar\omega=\mathrm{nm}$, $\mathrm{tr}^2-4\,\mathrm{nm}<0$, $p\nmid\mathrm{tr}^2-4\,\mathrm{nm}$,
  and $p\mid r^2-\mathrm{tr}\,r+\mathrm{nm}$ for some $r\in\Z$;
\item a ring $O$ which is an $\OO_F$-algebra, free of rank $m^2[F:\Q]$ as a $\Z$-module for some $m\ge1$,
  such that every central element of $O$ has a nonzero integral multiple in the image of $\OO_F$; a ring
  anti-automorphism $b\mapsto b^*$ of $O$ with $b^{**}=b$ restricting to complex conjugation on $\OO_F$; and an
  isomorphism of $\OO_F$-algebras $O\otimes\Z_p\cong M_m(\OO_F\otimes\Z_p)$;
\item an $O$-module $\Lambda$, free of rank $mn[F:\Q]$ over $\Z$ for some $n\ge1$, with an alternating
  form $E_\Lambda:\Lambda\times\Lambda\to\Z$ which is nondegenerate, satisfies $E_\Lambda(bx,y)=E_\Lambda(x,b^*y)$
  for $b\in O$, and is \emph{perfect up to an integer prime to $p$}: there is $e\ge0$ prime to $p$ such
  that every linear form $\Lambda\to\Z$ multiplied by $e$ is of the form $E_\Lambda(x,-)$;
\item an integer $a\ge0$, a root $\rho_c\in\C$ of $X^2-\mathrm{tr}\,X+\mathrm{nm}$ and distinct embeddings
  $\tau_0,\dots,\tau_{a-1}:F\to\C$ with $\tau_i(\omega)=\rho_c$;
\item (Kottwitz's determinant condition) an $\R$-linear endomorphism $J$ of $\Lambda_\R=\R\otimes\Lambda$
  with $J^2=-1$, $E_\Lambda(Jx,Jy)=E_\Lambda(x,y)$, $E_\Lambda(Jx,x)>0$ for $x\neq0$, commuting with $O$, such
  that for every $c\in\OO_F$ the multiplication $A_c$ by $c$ on $\Lambda_\R$ satisfies
  $\tr_\R(A_c)=2\Re S(c)$ and $\tr_\R(JA_c)=-2\Im S(c)$, where
  \[ S(c)=m\Bigl(\sum_{i<a}\tau_i(c)+n\!\!\sum_{\tau(\omega)=\bar\rho_c}\!\!\tau(c)-\sum_{i<a}\overline{\tau_i(c)}\Bigr).\]
\end{enumerate}
$D$ is of \emph{compact type} if $2a<[F:\Q]$, or $O\otimes\Q$ is a division algebra and $n<2m$.
\end{definition}

Condition (D6) says that $\tr_\C(c\mid(\Lambda_\R,J))=S(c)$: an embedding $\tau$ of $F$ occurs in the
complex vector space $(\Lambda_\R,J)$ with multiplicity $m$ if $\tau\in\{\tau_i\}$, $m(n-1)$ if
$\tau\in\{\bar\tau_i\}$, $mn$ if $\tau(\omega)=\bar\rho_c$ and $\tau\notin\{\bar\tau_i\}$, and $0$ otherwise.
The associated unitary group has signature $(1,n-1)$ at the places $\tau_i$ and is definite at the other
places; for $a=1$ and $O$ a maximal order of a division algebra this is the setting of
\cite{HarrisTaylor2001}.

\begin{definition}[Base data]\label{pf:ks1:def-base}
Let $D$ be a datum for $p$ and $R$ a ring. A \emph{base} $b$ for $D$ over $R$ consists of $\rho\in R$ with
$\rho^2-\mathrm{tr}\,\rho+\mathrm{nm}=0$ and ring homomorphisms $t_i:\OO_F\to R$ ($i<a$) with $t_i(\omega)=\rho$,
such that every integer prime to $p$ is a unit of $R$; for $i\neq j$ the elements $t_i(x)-t_j(x)$,
$x\in\OO_F$, generate the unit ideal; and for $i\neq j$ and every homomorphism $\beta:R\to k$ to a field of
characteristic $p$ the kernels of $\beta\circ t_i$ and $\beta\circ t_j$ are distinct.
\end{definition}

\begin{definition}[Objects and models]\label{pf:ks1:def-obj}
Let $b$ be a base for $D$ over $R$, $N\ge1$, and $s:S\to\ksoneSp R$ a scheme over $R$. An \emph{object of level
$N$} over $S$ consists of:
\begin{enumerate}[label=(O\arabic*)]
\item an abelian scheme $A\to S$ (a commutative group scheme over $S$ which is smooth and proper with
  geometrically connected fibres), with an action $\iota:O\to\End_S(A)$;
\item a line bundle $P$ on $A\times_SA$, rigidified along $0\times A$ and $A\times0$, with an isomorphism
  $P\cong\mathrm{sw}^*P$ ($\mathrm{sw}$ the swap), whose restriction to the diagonal is relatively ample; such
  that for some integer $e$ prime to $p$ the kernel of the associated polarisation is killed by $e$ (if
  $x\in A(T)$ and the restriction of $P$ to $\{x\}\times A_T$ is trivial, then $ex=0$); and such that
  $(\iota(c)\times1)^*P\cong(1\times\iota(c^*))^*P$ for $c\in O$ (the Rosati involution induces $*$);
\item (Kottwitz's condition) for every ring map $\beta:R\to R'$ and $x:\ksoneSp R'\to S$ over $\beta$, every
  tangent vector of $A_x$ on which $\omega$ acts by $\beta(\rho)$ is a sum $\sum_{i<a}w_i$ of tangent vectors on
  which $\OO_F$ acts through $\beta\circ t_i$; and for every geometric point $x:\ksoneSp k\to S$ over
  $\beta:R\to k$ and every $i<a$, the space of tangent vectors of $A_x$ on which $\OO_F$ acts through
  $\beta\circ t_i$ has dimension $m$;
\item an additive $O$-linear map $\eta:\Lambda\to A(S)$ with $\eta(N\Lambda)=0$ such that at every geometric
  point $x$ there are a compatible system of primitive roots of unity $\zeta_r\in k$ and maps
  $\mathrm{ev}_r:\Lambda\to A_x(k)$ ($r$ prime to $p$) which are additive, $O$-linear, with kernel $r\Lambda$
  and image $A_x[r](k)$, compatible in towers, with Weil pairing $e_r(\mathrm{ev}_rx,\mathrm{ev}_ry)=
  \zeta_r^{2E_\Lambda(x,y)}$ for the polarisation of the restriction of $P$ to the diagonal, and with
  $\eta_x=\mathrm{ev}_N$.
\end{enumerate}
A morphism of objects over $f:S'\to S$ is a cartesian square of group schemes $A'\to A$ over $f$
compatible with the $O$-actions, with $P$ (up to isomorphism of line bundles) and with $\eta$. A
\emph{representing family} of level $N$ is a scheme $X$ with $s:X\to\ksoneSp R$ and an object $u$ over $X$
such that for every object $o$ over every $S\to\ksoneSp R$ there is a unique $f:S\to X$ over $R$ admitting a
morphism $o\to u$ over $f$. A \emph{model} of level $N$ is a representing family such that $s$ is
quasi-projective and smooth of relative dimension $a(n-1)$.
\end{definition}

Here quasi-projective means: $s$ factors as an immersion followed by the projection of some
$\mathbb P^M_R$; smooth of relative dimension $d$ means smooth with $\Omega_{X/R}$ locally free of rank $d$.

% ---------------------------------------------------------------------------------------------
\subsection{Existence of quasi-projective representing families}\label{pf:ks1:sec-rep}

\begin{theorem}[Integral models of Harris--Taylor type]\label{pf:ks1:thm-models}
Let $p$ be a prime, $D$ a Harris--Taylor datum for $p$, $R$ a Noetherian ring, $b$ a base for $D$ over $R$
and $N\ge3$ an integer prime to $p$. Then a model of level $N$ for $b$ exists.
\formallabel{GlobalTheory.UPel.HT.model_nonempty}
\end{theorem}

The theorem is the combination of \cref{pf:ks1:prop-repfam} (a quasi-projective representing family
exists) and \cref{pf:ks1:thm-smooth} (every quasi-projective representing family is smooth of relative
dimension $a(n-1)$); its proof is given at the end of \cref{pf:ks1:sec-smooth}. For $a=1$, $O$ a maximal
order of a division algebra and $R$ a localisation of the ring of integers of the reflex field, this is
the theorem of Kottwitz \cite[\S5]{Kottwitz1992} used in \cite[Ch.~III]{HarrisTaylor2001}; see also
\cite[Ch.~1--2]{Lan2013}. The base ring is here an arbitrary Noetherian ring with base data.

The auxiliary objects for the Siegel moduli problem are as follows. For $M\ge1$ let
$\Z[\zeta_M,1/M]\subset\C$ be the subring generated by $\zeta_M=e^{2\pi i/M}$ and $1/M$. For $h\ge1$ and a
scheme $S$ over $\Z[\zeta_M,1/M]$, a \emph{pre-Siegel object} of genus $h$ and level $M$ over $S$ is an abelian
scheme $A/S$ with a rigidified symmetric line bundle $P$ on $A\times_SA$ as in (O2), together with a map
$\mathrm{lev}:(\Z/M)^{2h}\to A[M](S)$ which is additive and, at every geometric point, an isomorphism onto
$A[M]$ under which the Weil pairing of the polarisation of $P|_\Delta$ becomes $\zeta_M^{2\omega_{\mathrm{std}}}$,
$\omega_{\mathrm{std}}$ the standard symplectic form. Its polarisation kernel is \emph{killed by $d$} if
$P|_{\Delta}$ is relatively ample and the kernel of the polarisation is killed by $d$. Morphisms are cartesian
squares compatible with $P$ and $\mathrm{lev}$.

\begin{theorem}[Siegel moduli schemes]\label{pf:ks1:thm-siegel}
Let $h\ge1$, $M\ge3$ odd and $d>0$ with $\gcd(M,d)=1$. The functor which assigns to a scheme $S$ over
$\Z[\zeta_M,1/M]$ on which $d$ is invertible the isomorphism classes of pre-Siegel objects of genus $h$ and
level $M$ whose polarisation kernel is killed by $d$ is represented by a quasi-projective
$\Z[\zeta_M,1/M]$-scheme on which $d$ is invertible.
\formallabel{GlobalTheory.T.siegelKernelBoundedRepresentable_of_coprime}
\end{theorem}

\begin{proof}
This is Mumford's theorem \cite[Ch.~7]{MumfordFogartyKirwan1994} for principally polarised
abelian schemes with level structure, combined with Zarhin's trick \cite{Zarhin1977}
($(A\times A^\vee)^4$ carries a principal polarisation) to pass to polarisations whose kernel is killed by $d$.
\end{proof}

\begin{remark}\label{pf:ks1:rem-siegel}
The formal proof of \cref{pf:ks1:thm-siegel} (about $37\,000$ lines) does not use geometric invariant
theory. It starts from a smooth family of principally polarised abelian schemes with level structure over
$\Z[\zeta_M,1/M]$, provided by a literature-labelled statement on integral toroidal compactifications of
Siegel moduli \cite{FaltingsChai1990}; it proves that this family is universal by deformation theory
(classification over algebraically closed fields, lifting to Artinian, complete Noetherian local, Noetherian
local and Noetherian bases, and gluing); and it treats non-principal polarisations by Zarhin's trick and an
open-locus argument for images of idempotents.
\end{remark}

\begin{proposition}[Relative representability over Siegel spaces]\label{pf:ks1:prop-relrep}
Let $b$ be a base for $D$ over $R$, $N\ge3$ prime to $p$, $e$ prime to $p$, $h,M\ge1$, $B$ a ring with
$\zeta\in B$, and $S$ a locally Noetherian scheme with morphisms $s_1:S\to\ksoneSp B$ and $t:S\to\ksoneSp R$.
Let $c=(A_c,P_c,\mathrm{lev}_c)$ be a pre-Siegel object over $S$ whose polarisation kernel is killed by $e$.
Then there are a finite morphism $\pi:Z\to S$ and an object $u$ of level $N$ over $Z$ (for $\pi\circ t$)
together with a cartesian square $A_u\to A_c$ over $\pi$ identifying $P_u$ with the pull-back of $P_c$, such
that: for every $g:T\to S$ and every object $o$ of level $N$ over $T$ (for $g\circ t$) together with a
cartesian square $A_o\to A_c$ over $g$ identifying $P_o$ with the pull-back of $P_c$, there is a unique
$z:T\to Z$ with $\pi z=g$ admitting a morphism $o\to u$ over $z$ compatible with the given squares.
\formallabel{GlobalTheory.T.harrisTaylor_structures_relatively_representable}
\end{proposition}

\begin{proof}
The proof has three steps.

(1) The functor of pairs $(\iota,\eta)$, where $\iota$ is an action of $O$ on the pull-back of $A_c$ satisfying the
Rosati condition for the pull-back of $P_c$ (the Rosati involution induces $*$), and $\eta:\Lambda\to A[N]$ is additive,
$O$-linear and kills $N\Lambda$, is represented by a finite $S$-scheme $E$, which is locally Noetherian
(\cref{pf:ks1:prop-actlev}). The input for the $O$-actions is the classical fact that actions of an order compatible
with a polarisation are parametrised by a finite unramified scheme (\cref{pf:ks1:thm-actrep}; cf.\ \cite[\S5]{Kottwitz1992}).

(2) Inside $E$, the locus where Kottwitz's condition (O3) holds at all points and the Tate condition (O4) holds
at all geometric points is closed: there is a closed immersion $j:Z\to E$ such that a morphism $T\to E$ factors
through $Z$ if and only if the pulled-back data satisfy (O3) and (O4) pointwise
(\cref{pf:ks1:prop-lieclosed}; closedness of Kottwitz's condition is classical \cite[\S5]{Kottwitz1992}).

(3) Over $Z$, the universal data form an object $u$ of level $N$ whose abelian scheme is the pull-back of
$A_c$ along $\pi:=$ (the composite $Z\to E\to S$), which is finite
(\lref{GlobalTheory.T.exists_harrisTaylorObject_over_lieTate_locus}). Given $g$ and $o$ as in the
statement, the $O$-action and the level map of $o$ transported to the pull-back of $A_c$ define, by (1), a
unique $z_0:T\to E$ over $g$ with an identification of the data; by the pointwise characterisation in (2)
(conditions (O3), (O4) for $o$ are equivalent to the same conditions for the pulled-back data,
\lref{GlobalTheory.T.harrisTaylorLie_tate_conditions_iff_pullbacks}), $z_0$ factors uniquely as $j\circ z$; and
the identifications assemble to a morphism $o\to u$ over $z$
(\lref{GlobalTheory.T.exists_harrisTaylorObjectMap_of_common_pullback}). Uniqueness of $z$ follows from the
uniqueness in (1) and the fact that $j$ is a monomorphism.
\end{proof}

\begin{proposition}[Auxiliary symplectic level]\label{pf:ks1:prop-aux}
Let $b$ be a base for $D$ over a Noetherian ring $R$, let $N\ge3$ and $M\ge3$ be prime to $p$ with $M$ odd,
let $e$ be prime to $p$ with $\gcd(M,e)=1$ and $e\Lambda^\vee\subseteq E_\Lambda(\Lambda)$, and let
$h=\rank_\Z\Lambda/2$. Assume the conclusion of \cref{pf:ks1:thm-siegel} for $(h,M,e)$. Then there are a
quasi-projective $R$-scheme $X$ with a morphism $s_1:X\to\ksoneSp\Z[\zeta_M,1/M]$, an object $u$ of level $N$
over $X$ and a symplectic level-$M$ structure $\mathrm{lev}$ on $A_u$ (with respect to $\zeta_M$ via $s_1$)
which are universal: for every $S$ with $t:S\to\ksoneSp R$, $t_1:S\to\ksoneSp\Z[\zeta_M,1/M]$, every object
$o$ of level $N$ over $S$ and every symplectic level-$M$ structure $l$ on $A_o$, there is a unique
$f:S\to X$ compatible with $t,t_1$ admitting a morphism $o\to u$ over $f$ carrying $l$ to $\mathrm{lev}$.
\formallabel{GlobalTheory.T.harrisTaylor_auxiliary_symplectic_level_representable}
\end{proposition}

\begin{proof}
Let $(X_0,c)$ represent the Siegel functor of \cref{pf:ks1:thm-siegel}. Since $M$ is prime to $p$ it is a
unit in $R$, so $X_0':=X_0\times_{\ksoneSp\Z}\ksoneSp R$ is quasi-projective over $R$
(\lref{GlobalTheory.T.isQuasiProjective_of_finite_over_cyclotomicAway_pullback}) and in particular locally
Noetherian. Apply \cref{pf:ks1:prop-relrep} to the pull-back of $c$ to $X_0'$: we obtain a finite
$\pi:X\to X_0'$ and an object $u$ over $X$; the level structure $\mathrm{lev}$ is the pull-back of
$\mathrm{lev}_c$. As $X$ is finite over the quasi-projective $R$-scheme $X_0'$, it is quasi-projective over $R$
(\lref{GlobalTheory.T.isQuasiProjective_of_finite_over_cyclotomicAway_pullback}). Given $(o,l)$ over $S$, the
pre-Siegel object $(A_o,P_o,l)$ has polarisation kernel killed by $e$
(\lref{GlobalTheory.T.harrisTaylor_siegelObject_polarisationKernel_killed}: this uses
$e\Lambda^\vee\subseteq E_\Lambda(\Lambda)$ and the Tate condition (O4)), and $e$ is a unit on $S$ since it is
prime to $p$. Hence it is classified by a unique $f_0:S\to X_0$, and the pair $(f_0,t)$ is a morphism
$S\to X_0'$ along which $(A_o,P_o,l)$ is the pull-back of $c$. By \cref{pf:ks1:prop-relrep} there is a
unique $f:S\to X$ over $(f_0,t)$ with a morphism $o\to u$ compatible with the identifications; compatibility
with the level structures is automatic since $\mathrm{lev}$ is pulled back from $c$. Conversely every $f$ with
the stated properties lies over $(f_0,t)$ by the uniqueness of $f_0$, so $f$ is unique.
\end{proof}

\begin{proposition}[Descent of the auxiliary level]\label{pf:ks1:prop-descent}
Under the hypotheses of \cref{pf:ks1:prop-aux}, if the conclusion of \cref{pf:ks1:prop-aux} holds, then
there is a representing family of level $N$ for $b$ whose structure morphism is quasi-projective.
\formallabel{GlobalTheory.T.exists_representingFamily_of_auxiliaryRepresentation}
\end{proposition}

\begin{proof}
Let $(X,s,s_1,u,\mathrm{lev})$ be as in \cref{pf:ks1:prop-aux} and let $\Gamma_M$ be the finite group of
symplectic similitudes of $((\Z/M)^{2h},\omega_{\mathrm{std}})$. By the universal property, $\Gamma_M$ acts on
$X$ over $R$ by changing the level structure: there are morphisms $\varrho(a):X\to X$, $a\in\Gamma_M$, with
$\varrho(1)=\mathrm{id}$, $\varrho(ac)=\varrho(c)\circ\varrho(a)$, $s\varrho(a)=s$, and morphisms of objects $u\to u$ over
$\varrho(a)$ (\lref{GlobalTheory.T.exists_symplecticSimilitude_action_on_levelFamily}; the similitude
factor of $a$ acts on $\zeta_M$ through the Galois group of $\Q(\zeta_M)/\Q$, which is why $s_1$ is not
preserved). The following three properties hold.
\begin{enumerate}[label=(\alph*)]
\item The action is free on geometric points: if $\varrho(a)\circ y=y$ for a geometric point $y$ of $X$, then $a=1$
  (\lref{GlobalTheory.T.harrisTaylor_similitude_fixing_point_eq_one}; the morphism $u\to u$ over
  $\varrho(a)$ restricts at $y$ to an endomorphism of the fibre which, by the uniqueness of morphisms of objects of
  level $N\ge3$, is the identity; since it carries the level-$M$ structure twisted by $a$ to the original one, $a$
  fixes every element of $(\Z/M)^{2h}$).
\item Every object $o$ over $S$ acquires a symplectic level-$M$ structure on a Galois cover $S_1\to S$ with group
  $\Gamma_M$ (in the sense of \cref{pf:ks1:sec-repsteps}), and the resulting classifying map $S_1\to X$ is
  $\Gamma_M$-equivariant (\cref{pf:ks1:prop-galcover}).
\item Two classifying maps $x,x':T\to X$ of the same object (with different level-$M$ structures) differ
  locally on $T$ by an element of $\Gamma_M$: every point of $T$ has an open neighbourhood $V$ and $a\in\Gamma_M$
  with $\varrho(a)\circ x|_V=x'|_V$ (\lref{GlobalTheory.T.exists_local_symplecticSimilitude_classifyingMaps}).
\end{enumerate}
A finite group acting freely on a quasi-projective $R$-scheme, together with (b) and (c), yields a quotient
$X\to X/\Gamma_M$ which is a quasi-projective representing family of level $N$: the quotient exists and is
quasi-projective, $X\to X/\Gamma_M$ is a finite \'etale Galois cover, the object $u$ descends to $X/\Gamma_M$ along
it, and (b), (c) show that $X/\Gamma_M$ represents objects of level $N$
(\cref{pf:ks1:prop-quot}).
\end{proof}

\begin{proposition}\label{pf:ks1:prop-repfam}
Let $D$, $R$, $b$, $N$ be as in \cref{pf:ks1:thm-models}. There is a representing family of level $N$ for
$b$ whose structure morphism is quasi-projective.
\formallabel{GlobalTheory.T.exists_quasiProjective_harrisTaylor_family}
\formallabel{GlobalTheory.UPel.HT.exists_repFam_qproj}
\end{proposition}

\begin{proof}
Let $e$ be as in (D4). Then $e>0$, since $\gcd(0,p)=p\neq1$. Choose a prime $M\ge p+e+3$; then $M\ge3$, $M$ is
odd, and $M$ is prime to $p$ and to $e$. Put $h=\rank_\Z\Lambda/2$. We have $\rank_\Z\Lambda\ge2$: it is positive
by (D4), and if it were $1$ the alternating form $E_\Lambda$ would vanish on a basis vector, hence vanish, contradicting
nondegeneracy; so $h\ge1$. Now
\cref{pf:ks1:thm-siegel} applies to $(h,M,e)$, \cref{pf:ks1:prop-aux} then gives a quasi-projective
representing family with auxiliary level $M$, and \cref{pf:ks1:prop-descent} removes the auxiliary level.
\end{proof}


% ---------------------------------------------------------------------------------------------
\subsection{The representability steps}\label{pf:ks1:sec-repsteps}

This subsection proves the results used in the proofs of \cref{pf:ks1:prop-relrep,pf:ks1:prop-descent}. We use the following
terminology. Let $A\to S$ be an abelian scheme, $O$ a ring with an additive map $b\mapsto\bar b$ and $Q$ a line bundle on $A\times_SA$. An
action $\iota:O\to\End_S(A)$ (a ring homomorphism) is \emph{adjoint for $Q$} if $(\iota(b)\times1)^*Q\cong(1\times\iota(\bar b))^*Q$ for every
$b\in O$; for the ring $O$ of a datum, $\bar b=b^*$ and the bundle $P$ of (O2) this is the \emph{Rosati condition}. For a finite group $\Gamma$, a
\emph{Galois cover with group $\Gamma$} is a finite, \'etale and surjective morphism $\pi_1:S_1\to S$ with morphisms $\varrho_1(a):S_1\to S_1$
($a\in\Gamma$) such that $\varrho_1(1)=\mathrm{id}$, $\varrho_1(ac)=\varrho_1(c)\circ\varrho_1(a)$ and $\pi_1\circ\varrho_1(a)=\pi_1$, and such that
for any two geometric points $y,y'$ of $S_1$ with $\pi_1y=\pi_1y'$ there is a unique $a\in\Gamma$ with $\varrho_1(a)\circ y=y'$.

\begin{theorem}[Adjoint actions form a finite unramified scheme]\label{pf:ks1:thm-actrep}
Let $S$ be a locally Noetherian scheme, $A\to S$ an abelian scheme and $Q$ a line bundle on $A\times_SA$ which is rigidified along
$0\times A$ and $A\times0$, satisfies $Q\cong\mathrm{sw}^*Q$, and whose restriction $\Delta^*Q$ to the diagonal is relatively ample. Let $O$
be a ring with an additive map $b\mapsto\bar b$, let $e_1,\dots,e_r$ generate $O$ as an abelian group, and let $W\in M_r(\Z)$ and $d>0$ be
such that $W+W^{\mathsf T}$ is positive definite and $\sum_{i,j}W_{ij}\,\bar e_ie_j=d$ in $O$. Then there are a finite and formally unramified
morphism $\pi_H:H\to S$ and an action $\iota_H$ of $O$ on $A_H=A\times_SH$ adjoint for the pull-back of $Q$ which are universal: for every
$g:T\to S$ and every action $\iota$ of $O$ on $A_T$ adjoint for the pull-back of $Q$ there is a unique $h:T\to H$ with $\pi_Hh=g$ and
$\iota=h^*\iota_H$.
\formallabel{GlobalTheory.T.adjointActions_order_finite_unramified_scheme}
\formallabel{GlobalTheory.EndLoc.actRep_finite_unramified}
\formallabel{GlobalTheory.ActFin.exists_actRep_finite}
\formallabel{GlobalTheory.EndLoc.formallyUnramified_actRep}
\end{theorem}

In the formal statement the actions are right actions, $\iota(ab)=\iota(b)\circ\iota(a)$, and the relation reads $\sum_{i,j}W_{ij}e_j\bar e_i=d$;
this is the statement above for the opposite ring, and it is applied in this form in \cref{pf:ks1:prop-actlev}.

\begin{proof}
Put $L=(\Delta^*Q)^{\otimes d}$ on $A$ and, on the abelian scheme $A^r=A\times_S\dots\times_SA$,
$M=\bigotimes_{i,j}\bigl((\mathrm{pr}_i,\mathrm{pr}_j)^*Q\bigr)^{\otimes W_{ij}}$. Both line bundles are rigidified and relatively ample; for
$M$ this uses that $W+W^{\mathsf T}$ is positive definite (\lref{GlobalTheory.ActFin.formLB_props}: $M$ is a rigidified, relatively ample
line bundle on $A^r$). Since $Q$ is rigidified along both axes, the theorem of the cube makes it biadditive: for homomorphisms
$u,u',w:A_T\to A_T$ we have $(u+u',w)^*Q\cong(u,w)^*Q\otimes(u',w)^*Q$, and similarly in the second variable.

(a) The functor of homomorphisms $\Phi:A_T\to A^r_T$ with $\Phi^*M\cong L$ is represented by a finite $S$-scheme $R_0$ with a universal
homomorphism $\Phi_0$ (\lref{GlobalTheory.T.exists_finite_unramified_representing_polarizedHoms}: for abelian schemes $A$, $B$ over a
locally Noetherian base with rigidified relatively ample line bundles $L$, $M$, the homomorphisms $\Phi$ with $\Phi^*M\cong L$ are represented
by a finite and formally unramified scheme). Let $f_k=\mathrm{pr}_k\circ\Phi_0\in\End(A_{R_0})$.

(b) For $c\in\Z^r$ and endomorphisms $g_1,\dots,g_r$ of a commutative group scheme write $g(c)=\sum_kc_kg_k$. Choose vectors
$N_1,\dots,N_s\in\Z^r$ generating the kernel of $\Z^r\to O$, $c\mapsto\sum_kc_ke_k$ (possible since $\Z^r$ is Noetherian), $c^{(0)}\in\Z^r$ with
$\sum_kc^{(0)}_ke_k=1$, and vectors $c^{(ij)}$, $\bar c^{(i)}$ with $e_ie_j=\sum_kc^{(ij)}_ke_k$ and $\bar e_i=\sum_k\bar c^{(i)}_ke_k$. Then there is an
action $\iota$ with $\iota(e_k)=g_k$ for all $k$ if and only if $g(N_l)=0$ for all $l$, $g(c^{(0)})=1$ and $g_ig_j=g(c^{(ij)})$ for all $i,j$: the first condition says that $\sum_kc_ke_k\mapsto g(c)$ is a
well-defined additive map, and the other two that it is unital and multiplicative on generators, hence everywhere, composition being
biadditive.

(c) Each condition of (b) for $g_k=f_k$ is an equality of two homomorphisms of abelian schemes over $R_0$, and the locus where two
homomorphisms agree is a closed subscheme (\lref{GlobalTheory.T.abelianScheme_homAgreeLocus_closedSubscheme}: for homomorphisms $u,w$
of abelian schemes over a locally Noetherian base there is a closed subscheme through which $T\to S$ factors if and only if $u_T=w_T$). For
each $i$ the condition $(f_i\times1)^*Q\cong(1\times f(\bar c^{(i)}))^*Q$ also defines a closed subscheme
(\lref{GlobalTheory.T.exists_closedSubscheme_pullback_adjoint_endomorphism_pair}: for endomorphisms $u,u'$ the locus where
$(u\times1)^*Q\cong(1\times u')^*Q$ is closed). Let $H\subseteq R_0$ be the intersection of these finitely many closed subschemes. Over $H$, (b)
gives an action $\iota_H$ with $\iota_H(e_k)=f_k|_H$. It is adjoint for $Q$: for the generators $e_i$ this holds by construction, as
$\iota_H(\bar e_i)=f(\bar c^{(i)})|_H$, and the set of $b$ for which it holds is a subgroup, because $b\mapsto\bar b$ is additive and $Q$ is
biadditive.

(d) Let $\iota$ be an action on $A_T$ adjoint for $Q$ and put $\Phi=(\iota(e_1),\dots,\iota(e_r)):A_T\to A_T^r$. Then $\Phi^*M\cong L$
(\lref{GlobalTheory.T.formBundle_pullback_adjointAction_eq_diagonalPower}). Indeed,
$\Phi^*M=\bigotimes_{i,j}((\iota(e_i),\iota(e_j))^*Q)^{\otimes W_{ij}}$; pulling back the adjointness isomorphism for $e_i$ along $(1,\iota(e_j))$
gives $(\iota(e_i),\iota(e_j))^*Q\cong(1,\iota(\bar e_ie_j))^*Q$; and by biadditivity the tensor product is
$(1,\iota(\sum_{i,j}W_{ij}\bar e_ie_j))^*Q=(1,d)^*Q\cong(\Delta^*Q)^{\otimes d}$. Hence $\Phi$ is classified by a unique $h_0:T\to R_0$ over $S$. As
$\iota$ satisfies the relations of (b) and is adjoint for $Q$, $h_0$ factors through $H$, and $\iota=h^*\iota_H$ for the factorisation $h$.
Conversely, if $\iota=h'^*\iota_H$, then $\Phi$ is the pull-back of $\Phi_0$ along $T\to H\to R_0$, so this composite is $h_0$, and $h'=h$
because $H\to R_0$ is a monomorphism.

(e) $\pi_H$ is finite, as the composite of a closed immersion and the finite morphism $R_0\to S$. It is formally unramified; this holds for
every universal pair $(H,\iota_H)$ as in the statement. Let $B$ be a ring, $J\subseteq B$ an ideal with $J^2=0$, and $h_1,h_2:\ksoneSp B\to H$
morphisms with $\pi_Hh_1=\pi_Hh_2$ which agree on $\ksoneSp B/J$. The actions $h_1^*\iota_H$ and $h_2^*\iota_H$ on the same abelian scheme over
$B$ agree over $B/J$; since homomorphisms of abelian schemes which agree modulo a nilpotent ideal are equal
(\lref{GlobalTheory.NilRigidity.eq_of_pullback_quotient_eq}, rigidity), they are equal, and the universal property gives $h_1=h_2$.
\end{proof}

\begin{proposition}[Actions and level maps over Siegel pull-backs]\label{pf:ks1:prop-actlev}
Let $b$, $N$, $e$, $h$, $M$, $B$, $\zeta$, $S$, $s_1$, $t$ and $c=(A_c,P_c,\mathrm{lev}_c)$ be as in \cref{pf:ks1:prop-relrep}, except that $N\ge3$ is
not required. There are a finite morphism $\pi_E:E\to S$ with $E$ locally Noetherian, an action $\iota_E$ of $O$ on $A_E=A_c\times_SE$ satisfying
the Rosati condition for the pull-back of $P_c$, and an additive $O$-linear map $\eta_E:\Lambda\to A_E(E)$ with $\eta_E(N\Lambda)=0$ which are
universal: for every $g:T\to S$, every action $\iota$ of $O$ on $A_T=A_c\times_ST$ satisfying the Rosati condition and every additive $O$-linear
$\eta:\Lambda\to A_T(T)$ with $\eta(N\Lambda)=0$ there is a unique $z:T\to E$ with $\pi_Ez=g$ and $(\iota,\eta)=z^*(\iota_E,\eta_E)$.
\formallabel{GlobalTheory.T.exists_rep_adjointAction_level_siegelPullback}
\formallabel{GlobalTheory.UPel.HT.exists_finite_rep_rosati_level}
\formallabel{GlobalTheory.UPel.HT.exists_finite_rep_oActN_rosati}
\end{proposition}

\begin{proof}
(a) There are generators $e_1,\dots,e_r$ of the abelian group $O$, a matrix $W\in M_r(\Z)$ with $W+W^{\mathsf T}$ positive definite and an
integer $d>0$ with $\sum_{i,j}W_{ij}e_i^*e_j=d$ (\lref{GlobalTheory.T.exists_spanning_posDef_relation_harrisTaylorOrder}: such data exist for
the ring $O$ of every Harris--Taylor datum).

(b) The rigidifications and the symmetry of $P_c$ and the relative ampleness of its restriction to the diagonal are part of the data of a
pre-Siegel object, so \cref{pf:ks1:thm-actrep}, applied to $A_c$, $Q=P_c$, the ring $O$ with $\bar b=b^*$ and the data of (a), gives a finite
$\pi_H:H\to S$ with an action $\iota_H$ on $A_H$ satisfying the Rosati condition which is universal among such actions. Being finite over $S$,
$H$ is locally Noetherian.

(c) Every integer prime to $p$ is a unit in $R$, hence on $H$; in particular $N$ is. For an abelian scheme with an action of $O$ over a base on
which $N$ is invertible, the functor of additive $O$-linear maps $\eta:\Lambda\to A(T)$ with $\eta(N\Lambda)=0$ is represented by a finite scheme
over the base (\lref{GlobalTheory.T.exists_finite_representing_linearLevels_abelianScheme}; such maps take values in the finite \'etale group
scheme $A[N]$, additive maps $\Lambda\to A[N]$ are parametrised by a finite \'etale scheme, and $O$-linearity, which it suffices to check on
products of basis elements of $O$ and $\Lambda$, cuts out a closed subscheme). Applied to $(A_H,\iota_H)$ this gives a finite $E\to H$ with a
universal $\eta_E$.

(d) The composite $\pi_E:E\to H\to S$ is finite, $E$ is locally Noetherian, and $\iota_E=\iota_H|_E$ satisfies the Rosati condition, which is
stable under pull-back. Given $(g,\iota,\eta)$, (b) gives a unique $z_1:T\to H$ over $S$ with $\iota=z_1^*\iota_H$, and (c), applied to $\eta$,
which is $O$-linear for $z_1^*\iota_H$, gives a unique $z:T\to E$ over $z_1$ with $\eta=z^*\eta_E$. If $z'$ is another morphism with the stated
properties, its composite with $E\to H$ has the property of (b), hence equals $z_1$, and $z'=z$ by the uniqueness in (c).
\end{proof}

\begin{lemma}[Kottwitz's condition is closed]\label{pf:ks1:lem-lieher}
Let $b$ be a base for $D$ over $R$, $S$ a scheme with $s:S\to\ksoneSp R$, $A\to S$ a smooth commutative group scheme and $\iota$ an action of
$O$ on $A$. For a ring $R'$ and $x:\ksoneSp R'\to S$ lying over $\beta:R\to R'$, say that $x$ satisfies $\mathcal K$ if the splitting and
dimension conditions of (O3) hold for $(A,\iota)$ at $x$ and at all its base changes $\ksoneSp R''\to\ksoneSp R'\to S$. Then there is a closed
immersion $i:Z_1\to S$ such that $x$ factors through $i$ if and only if $x$ satisfies $\mathcal K$.
\formallabel{GlobalTheory.UPel.HT.lieHer_closed}
\end{lemma}

\begin{proof}
The property $\mathcal K$ is stable under base change by definition. It is local on $\ksoneSp R'$: if $r_1,\dots,r_l$ generate the unit ideal of
$R'$ and $\mathcal K$ holds after each localisation $R'\to R'_{r_k}$, then it holds for $x$. For the splitting condition this follows from
\lref{GlobalTheory.T.harrisTaylor_lie_condition_of_basicOpen_cover} (the splitting condition can be checked on a cover by basic open
subsets), applied to all base changes $R'\to R''$ and the images of the $r_k$; the dimension condition concerns geometric points
$R'\to k$, and each of them factors through some $R'_{r_k}$, since some $r_k$ has nonzero image in $k$. Finally, $\mathcal K$ is locally cut out
by an ideal: every point of $S$ has an affine open neighbourhood $U$ with an ideal $J\subseteq\Gamma(U,\OO_S)$ such that the point
$\ksoneSp R'\to U$ given by $f:\Gamma(U,\OO_S)\to R'$ satisfies $\mathcal K$ if and only if $f(J)=0$
(\lref{GlobalTheory.T.kottwitzLieCondition_locally_cut_out_by_ideal}; on an affine open subset on which $\Lie A$ is free, (O3) is a condition
on the matrices of the action of $\OO_F$, and both the splitting condition and the dimension condition are equivalent to the vanishing of an
ideal). A property of points with these three features is represented by a closed subscheme
(\lref{GlobalTheory.T.isClosed_of_local_ideal_point_property}: the closed subschemes defined locally by the ideals $J$ represent the property
on each affine open and therefore glue).
\end{proof}

\begin{proposition}[Closedness of Kottwitz's condition and of the Tate condition]\label{pf:ks1:prop-lieclosed}
Let $b$, $N\ge3$, $e$, $h$, $M$, $B$, $\zeta$, $S$, $s_1$, $t$ and $c$ be as in \cref{pf:ks1:prop-relrep}, let $X$ be a locally Noetherian scheme with
$\pi_X:X\to S$, let $\iota_X$ be an action of $O$ on $A_X=A_c\times_SX$ satisfying the Rosati condition for the pull-back of $P_c$, and let
$\eta_X:\Lambda\to A_X(X)$ be additive and $O$-linear with $\eta_X(N\Lambda)=0$. Then there is a closed immersion $j:Z\to X$ such that a morphism
$g:T\to X$ factors through $j$ if and only if the data pulled back along $g$ satisfy Kottwitz's condition (O3) (with respect to $t\pi_Xg$) and
satisfy the Tate condition (O4) at every geometric point of $T$.
\formallabel{GlobalTheory.T.harrisTaylor_lie_tate_loci_closed_siegelPullback}
\formallabel{GlobalTheory.UPel.HT.lieGeom_locus_closed}
\end{proposition}

\begin{proof}
The abelian scheme $A_X$ is smooth over $X$, so \cref{pf:ks1:lem-lieher} applies to $(A_X,\iota_X)$ and $t\pi_X$ and gives a closed immersion
$i:Z_1\to X$ representing $\mathcal K$. The data $(A_X,\iota_X,P,\eta_X)$, with $P$ the pull-back of $P_c$, satisfy all conditions of an object of level
$N$ except possibly (O3) and (O4), and (O4) at a geometric point depends only on these data
(\lref{GlobalTheory.T.exists_harrisTaylorPreObject_of_siegelPullback}). Since $X$ is locally Noetherian and $N\ge3$ is prime to $p$, there is
an open and closed subset $U\subseteq X$ such that a geometric point of $X$ satisfies (O4) if and only if its image lies in $U$
(\lref{GlobalTheory.UPel.HT.PreEta.exists_clopen_geomAt}). Let $Z=Z_1\cap U$, a closed subscheme of $X$. A morphism $g:T\to X$ factors through
$Z$ if and only if it factors through $Z_1$ and has image in $U$. It factors through $Z_1$ if and only if every point $\ksoneSp R'\to T$ does (a
morphism factors through a closed subscheme if and only if all its affine points do), i.e.\ if and only if the pulled-back data satisfy (O3). It
has image in $U$ if and only if (O4) holds at every geometric point of $T$, since every point of $T$ is the image of a geometric point.
\end{proof}

\begin{proposition}[Level structures on Galois covers]\label{pf:ks1:prop-galcover}
Let $b$, $N$, $M$, $e$ and $h$ be as in \cref{pf:ks1:prop-aux}, and let $(X,s,s_1,u,\mathrm{lev})$ satisfy the conclusion of \cref{pf:ks1:prop-aux}.
For $a\in\Gamma_M$ let $\nu(a)\in\N$ represent its multiplier and $\sigma_a$ the endomorphism of $\Z[\zeta_M,1/M]$ with
$\sigma_a(\zeta_M)=\zeta_M^{\nu(a)}$. Let $\varrho(a):X\to X$ ($a\in\Gamma_M$) be morphisms with $s\varrho(a)=s$,
$s_1\varrho(a)=\ksoneSp(\sigma_a)\circ s_1$ and a morphism $u\to u$ over $\varrho(a)$ carrying the level structure $\mathrm{lev}\circ a$ to
$\mathrm{lev}$. Then for every object $o$ of level $N$ over a scheme $S$ (for $t:S\to\ksoneSp R$) there are a Galois cover $\pi_1:S_1\to S$ with
group $\Gamma_M$, a morphism $y:S_1\to X$ with $sy=t\pi_1$ and $y\circ\varrho_1(a)=\varrho(a)\circ y$ for all $a$, and an object $o'$ over $S_1$ with
morphisms $o'\to o$ over $\pi_1$ and $o'\to u$ over $y$.
\formallabel{GlobalTheory.T.symplectic_level_over_cyclotomic}
\formallabel{GlobalTheory.HTSympCover.exists_galoisCover}
\formallabel{GlobalTheory.HTSympCover.exists_galoisCover_of_geom}
\end{proposition}

\begin{proof}
Since $h\ge1$, the multiplier of $a$ is determined by $a$; hence $\nu(1)\equiv1$ and $\nu(ac)\equiv\nu(a)\nu(c)$ modulo $M$, and since an
endomorphism of $\Z[\zeta_M,1/M]$ is determined by the image of $\zeta_M$, $\sigma_1=\mathrm{id}$ and $\sigma_a\sigma_c=\sigma_{ac}$.

(1) There are a finite \'etale $S$-scheme $Y$, a morphism $t_1:Y\to\ksoneSp\Z[\zeta_M,1/M]$ and a symplectic level-$M$ structure $\mathrm{lu}$ on
$A_o\times_SY$ with respect to $t_1$ which are universal: for every $S$-scheme $T$ with $t_1':T\to\ksoneSp\Z[\zeta_M,1/M]$, the symplectic
level-$M$ structures on $A_o\times_ST$ with respect to $t_1'$ are exactly the pull-backs of $\mathrm{lu}$ along $S$-morphisms $f:T\to Y$ with
$t_1f=t_1'$, and $f$ is unique (\lref{GlobalTheory.T.exists_finiteEtale_symplecticLevel_scheme}; $M$ is invertible on $S$, so that $A_o[M]$
is finite \'etale).

(2) If $l$ is a symplectic level structure with respect to $t_1'$, then $l\circ a$ is one with respect to $\ksoneSp(\sigma_a)\circ t_1'$, since the
Weil pairing is multiplied by the multiplier of $a$ (\lref{GlobalTheory.T.symplecticLevel_twist_by_similitude}). By (1) there is a unique
$\varrho_1(a):Y\to Y$ over $S$ with $t_1\varrho_1(a)=\ksoneSp(\sigma_a)\circ t_1$ and $\varrho_1(a)^*\mathrm{lu}=\mathrm{lu}\circ a$. The uniqueness in (1)
gives $\varrho_1(1)=\mathrm{id}$ and $\varrho_1(ac)=\varrho_1(c)\circ\varrho_1(a)$.

(3) The morphism $\pi_1:Y\to S$ is a Galois cover with group $\Gamma_M$. It is finite and \'etale by (1). It is surjective: every point of $S$ is the
image of a geometric point $x:\ksoneSp k\to S$, and the fibre $A_{o,x}$ has a symplectic level-$M$ structure
(\lref{GlobalTheory.T.exists_symplectic_level_geometric_point_harrisTaylor}: by the Tate condition (O4), $A_{o,x}[M](k)$ is isomorphic to
$\Lambda/M\Lambda$ with the pairing induced by $E_\Lambda$, which is perfect because $e\Lambda^\vee\subseteq E_\Lambda(\Lambda)$ and $\gcd(M,e)=1$;
such a module has a symplectic basis, and $k$ contains a primitive $M$-th root of unity), which by (1) gives a point of $Y$ over $x$. Two
geometric points $y,y'$ of $Y$ over the same geometric point of $S$ correspond to two symplectic level structures on the same fibre, and these
differ by a unique $a\in\Gamma_M$, compatibly with the roots of unity
(\lref{GlobalTheory.T.existsUnique_similitude_between_symplecticLevels}); by (1) and (2) this means that $\varrho_1(a)\circ y=y'$ for a unique
$a$.

(4) Let $S_1=Y$, let $o'$ be the pull-back of $o$ to $S_1$ and $\mathrm{lv}'$ the level structure on $A_{o'}$ given by $\mathrm{lu}$. By the universal
property of $(X,u,\mathrm{lev})$ there is $y:S_1\to X$ with $sy=t\pi_1$, $s_1y=t_1$ and a morphism $o'\to u$ over $y$ carrying $\mathrm{lv}'$ to
$\mathrm{lev}$. Both $y\circ\varrho_1(a)$ and $\varrho(a)\circ y$ classify the object $o'$ with the level structure $\mathrm{lv}'\circ a$ (with respect
to $\ksoneSp(\sigma_a)\circ t_1$), so they are equal by the uniqueness in \cref{pf:ks1:prop-aux}
(\lref{GlobalTheory.T.classifyingMap_symplecticLevel_equivariant}).
\end{proof}

\begin{proposition}[Quotients by free actions]\label{pf:ks1:prop-quot}
Let $R$ be a Noetherian ring, $b$ a base for $D$ over $R$, $N\ge3$ prime to $p$, $X$ a quasi-projective $R$-scheme and $u$ an object of level $N$
over $X$. Let a finite group $\Gamma$ act on $X$ over $R$ by morphisms $\varrho(a)$ with $\varrho(1)=\mathrm{id}$ and
$\varrho(ac)=\varrho(c)\circ\varrho(a)$, and assume that for every $a$ there is a morphism $u\to u$ over $\varrho(a)$, and:
\begin{enumerate}[label=(\alph*)]
\item if $y$ is a geometric point of $X$ and $\varrho(a)\circ y=y$, then $a=1$;
\item for every object $o$ of level $N$ over a scheme $S$ there are a Galois cover $\pi_1:S_1\to S$ with group $\Gamma$, a morphism $y:S_1\to X$
  over $R$ with $y\circ\varrho_1(a)=\varrho(a)\circ y$, and an object $o'$ over $S_1$ with morphisms $o'\to o$ over $\pi_1$ and $o'\to u$ over $y$;
\item if $x,x':T\to X$ are morphisms over $R$ admitting morphisms $o\to u$ from the same object $o$ over $T$, then every point of $T$ has an open
  neighbourhood $V$ with $\varrho(a)\circ x|_V=x'|_V$ for some $a\in\Gamma$.
\end{enumerate}
Then there is a representing family of level $N$ for $b$ whose structure morphism is quasi-projective.
\formallabel{GlobalTheory.HTQuot.exists_repFam_of_levelCover}
\end{proposition}

\begin{proof}
(1) Since $X$ is quasi-projective over the Noetherian ring $R$, every orbit of $\Gamma$ lies in an affine open subset, and the quotient
$\pi:X\to Y=X/\Gamma$ exists; it is obtained by gluing the spectra of the rings of invariants of $\Gamma$-stable affine open subsets, and by (a)
it is a Galois cover with group $\Gamma$ (\lref{GlobalTheory.ZM.exists_isGaloisCover_of_isQProjOver_noetherian}; cf.\
\cite[Exp.~V, \S1]{Grothendieck1971}). The scheme $Y$ is quasi-projective over $R$
(\lref{GlobalTheory.GaloisQuotientRing.isQProjOver_of_isGaloisCover}: the base of a Galois cover whose total space is quasi-projective over
$R$ is quasi-projective over $R$).

(2) The morphisms $u\to u$ over $\varrho(a)$ are descent data for $u$ along $\pi$, and $u$ descends: there are an object $u_Y$ of level $N$ over $Y$
and a morphism $u\to u_Y$ over $\pi$ (\lref{GlobalTheory.HTQuot.obj_descends_galoisCover}: the abelian scheme, the action, the line bundle and
the level map descend along a Galois cover, and so do the conditions (O2)--(O4)).

(3) Let $o$ be an object of level $N$ over $S$ (for $t:S\to\ksoneSp R$) and let $(S_1,\pi_1,y,o')$ be as in (b). The morphism $\pi y$ is invariant:
$\pi y\varrho_1(a)=\pi\varrho(a)y=\pi y$. Hence it factors uniquely as $f\pi_1$ with $f:S\to Y$, as $S$ is the quotient of $S_1$ by $\Gamma$; and
$f$ lies over $R$ since $\pi_1$ is an epimorphism. Composing $o'\to u\to u_Y$ gives a morphism $o'\to u_Y$ over $f\pi_1$, which together with
$o'\to o$ over $\pi_1$ descends to a morphism $o\to u_Y$ over $f$ (\lref{GlobalTheory.T.harrisTaylor_moduliMap_descends_galoisCover}: a
morphism from the pull-back of $o$ to a Galois cover of $S$ into $u_Y$ descends to a morphism from $o$).

Uniqueness: let $f':S\to Y$ over $R$ admit a morphism $o\to u_Y$. Let $T=S_1\times_YX$ be the fibre product of $f'\pi_1$ and $\pi$, with
projections $q_1$, $q_2$, and let $o_T$ be the pull-back of $o'$ to $T$. Both $yq_1$ and $q_2$ admit morphisms $o_T\to u$: the first through
$o'\to u$, the second because the morphism $o_T\to o\to u_Y$ lies over $f'\pi_1q_1=\pi q_2$ and $u\to u_Y$ is a cartesian square over $\pi$. By (c),
every point of $T$ has a neighbourhood on which $q_2=\varrho(a)\circ y\circ q_1$ for some $a$, hence $\pi q_2=\pi yq_1$, i.e.\ $f'\pi_1q_1=f\pi_1q_1$.
As $q_1$ is a base change of the finite \'etale surjective $\pi$, it is an epimorphism, so $f'\pi_1=f\pi_1$ and $f'=f$.

Hence $(Y,u_Y)$ is a representing family of level $N$, and its structure morphism is quasi-projective by (1).
\end{proof}

% ---------------------------------------------------------------------------------------------
\subsection{Smoothness and relative dimension}\label{pf:ks1:sec-smooth}

\begin{theorem}[Smoothness]\label{pf:ks1:thm-smooth}
Let $D$, $R$, $b$, $N$ be as in \cref{pf:ks1:thm-models} and let $(X,s,u)$ be a representing family of level
$N$ for $b$ with $s$ quasi-projective. Then $s$ is smooth of relative dimension $a(n-1)$.
\formallabel{GlobalTheory.T.smoothOfRelativeDimension_repFam_of_isQProjOver}
\formallabel{GlobalTheory.UPel.HT.RepFam.smoothOfRelativeDimension}
\end{theorem}

The proof uses four results of the formal development, each of which is a form of a classical statement.

\begin{enumerate}[label=(S\arabic*)]
\item\label{pf:ks1:S1} (Infinitesimal lifting criterion; \lref{GlobalTheory.T.smooth_of_artinian_squareZero_lifting_noetherian_base},
  cf.\ \cite[\S17]{Grothendieck1967}.) Let $B$ be a Noetherian ring and $s:X\to\ksoneSp B$ locally of finite type.
  Assume that for every local Artinian ring $C$, every surjection $\pi:C\to C_0$ with $(\ker\pi)^2=0$, every
  $t:\ksoneSp C\to\ksoneSp B$ and every $x_0:\ksoneSp C_0\to X$ with $sx_0=t\circ\ksoneSp\pi$ there is
  $x:\ksoneSp C\to X$ with $x\circ\ksoneSp\pi=x_0$ and $sx=t$. Then $s$ is smooth.
\item\label{pf:ks1:S2} (Lifting of objects; \cref{pf:ks1:thm-sqzero}.)
  Let $N\ge3$ be prime to $p$, $C$ a local Artinian ring, $\pi:C\to C_0$ a surjection with $(\ker\pi)^2=0$ and
  $\beta:R\to C$. Every object of level $N$ over $\ksoneSp C_0$ (for $\pi\beta$) is the pull-back along
  $\ksoneSp\pi$ of an object of level $N$ over $\ksoneSp C$ (for $\beta$).
\item\label{pf:ks1:S3} (Dimension from a universal tangent; \lref{GlobalTheory.T.smoothOfRelativeDimension_of_universal_firstOrder_tangent}.)
  Let $s:X\to\ksoneSp B$ be smooth and $d\ge0$. Assume that for every algebraically closed field $k$ and every
  $x:\ksoneSp k\to X$ there is $u:\ksoneSp(k\oplus k^d)\to X$ (with $k\oplus k^d$ the square-zero extension)
  restricting to $x$ and lying over $s x$, such that every $v:\ksoneSp k[\varepsilon]\to X$ restricting to $x$
  and lying over $sx$ is $u\circ\ksoneSp(\mathrm{id}+c)$ for a unique $k$-linear $c:k^d\to k$. Then $s$ is smooth
  of relative dimension $d$.
\item\label{pf:ks1:S4} (Universal first-order deformations; \cref{pf:ks1:prop-univtan}.)
  If $(X,s,u)$ is as in \cref{pf:ks1:thm-smooth} and $s$ is smooth, the hypothesis of \ref{pf:ks1:S3} holds
  with $d=a(n-1)$.
\end{enumerate}

Statements \ref{pf:ks1:S2} and \ref{pf:ks1:S4} are proved in \cref{pf:ks1:sec-deform} with Grothendieck--Messing theory
\cite{Messing1972}: a lift of an abelian scheme over a square-zero thickening is determined by a lift of its Hodge filtration, the
$O$-action and the polarisation lift if and only if the lifted filtration is $O$-stable and isotropic, and Kottwitz's condition
becomes a splitting condition modulo the lifted filtration. The linear-algebra core of \ref{pf:ks1:S2} is the existence of
$O$-stable Lagrangian lifts satisfying this splitting condition (\cref{pf:ks1:prop-laglift}). For \ref{pf:ks1:S4}, the first-order
deformations of an object at a geometric point are described by deformations of the Hodge filtration which are $O$-linear,
symmetric for the polarisation pairing (for which the Hodge filtration is Lagrangian, \cref{pf:ks1:prop-fillagr}) and compatible
with Kottwitz's condition, and these Kodaira--Spencer directions form a space of dimension $a(n-1)$ (\cref{pf:ks1:prop-ksdim}). In
characteristic $0$ with $m=1$ this is the familiar count: the $\tau_i$-part of $\Lie A$ is a line, the $\tau_i$-part of the Hodge
filtration has dimension $n-1$, symmetric deformations are determined by their $\tau_i$-components, and the parts for the other
embeddings do not contribute since there the Hodge filtration or $\Lie A$ vanishes; in general the formal proof works with the
eigenspaces of $\omega$, which are separated because $\mathrm{tr}^2-4\,\mathrm{nm}$ is a unit, and uses that the $t_i$ are pairwise
distinct modulo every prime of characteristic $p$ (\cref{pf:ks1:def-base}). These are the classical arguments of
\cite[\S5]{Kottwitz1992} and \cite[Ch.~2]{Lan2013}.

\begin{proof}[Proof of \cref{pf:ks1:thm-smooth}]
A quasi-projective morphism is the composite of an immersion and the projection of a projective space, hence
locally of finite type. We verify the hypothesis of \ref{pf:ks1:S1} with $B=R$. Let $C,\pi,t,x_0$ be given and
write $t=\ksoneSp\beta$ for $\beta:R\to C$. The pull-back of the universal object $u$ along $x_0$ is an object
$o_0$ of level $N$ over $\ksoneSp C_0$ (for $\pi\beta$)
(\lref{GlobalTheory.T.exists_harrisTaylorObjectPullback}). By \ref{pf:ks1:S2} there is an object $o'$ over
$\ksoneSp C$ (for $\beta$) with a morphism $o_0\to o'$ over $\ksoneSp\pi$. By the universal property there is
$x:\ksoneSp C\to X$ over $R$ with a morphism $o'\to u$ over $x$. Composing, $o_0\to u$ is a morphism over
$x\circ\ksoneSp\pi$; since $x_0$ is the unique morphism over $R$ admitting a morphism $o_0\to u$, we get
$x\circ\ksoneSp\pi=x_0$. Hence $s$ is smooth, and \ref{pf:ks1:S4} and \ref{pf:ks1:S3} show that it is smooth of
relative dimension $a(n-1)$.
\end{proof}

\begin{proof}[Proof of \cref{pf:ks1:thm-models}]
By \cref{pf:ks1:prop-repfam} there is a representing family $(X,s,u)$ with $s$ quasi-projective; by
\cref{pf:ks1:thm-smooth} it is smooth of relative dimension $a(n-1)$; hence it is a model.
\end{proof}

% ---------------------------------------------------------------------------------------------
\subsection{Deformation theory}\label{pf:ks1:sec-deform}

This subsection proves \ref{pf:ks1:S2} and \ref{pf:ks1:S4}. For an abelian scheme $A$ over a ring $C_0$ let $H^{\mathrm{dR}}_1(A)$ be its first de
Rham homology, the $C_0$-linear dual of the module of primitive classes in $H^1_{\mathrm{dR}}(A/C_0)$, and let $\mathrm{Fil}(A)\subseteq
H^{\mathrm{dR}}_1(A)$ be the Hodge filtration, the annihilator of the invariant differentials, so that $H^{\mathrm{dR}}_1(A)/\mathrm{Fil}(A)=\Lie A$.
Endomorphisms of $A$ act on $H^{\mathrm{dR}}_1(A)$ and preserve $\mathrm{Fil}(A)$. For abelian schemes $A$, $B$ and a line bundle $P$ on $A\times B$,
a bilinear form $\varphi$ on $H^{\mathrm{dR}}_1(A)\times H^{\mathrm{dR}}_1(B)$ is \emph{induced by $P$} if the de Rham Chern class of $P$ is
$\sum_k\mathrm{pr}_1^*a_k\cup\mathrm{pr}_2^*b_k$ with primitive classes $a_k$, $b_k$ and $\varphi(x,y)=\sum_kx(a_k)\,y(b_k)$.

Let $\pi:C\to C_0$ be a surjection of rings with $I=\ker\pi$ and $I^2=0$, and give $I$ the trivial divided powers. For an abelian scheme $A_0$
over $C_0$ let $\mathbb D(A_0)_C$ be the value of its crystal on the thickening $\ksoneSp C_0\to\ksoneSp C$. It is functorial in $A_0$, satisfies
$\mathbb D(A_0)_C\otimes_CC_0=H^{\mathrm{dR}}_1(A_0)$, and for an abelian scheme $A$ over $C$ lifting $A_0$ it is identified with $H^{\mathrm{dR}}_1(A)$;
the image $E_A$ of $\mathrm{Fil}(A)$ is a direct summand lifting $\mathrm{Fil}(A_0)$, the \emph{Hodge filtration of the lift}. By
Grothendieck--Messing theory \cite{Messing1972}, lifts of $A_0$ correspond to direct summands of $\mathbb D(A_0)_C$ lifting $\mathrm{Fil}(A_0)$, and an
endomorphism of $A_0$ lifts to $A$ if and only if it preserves $E_A$. For $C=k[\varepsilon]$ and $C_0=k$ a field, $\mathbb D(A_0)_{k[\varepsilon]}=
H^{\mathrm{dR}}_1(A_0)\otimes_kk[\varepsilon]$, and a linear map $f$ of $H^{\mathrm{dR}}_1(A_0)$ to itself is a \emph{Kodaira--Spencer class} of a lift $A$ if
$x+\varepsilon f(x)\in E_A$ for every $x\in\mathrm{Fil}(A_0)$; similarly, for the square-zero extension $k\oplus k^d$ a $d$-tuple $(f_1,\dots,f_d)$ is a
Kodaira--Spencer class if $x+\sum_j\varepsilon_jf_j(x)\in E_A$ for $x\in\mathrm{Fil}(A_0)$, $\varepsilon_j$ the basis vectors of $k^d$. If $f$ is a
Kodaira--Spencer class of a first-order lift, then the Kodaira--Spencer classes of that lift are exactly the maps $f'$ with
$f'(x)-f(x)\in\mathrm{Fil}(A_0)$ for $x\in\mathrm{Fil}(A_0)$ (\lref{GlobalTheory.T.hodgeFiltration_lift_determines_mod_fil}). A \emph{deformation}
of an object $o_0$ over $\ksoneSp k$ to $k[\varepsilon]$ (or $k\oplus k^d$) is an object $o$ over it with a morphism $o_0\to o$ over the closed point;
its Kodaira--Spencer classes are those of its abelian scheme.

\begin{proposition}[The Hodge filtration is Lagrangian]\label{pf:ks1:prop-fillagr}
Let $b$ be a base for $D$ over $R$, $C_0$ a ring with $\ksoneSp C_0\to\ksoneSp R$, and $o_0$ an object of level $N$ over $\ksoneSp C_0$ whose abelian
scheme $A_0$ has a relative dimension. Let $\varphi$ be a bilinear form on $H=H^{\mathrm{dR}}_1(A_0)$ induced by $P_{o_0}$. Then
$\varphi(\mathrm{Fil},\mathrm{Fil})=0$, and every $x\in H$ with $\varphi(x,\mathrm{Fil})=0$ lies in $\mathrm{Fil}=\mathrm{Fil}(A_0)$.
\formallabel{GlobalTheory.HTLagr.fil_lagrangian}
\end{proposition}

\begin{proof}
The first assertion holds for every line bundle $P$ on a product $A\times B$ of abelian schemes of the same relative dimension and every form
induced by $P$ (\lref{GlobalTheory.PtExt.HodgeOrth.fil_orthogonal}: such a form vanishes on $\mathrm{Fil}(A)\times\mathrm{Fil}(B)$). In its formal
proof the elements of the Hodge filtrations are represented by \v{C}ech cocycles for an open cover trivialising $P$, and the cup product of such
cocycles with the Chern cocycle of $P$ is shown to be a coboundary.

For the second assertion we may assume $C_0\neq0$. Let $e$ be an integer prime to $p$ which kills the polarisation kernel as in (O2); it is a unit
of $R$, hence of $C_0$. Put $P'=\mathrm{sw}^*P_{o_0}$. Then $P'$ is trivial on $A_0\times0$, since $P_{o_0}$ is rigidified along $0\times A_0$; the
form $\varphi'(x,y)=-\varphi(y,x)$ is induced by $P'$ (\lref{GlobalTheory.SymmPoincare.isCorrForm_swap}: interchanging the factors replaces the
induced form by its negative transpose); and every point $v$ of $A_0$ with values in $C_0[\varepsilon]$ such that $P'$ is trivial on $A_0\times v$ is
killed by $e$, by the condition on the polarisation kernel in (O2). Under these conditions every element of $H^{\mathrm{dR}}_1(A_0)$ which is
orthogonal to $\mathrm{Fil}$ for $\varphi'$ (in the first variable) lies in $\mathrm{Fil}$ (\lref{GlobalTheory.HTLagr.mem_fil_of_orth_of_ker}: for a
line bundle on $A\times B$ which is trivial on $A\times0$ and such that every first-order point $v$ of $B$ with trivial restriction to $A\times v$
is killed by a unit $e$, every $y\in H^{\mathrm{dR}}_1(B)$ with $\varphi'(\mathrm{Fil}(A),y)=0$ lies in $\mathrm{Fil}(B)$). The idea of that statement:
the image of $y$ in $\Lie B=H^{\mathrm{dR}}_1(B)/\mathrm{Fil}(B)$ is a first-order point $v$ of $B$; orthogonality to $\mathrm{Fil}(A)$ means that a
cocycle describing the restriction of the line bundle to $A\times v$ is a coboundary, so this restriction is trivial; hence $ev=0$, so $v=0$ as
$e$ is a unit, i.e.\ $y\in\mathrm{Fil}(B)$. Since $\varphi'(z,x)=-\varphi(x,z)=0$ for $z\in\mathrm{Fil}$, we obtain $x\in\mathrm{Fil}$.
\end{proof}

\begin{proposition}[Lagrangian lifts]\label{pf:ks1:prop-laglift}
Let $b$ be a base for $D$ over $R$, $C$ a local Artinian ring with $\beta:R\to C$, $I\subseteq C$ an ideal with $I^2=0$, and $H$ a finite projective
$C$-module with a ring homomorphism $O\to\End_C(H)$, $b\mapsto(x\mapsto bx)$, such that $H$ is projective over $C\otimes\OO_F$. Let $\varphi$ be a
perfect $C$-bilinear form on $H$ with $\varphi(x,y)=-\varphi(y,x)$ and $\varphi(bx,y)=\varphi(x,b^*y)$ for $b\in O$. Let $F\subseteq H$ be an $O$-stable submodule with
$IH\subseteq F$ and $\varphi(F,F)\subseteq I$, such that every $x$ with $\varphi(x,F)\subseteq I$ lies in $F$, such that $F=E'+IH$ for a direct
summand $E'$ of $H$, and such that the \emph{splitting condition modulo $F$} holds:
\begin{itemize}
\item[$(*)_F$] for every $d\in H$ with $\omega d-\beta(\rho)d\in F$ there are $w_0,\dots,w_{a-1}\in H$ with $cw_i-\beta(t_i(c))w_i\in F$ for all
  $c\in\OO_F$ and $i<a$, and $d-\sum_iw_i\in F$.
\end{itemize}
Then there is an $O$-stable direct summand $E\subseteq H$ with $E+IH=F$ and $\varphi(E,E)=0$ which satisfies $(*)_E$.
\formallabel{GlobalTheory.T.exists_lagrangian_lift_lie_squareZero_harrisTaylor}
\formallabel{GlobalTheory.HTSummand.exists_lagrangian_lift_lie}
\end{proposition}

\begin{proof}
Write $\rho$ for $\beta(\rho)$ and $\rho'=\mathrm{tr}-\rho$.

(1) Eigenspaces of $\omega$. In $C$ we have $(\rho-\rho')^2=(2\rho-\mathrm{tr})^2=\mathrm{tr}^2-4\,\mathrm{nm}$, as $\rho^2=\mathrm{tr}\,\rho-\mathrm{nm}$; this
integer is prime to $p$ by (D2), hence a unit of $R$ and of $C$. As $\omega^2-\mathrm{tr}\,\omega+\mathrm{nm}=0$, $\rho+\rho'=\mathrm{tr}$ and
$\rho\rho'=\mathrm{nm}$, the endomorphism $\omega$ of $H$ satisfies $(\omega-\rho)(\omega-\rho')=0$. Hence $e=(\rho-\rho')^{-1}(\omega-\rho')$ is an
idempotent, $H=H_\rho\oplus H_{\rho'}$ with $H_\rho=eH=\ker(\omega-\rho)$ and $H_{\rho'}=(1-e)H=\ker(\omega-\rho')$, and $\omega-\rho$ acts on
$H_{\rho'}$ by the unit $\rho'-\rho$. Since $\omega$ is central in $O$, $e$ commutes with $O$; so every $O$-stable submodule $V$ satisfies
$V=(V\cap H_\rho)\oplus(V\cap H_{\rho'})$, and $IH\cap H_\rho=IH_\rho$. Since $\omega^*=\bar\omega=\mathrm{tr}-\omega$ by (D3), for $x,y\in H_\rho$ we
get $\rho\varphi(x,y)=\varphi(\omega x,y)=\varphi(x,(\mathrm{tr}-\omega)y)=\rho'\varphi(x,y)$, so $\varphi(x,y)=0$; likewise $H_{\rho'}$ is isotropic.
As $\varphi$ is perfect, $y\mapsto\varphi(\cdot,y)|_{H_\rho}$ is an isomorphism $H_{\rho'}\to\Hom_C(H_\rho,C)$: it is injective because a
$y\in H_{\rho'}$ with $\varphi(H_\rho,y)=0$ satisfies $\varphi(H,y)=0$, $H_{\rho'}$ being isotropic; and it is surjective because every linear form
on $H_\rho$, extended by $0$ on $H_{\rho'}$, is of the form $\varphi(\cdot,y)$, and then also of the form $\varphi(\cdot,(1-e)y)$ on $H_\rho$.

(2) The $\rho$-part. Every $d\in H_\rho$ satisfies $\omega d-\rho d=0$, so it satisfies the conclusion of $(*)_F$. There is an $O$-stable submodule
$G\subseteq H_\rho$ with a complement $G_1$ in $H_\rho$ such that $G+IH_\rho=F\cap H_\rho$ and every $d\in H_\rho$ is of the form $\sum_iw_i+g$ with
$g\in G$ and $cw_i-\beta(t_i(c))w_i\in G$ for all $c$ and $i$ (\lref{GlobalTheory.T.exists_liftSummand_lie_residueCharZero} if $p$ is a unit
in $C$, and \lref{GlobalTheory.T.exists_lieCondition_summand_residueCharP} otherwise: an $O$-stable direct summand of $H_\rho$ lifting $F\cap
H_\rho$ modulo $I$ and satisfying the splitting condition exists; in residue characteristic $0$ it is constructed with separability idempotents of
$O\otimes\Q$, in residue characteristic $p$ with the isomorphism $O\otimes\Z_p\cong M_m(\OO_F\otimes\Z_p)$ of (D3) and Morita equivalence, using that
$H$ is projective over $C\otimes\OO_F$).

(3) Lagrangian completion. Let $G'=\{y\in H_{\rho'}:\varphi(G,y)=0\}$ and $E=G\oplus G'$. Then $E$ is totally isotropic, since $H_\rho$ and $H_{\rho'}$
are isotropic and $\varphi(G,G')=0=\varphi(G',G)$. It is $O$-stable: for $y\in G'$, $b\in O$ and $x\in G$ we have $by\in H_{\rho'}$ and
$\varphi(x,by)=\varphi(b^*x,y)=0$, as $b^{**}=b$ and $b^*x\in G$. It is a direct summand: under $H_{\rho'}\cong\Hom_C(H_\rho,C)$ the submodule $G'$
corresponds to the forms vanishing on $G$, so $H_{\rho'}=G'\oplus G_1'$ with $G_1'$ corresponding to the forms vanishing on $G_1$, and
$H=E\oplus(G_1\oplus G_1')$. Finally $E+IH=F$. Since $F$ and $IH$ are $O$-stable and $F\cap H_\rho=G+IH_\rho$, it suffices to show that
$F\cap H_{\rho'}=G'+IH_{\rho'}$. The two conditions on $F$ say that $F=\{x:\varphi(x,F)\subseteq I\}$. For $y\in H_{\rho'}$ we have
$\varphi(F\cap H_{\rho'},y)=0$ and $\varphi(IH_\rho,y)\subseteq I$, so $y\in F$ if and only if $\varphi(G,y)\subseteq I$. This shows $G'+IH_{\rho'}\subseteq
F$. Conversely, let $y\in F\cap H_{\rho'}$, so that $\lambda=\varphi(\cdot,y)$ maps $G$ into $I$. Extending $\lambda|_G$ by $0$ on $G_1$ gives
$\tilde\lambda\in\Hom_C(H_\rho,I)=I\cdot\Hom_C(H_\rho,C)$ (as $H_\rho$ is finite projective), so $\tilde\lambda=\varphi(\cdot,z)|_{H_\rho}$ for some
$z\in IH_{\rho'}$ by (1). Then $y-z\in G'$, and $y\in G'+IH_{\rho'}$.

(4) The splitting condition modulo $E$. Let $d\in H$ with $\omega d-\rho d\in E$, and write $d=ed+k$ with $k=(1-e)d\in H_{\rho'}$. As $E$ is stable
under $e$, $(1-e)(\omega-\rho)d=(\omega-\rho)k=(\rho'-\rho)k$ lies in $E$, hence $k\in E$. By (2), $ed=\sum_iw_i+g$ with $g\in G$ and
$cw_i-\beta(t_i(c))w_i\in G\subseteq E$. Hence $d-\sum_iw_i=g+k\in E$.
\end{proof}

\begin{lemma}[Lifting the Hodge filtration]\label{pf:ks1:lem-hodgelift}
Let $b$ be a base for $D$ over $R$, $C$ a local Artinian ring with $\beta:R\to C$, $\pi:C\to C_0$ a surjection with $I=\ker\pi$, $I^2=0$, and $o_0$
an object of level $N$ over $\ksoneSp C_0$ (for $\pi\beta$) with abelian scheme $A_0$. Let $H=\mathbb D(A_0)_C$, with the action of $O$ induced by
$\iota_{o_0}$. There is an $O$-stable direct summand $E\subseteq H$ whose image in $H\otimes_CC_0=H^{\mathrm{dR}}_1(A_0)$ is $\mathrm{Fil}(A_0)$, which
is totally isotropic for every form on $H$ induced by $P_{o_0}$, and which satisfies the splitting condition $(*)_E$ of
\cref{pf:ks1:prop-laglift}.
\formallabel{GlobalTheory.T.harrisTaylor_exists_isotropic_stable_summand}
\formallabel{GlobalTheory.HTSummand.exists_summand}
\end{lemma}

\begin{proof}
(1) Let $\mathrm{red}:H\to H^{\mathrm{dR}}_1(A_0)$ be the reduction map. The module $H$ is finite projective over $C$, $\mathrm{red}$ is surjective,
$O$-equivariant and has kernel $IH$, and there is a direct summand $E'\subseteq H$ with $\mathrm{red}(E')=\mathrm{Fil}(A_0)$
(\lref{GlobalTheory.CrysRed.exists_crysRed}). The module $H$ is projective over $C\otimes\OO_F$
(\lref{GlobalTheory.T.crystal_projective_over_tensor_ringOfIntegers}). There is a unique form $\varphi_0$ on $H$ induced by $P_{o_0}$; it
satisfies $\varphi_0(bx,y)=\varphi_0(x,b^*y)$, and $\pi(\varphi_0(x,y))=\varphi'(\mathrm{red}\,x,\mathrm{red}\,y)$ for a form $\varphi'$ on
$H^{\mathrm{dR}}_1(A_0)$ induced by $P_{o_0}$ (\lref{GlobalTheory.HTDefLift.exists_crysForm}); and $\varphi_0$ is skew-symmetric and perfect
(\lref{GlobalTheory.HTPerf.skew_bijective}).

(2) Let $F=\mathrm{red}^{-1}(\mathrm{Fil}(A_0))$. Then $IH\subseteq F$; $F$ is $O$-stable, as $\mathrm{Fil}(A_0)$ is stable under endomorphisms; and
$F=E'+IH$. Since $C_0$ is local, $A_0$ has a relative dimension, and \cref{pf:ks1:prop-fillagr} gives $\varphi_0(F,F)\subseteq I$. If
$\varphi_0(x,F)\subseteq I$, then $\varphi'(\mathrm{red}\,x,\mathrm{Fil}(A_0))=0$ because $\mathrm{red}(F)=\mathrm{Fil}(A_0)$, so
$\mathrm{red}\,x\in\mathrm{Fil}(A_0)$ by \cref{pf:ks1:prop-fillagr} and $x\in F$. Condition (O3) for $o_0$ says that the
splitting condition holds in $H^{\mathrm{dR}}_1(A_0)$ modulo $\mathrm{Fil}(A_0)$ (\lref{GlobalTheory.HTSummand.lie_mod_fil}: Kottwitz's condition
for $o_0$, transported through $\Lie A_0=H^{\mathrm{dR}}_1(A_0)/\mathrm{Fil}(A_0)$); lifting the elements $w_i$ along the surjection $\mathrm{red}$ gives
$(*)_F$.

(3) \Cref{pf:ks1:prop-laglift} gives an $O$-stable direct summand $E$ with $E+IH=F$, hence $\mathrm{red}(E)=\mathrm{Fil}(A_0)$, with
$\varphi_0(E,E)=0$ and $(*)_E$. By the uniqueness in (1), $E$ is isotropic for every form induced by $P_{o_0}$.
\end{proof}

\begin{theorem}[Square-zero lifts of objects]\label{pf:ks1:thm-sqzero}
Let $b$ be a base for $D$ over $R$, $N$ prime to $p$, $C$ a local Artinian ring, $\pi:C\to C_0$ a surjective ring homomorphism with $(\ker\pi)^2=0$
and $\beta:R\to C$. For every object $o_0$ of level $N$ over $\ksoneSp C_0$ (for $\pi\beta$) there are an object $o$ of level $N$ over $\ksoneSp C$
(for $\beta$) and a morphism $o_0\to o$ over $\ksoneSp\pi$.
\formallabel{GlobalTheory.T.harrisTaylor_formally_smooth_squareZero_lift}
\formallabel{GlobalTheory.UPel.HT.obj_lift_sqZero}
\formallabel{GlobalTheory.T.harrisTaylor_exists_lift_squareZero_artinian}
\formallabel{GlobalTheory.UPel.HT.exists_sqZero_lift}
\end{theorem}

The first two formal statements also assume $N\ge3$.

\begin{proof}
Let $E\subseteq\mathbb D(A_{o_0})_C$ be as in \cref{pf:ks1:lem-hodgelift}.

(a) By Grothendieck--Messing theory there are an abelian scheme $A$ over $C$ lifting $A_{o_0}$, with Hodge filtration $E$, and an action $\iota$ of $O$
on $A$ lifting $\iota_{o_0}$ (\lref{GlobalTheory.HTDeform.exists_lift_oAct}: for an $O$-stable direct summand lifting
the Hodge filtration there is such a lift, and the action lifts because $E$ is $O$-stable).

(b) As $E$ is isotropic, the polarisation bundle lifts: there are a line bundle $P$ on $A\times_CA$ whose pull-back to $A_{o_0}\times A_{o_0}$ is
$P_{o_0}$ and which satisfies the conditions of (O2) (rigidifications, symmetry, relative ampleness of the restriction to the diagonal, polarisation
kernel killed by an integer prime to $p$, Rosati condition) (\lref{GlobalTheory.T.harrisTaylor_exists_polarisation_bundle_lift}).

(c) The level map $\eta_{o_0}$ lifts to an additive $O$-linear map $\eta:\Lambda\to A(C)$ with $\eta(N\Lambda)=0$ which satisfies the Tate condition (O4)
for $P$ (\lref{GlobalTheory.T.harrisTaylor_primeToP_level_lift_squareZero}; $N$ is invertible in $C$).

(d) Kottwitz's condition (O3) for $(A,\iota)$: since $\Lie A=\mathbb D(A_{o_0})_C/E$ and $(*)_E$ holds, the splitting part of (O3) holds at all points of
$\ksoneSp C$ (\lref{GlobalTheory.T.harrisTaylor_lift_lie_of_hodge_summand_lie}); the dimension part concerns geometric points, which factor
through $\ksoneSp C_0$ as $(\ker\pi)^2=0$, and there it holds for $o_0$ (\lref{GlobalTheory.HTLiftLie.lieMult_of_sqZeroLift}).

Hence $o=(A,\iota,P,\eta)$ is an object of level $N$ over $\ksoneSp C$, and the closed immersion $A_{o_0}\to A$, which is compatible with the actions,
the line bundles and the level maps, is a morphism $o_0\to o$ over $\ksoneSp\pi$.
\end{proof}

\begin{definition}[Kodaira--Spencer directions]\label{pf:ks1:def-ksdir}
Let $k$ be a field, $H$ a finite-dimensional $k$-vector space with a ring homomorphism $O\to\End_k(H)$, $\varphi$ a bilinear form on $H$, and
$\mathrm{Fil},Z\subseteq H$ subspaces. A $k$-linear map $f:H\to H$ is a \emph{Kodaira--Spencer direction} for $(\varphi,\mathrm{Fil},Z)$ if
$f(bx)-bf(x)\in\mathrm{Fil}$ for all $b\in O$ and $x\in\mathrm{Fil}$, if $f(Z)\subseteq\mathrm{Fil}$, and if $\varphi(fx,y)+\varphi(x,fy)=0$ for all
$x,y\in\mathrm{Fil}$. The Kodaira--Spencer directions for $(\varphi,\mathrm{Fil},Z)$ \emph{have dimension $d$} if there are directions $f_1,\dots,f_d$
such that for every direction $g$ there is a unique $c\in k^d$ with $g(x)-\sum_ic_if_i(x)\in\mathrm{Fil}$ for all $x\in\mathrm{Fil}$.
\formallabel{GlobalTheory.UPel.Sm.KSDim}
\end{definition}

Equivalently, the image of the space of directions in $\Hom_k(\mathrm{Fil},H/\mathrm{Fil})$ has dimension $d$.

\begin{lemma}[Reduction to an eigenspace]\label{pf:ks1:lem-ksreduce}
Let $k$, $H$, $O$ be as in \cref{pf:ks1:def-ksdir}, let $\varphi$ be a nondegenerate form on $H$ with $\varphi(x,y)=-\varphi(y,x)$ and $b\mapsto b^\sharp$ a surjective map
$O\to O$ with $\varphi(bx,y)=\varphi(x,b^\sharp y)$, and let $\mathrm{Fil}\subseteq H$ be an $O$-stable isotropic subspace with
$2\dim\mathrm{Fil}=\dim H$. Let $w\in O$ act on $H$ commuting with $O$, with $(w-r)(w-r')=0$ on $H$ for some $r\neq r'$ in $k$ and
$\varphi(wx,y)=\varphi(x,(r+r'-w)y)$, and let $Z\subseteq(w-r')H$. If the Kodaira--Spencer directions for $(0,\mathrm{Fil}+(w-r)H,Z)$ have dimension
$d$, then so do the directions for $(\varphi,\mathrm{Fil},Z)$.
\formallabel{GlobalTheory.T.kodairaSpencer_skew_dim_semisimple_two_roots}
\formallabel{GlobalTheory.UPel.Sm.KSCount.ksDim_of_eigen_reduce}
\end{lemma}

\begin{proof}
Let $e=(r-r')^{-1}(w-r')$. It is an idempotent commuting with $O$, and $H=H_r\oplus H_{r'}$ with $H_r=eH=\ker(w-r)=(w-r')H$ and
$H_{r'}=(1-e)H=\ker(w-r')=(w-r)H$. As in step (1) of the proof of \cref{pf:ks1:prop-laglift}, $H_r$ and $H_{r'}$ are isotropic and $\varphi$
identifies $H_{r'}$ with the dual of $H_r$. Since $\mathrm{Fil}$ is $O$-stable, it is $e$-stable: $\mathrm{Fil}=\mathrm{Fil}_r\oplus\mathrm{Fil}_{r'}$
with $\mathrm{Fil}_r=\mathrm{Fil}\cap H_r$ and $\mathrm{Fil}_{r'}=\mathrm{Fil}\cap H_{r'}$. As $\mathrm{Fil}$ is isotropic of half dimension and
$\varphi$ is nondegenerate, $\mathrm{Fil}$ is its own orthogonal, whence $\mathrm{Fil}_{r'}=\{y\in H_{r'}:\varphi(\mathrm{Fil}_r,y)=0\}$. Put
$\mathrm{Fil}'=\mathrm{Fil}+H_{r'}=\mathrm{Fil}_r\oplus H_{r'}$. Directions for $(\varphi,\mathrm{Fil},Z)$ will be called $\varphi$-directions, and
directions for $(0,\mathrm{Fil}',Z)$ will be called $0$-directions.

(a) If $g$ is a $\varphi$-direction and $x\in\mathrm{Fil}_r$, then $(1-e)g(x)\in\mathrm{Fil}$: indeed $g(wx)-wg(x)=(r-w)g(x)\in\mathrm{Fil}$, and
applying $1-e$, which preserves $\mathrm{Fil}$, gives $(r-r')(1-e)g(x)\in\mathrm{Fil}$. Likewise $eg(y)\in\mathrm{Fil}$ for $y\in\mathrm{Fil}_{r'}$.
(b) If $f$ is a $0$-direction and $y\in H_{r'}$, then $f(y)\in\mathrm{Fil}'$: $f(wy)-wf(y)=(r'-w)f(y)\in\mathrm{Fil}'$, and applying $e$ gives
$(r'-r)ef(y)\in\mathrm{Fil}_r$.

(c) For a $0$-direction $f$ define $\tilde f:H\to H$ by $\tilde f=ef$ on $H_r$ and, on $H_{r'}$, as the linear map $H_{r'}\to H_{r'}$ with
$\varphi(x,\tilde fy)=-\varphi(efx,y)$ for $x\in H_r$. Then $\varphi(\tilde fx,y)+\varphi(x,\tilde fy)=0$ for all $x,y\in H$, as one checks on the
components, using that $H_r$, $H_{r'}$ are isotropic and preserved by $\tilde f$. Moreover $\tilde f$ is a $\varphi$-direction. For $x\in\mathrm{Fil}_r$
and $b\in O$, $\tilde f(bx)-b\tilde f(x)=e(f(bx)-bf(x))\in e\,\mathrm{Fil}'=\mathrm{Fil}_r$. For $y\in\mathrm{Fil}_{r'}$, $b\in O$ and $x\in\mathrm{Fil}_r$,
choose $b'$ with $b'^\sharp=b$; then $\varphi(x,\tilde f(by)-b\tilde f(y))=\varphi(\tilde f(b'x)-b'\tilde f(x),y)=0$, since
$\tilde f(b'x)-b'\tilde f(x)\in\mathrm{Fil}_r$ and $\mathrm{Fil}$ is isotropic; as $\tilde f(by)-b\tilde f(y)\in H_{r'}$, it lies in $\mathrm{Fil}_{r'}$.
Finally $\tilde f(Z)=ef(Z)\subseteq e\,\mathrm{Fil}'=\mathrm{Fil}_r$, because $Z\subseteq H_r$.

(d) Let $f_1,\dots,f_d$ be $0$-directions as in the hypothesis. We show that $\tilde f_1,\dots,\tilde f_d$ have the required property. Let $g$ be a
$\varphi$-direction. Then $g'=ege$ is a $0$-direction: for $x=x_r+x_{r'}\in\mathrm{Fil}'$ and $b\in O$,
$g'(bx)-bg'(x)=e(g(bx_r)-bg(x_r))\in e\,\mathrm{Fil}=\mathrm{Fil}_r$, and $g'(Z)=eg(Z)\subseteq\mathrm{Fil}_r$. Hence there is $c\in k^d$ with
$g'(x)-\sum_ic_if_i(x)\in\mathrm{Fil}'$ for $x\in\mathrm{Fil}'$. For $x\in\mathrm{Fil}_r$,
$g(x)-\sum_ic_i\tilde f_i(x)=(1-e)g(x)+e\bigl(g'(x)-\sum_ic_if_i(x)\bigr)\in\mathrm{Fil}$ by (a). For $y\in\mathrm{Fil}_{r'}$ put
$z=g(y)-\sum_ic_i\tilde f_i(y)$. Then $ez=eg(y)\in\mathrm{Fil}$ by (a), and for $x\in\mathrm{Fil}_r$,
$\varphi(x,(1-e)z)=\varphi(x,z)=-\varphi\bigl(g(x)-\sum_ic_i\tilde f_i(x),y\bigr)=0$, by the symmetry conditions and the isotropy of $\mathrm{Fil}$; so
$(1-e)z\in\mathrm{Fil}_{r'}$ and $z\in\mathrm{Fil}$. Thus $g(x)-\sum_ic_i\tilde f_i(x)\in\mathrm{Fil}$ for all $x\in\mathrm{Fil}$. For uniqueness let
$\sum_ic_i\tilde f_i(x)\in\mathrm{Fil}$ for all $x\in\mathrm{Fil}$. For $x\in\mathrm{Fil}_r$ this gives $e\sum_ic_if_i(x)\in\mathrm{Fil}_r$, hence
$\sum_ic_if_i(x)\in\mathrm{Fil}'$; for $y\in H_{r'}$, $\sum_ic_if_i(y)\in\mathrm{Fil}'$ by (b). So the $0$-direction $\sum_ic_if_i$ is represented both by
$c$ and by $0$, and $c=0$ by the uniqueness in the hypothesis.
\end{proof}

\begin{proposition}[Dimension of the Kodaira--Spencer directions]\label{pf:ks1:prop-ksdim}
Let $b$ be a base for $D$ over $R$, $N\ge3$ prime to $p$, $k$ an algebraically closed field with $\beta:R\to k$, and $o_0$ an object of level $N$ over
$\ksoneSp k$. Let $H=H^{\mathrm{dR}}_1(A_{o_0})$ with the action of $O$ induced by $\iota_{o_0}$, $\varphi$ a form on $H$ induced by $P_{o_0}$, and
$\mathrm{Fil}=\mathrm{Fil}(A_{o_0})$. Write $\rho$ for $\beta(\rho)$, put $\rho'=\mathrm{tr}-\rho$ and $H_\rho=(\omega-\rho')H$, and let $Z\subseteq H$
be the span of the elements $\sum_l\lambda_lc_ly$ with $y\in H_\rho$, $c_l\in\OO_F$ and $\lambda_l\in k$ such that $\sum_l\lambda_l\beta(t_i(c_l))=0$
for every $i<a$. Then the Kodaira--Spencer directions for $(\varphi,\mathrm{Fil},Z)$ have dimension $a(n-1)$.
\formallabel{GlobalTheory.T.harrisTaylor_kodairaSpencer_directions_dim}
\formallabel{GlobalTheory.HTCount.ksDim_at_point}
\end{proposition}

\begin{proof}
(1) The space $H$ is finite-dimensional, $2\dim\mathrm{Fil}=\dim H$, $\varphi$ is skew-symmetric and $\mathrm{Fil}$ is isotropic
(\lref{GlobalTheory.T.harrisTaylor_pairing_skew_hodge_lagrangian}); $\varphi$ is nondegenerate
(\lref{GlobalTheory.T.harrisTaylor_polarisation_pairing_nondegenerate}); and $\varphi(bx,y)=\varphi(x,b^*y)$ by the Rosati condition
(\lref{GlobalTheory.T.harrisTaylor_crystal_action_adjoint}). The subspace $\mathrm{Fil}$ is $O$-stable, and $b\mapsto b^*$ is surjective since
$b^{**}=b$.

(2) As in step (1) of the proof of \cref{pf:ks1:prop-laglift}, $(\omega-\rho)(\omega-\rho')=0$ on $H$, and $\rho\neq\rho'$ in $k$, because
$(\rho-\rho')^2=\mathrm{tr}^2-4\,\mathrm{nm}$ is the image of an integer prime to $p$, which is a unit of $R$. The element $\omega$ is central in $O$, and
$\varphi(\omega x,y)=\varphi(x,(\rho+\rho'-\omega)y)$ since $\omega^*=\mathrm{tr}-\omega$. Each $c_l$ commutes with $\omega-\rho'$, so
$Z\subseteq(\omega-\rho')H$.

(3) By (1), (2) and \cref{pf:ks1:lem-ksreduce} with $w=\omega$, $r=\rho$ and $r'=\rho'$, it suffices that the Kodaira--Spencer directions for
$(0,\mathrm{Fil}+(\omega-\rho)H,Z)$ have dimension $a(n-1)$. This is \lref{GlobalTheory.T.harrisTaylor_rho_part_deformation_dim} (the
$O$-linear deformations of $\mathrm{Fil}+(\omega-\rho)H$ which vanish on $Z$ form a space of dimension $a(n-1)$); its formal proof identifies
$H/(\mathrm{Fil}+(\omega-\rho)H)$ with the part of $\Lie A_{o_0}$ on which $\omega$ acts by $\rho$, uses Kottwitz's condition (O3) and a Morita
equivalence for the action of $O$, and counts dimensions.
\end{proof}

\begin{lemma}[First-order deformations]\label{pf:ks1:lem-fodef}
Let $b$, $N$, $k$, $o_0$, $H$ and $\mathrm{Fil}$ be as in \cref{pf:ks1:prop-ksdim} and $d=a(n-1)$. There are linear maps $f_1,\dots,f_d:H\to H$ such
that each $f_j$ is a Kodaira--Spencer class of some deformation of $o_0$ to $k[\varepsilon]$, and every deformation of $o_0$ to $k[\varepsilon]$ has
$\sum_jc_jf_j$ as a Kodaira--Spencer class for a unique $c\in k^d$.
\formallabel{GlobalTheory.T.harrisTaylor_firstOrder_deformations_basis_count}
\formallabel{GlobalTheory.HTCount.count}
\end{lemma}

\begin{proof}
Choose a form $\varphi$ on $H$ induced by $P_{o_0}$ (\lref{GlobalTheory.T.harrisTaylor_exists_deRham_pairing}), let $Z$ be as in
\cref{pf:ks1:prop-ksdim}, and let $f_1,\dots,f_d$ be Kodaira--Spencer directions for $(\varphi,\mathrm{Fil},Z)$ as in the conclusion of
\cref{pf:ks1:prop-ksdim}. Every Kodaira--Spencer direction is a Kodaira--Spencer class of some deformation of $o_0$ to $k[\varepsilon]$
(\lref{GlobalTheory.T.harrisTaylor_realise_kodairaSpencer_direction}), and every deformation of $o_0$ to $k[\varepsilon]$ has a Kodaira--Spencer
class which is a Kodaira--Spencer direction (\lref{GlobalTheory.T.harrisTaylor_firstOrder_hodge_eq_kodairaSpencer_graph}); both statements
are Grothendieck--Messing theory, the three conditions of \cref{pf:ks1:def-ksdir} corresponding to the lifting of the action, of the polarisation
and to Kottwitz's condition. Let $o'$ be a deformation and $g$ a class of $o'$ which is a direction. By \cref{pf:ks1:prop-ksdim} there is a unique $c$
with $g(x)-\sum_jc_jf_j(x)\in\mathrm{Fil}$ for $x\in\mathrm{Fil}$. Since the classes of $o'$ are exactly the maps congruent to $g$ modulo $\mathrm{Fil}$
on $\mathrm{Fil}$, $\sum_jc_jf_j$ is a class of $o'$, and any $c'$ for which $\sum_jc'_jf_j$ is a class of $o'$ satisfies the same congruence, hence
$c'=c$.
\end{proof}

\begin{lemma}[Universal first-order deformations]\label{pf:ks1:lem-univdef}
Let $b$ be a base for $D$ over $R$, $N\ge3$ prime to $p$, $k$ an algebraically closed field with $\ksoneSp k\to\ksoneSp R$, $o_0$ an object of level
$N$ over $\ksoneSp k$, and $d=a(n-1)$. For a $k$-linear map $c:k^d\to k$ let $\mathrm{id}+c:k\oplus k^d\to k[\varepsilon]$, $(x,v)\mapsto
x+c(v)\varepsilon$. There is a deformation $o$ of $o_0$ to $k\oplus k^d$ such that for every deformation $o'$ of $o_0$ to $k[\varepsilon]$ there is a
unique $c$ for which there exists a morphism $o'\to o$ over $\ksoneSp(\mathrm{id}+c)$.
\formallabel{GlobalTheory.T.exists_universal_firstOrder_deformation_harrisTaylor}
\formallabel{GlobalTheory.HTUniv.main}
\end{lemma}

\begin{proof}
Let $f_1,\dots,f_d$ be as in \cref{pf:ks1:lem-fodef}. Deformations of $o_0$ to $k[\varepsilon]$ with Kodaira--Spencer classes $f_1,\dots,f_d$ can be
glued: there is a deformation $o$ of $o_0$ to $k\oplus k^d$ with Kodaira--Spencer class $(f_1,\dots,f_d)$
(\lref{GlobalTheory.T.harrisTaylor_exists_glued_firstOrder_deformation}). We use two further facts: if a deformation $o$ to $k\oplus k^d$ has
class $(f_j)_j$ and $o'\to o$ is a morphism over $\ksoneSp(\mathrm{id}+c)$ from a deformation $o'$ to $k[\varepsilon]$, then $\sum_jc(\varepsilon_j)f_j$
is a class of $o'$ (\lref{GlobalTheory.T.harrisTaylor_hodgeFiltration_naturality}, naturality of the Hodge filtration); and two deformations of
$o_0$ to $k[\varepsilon]$ with a common Kodaira--Spencer class are isomorphic as deformations
(\lref{GlobalTheory.T.harrisTaylor_firstOrder_iso_of_same_hodge_filtration}, Grothendieck--Messing theory).

Let $o'$ be a deformation to $k[\varepsilon]$, let $c_0\in k^d$ be such that $\sum_jc_{0,j}f_j$ is a class of $o'$ (\cref{pf:ks1:lem-fodef}), and
$c(v)=\sum_jc_{0,j}v_j$. The pull-back $o_c$ of $o$ along $\ksoneSp(\mathrm{id}+c)$ is a deformation of $o_0$ (the composite of the closed point with
$\mathrm{id}+c$ is the closed point of $k\oplus k^d$) with a morphism $o_c\to o$ over $\ksoneSp(\mathrm{id}+c)$, so it has the class
$\sum_jc(\varepsilon_j)f_j=\sum_jc_{0,j}f_j$. Hence $o'\cong o_c$, and composing gives a morphism $o'\to o$ over $\ksoneSp(\mathrm{id}+c)$. If
$o'\to o$ is a morphism over $\ksoneSp(\mathrm{id}+c')$, then $\sum_jc'(\varepsilon_j)f_j$ is a class of $o'$, so $c'(\varepsilon_j)=c_{0,j}$ for all $j$
by the uniqueness in \cref{pf:ks1:lem-fodef}, i.e.\ $c'=c$.
\end{proof}

\begin{proposition}[Universal tangents of representing families]\label{pf:ks1:prop-univtan}
Let $R$ be a Noetherian ring, $b$ a base for $D$ over $R$, $N\ge3$ prime to $p$, and $(X,s,u)$ a representing family of level $N$ with $s$
quasi-projective and smooth. Let $k$ be an algebraically closed field, $x:\ksoneSp k\to X$ and $d=a(n-1)$. There is $u_x:\ksoneSp(k\oplus k^d)\to X$
which restricts to $x$ on the closed point and lies over $sx$, such that every $v:\ksoneSp k[\varepsilon]\to X$ which restricts to $x$ and lies over
$sx$ is $u_x\circ\ksoneSp(\mathrm{id}+c)$ for a unique $k$-linear $c:k^d\to k$.
\formallabel{GlobalTheory.T.exists_universal_tangent_harrisTaylor_representingFamily}
\formallabel{GlobalTheory.UPel.HT.RepTangent.exists_univ_tangent}
\end{proposition}

\begin{proof}
Let $o_0$ be the pull-back of $u$ along $x$, an object over $\ksoneSp k$ with a morphism $o_0\to u$ over $x$
(\lref{GlobalTheory.T.exists_harrisTaylorObject_pullback_along}), and let $o$ be the deformation of \cref{pf:ks1:lem-univdef}. By the universal
property, $o$ is classified by $u_x:\ksoneSp(k\oplus k^d)\to X$ over $R$, with a morphism $o\to u$ over $u_x$; $u_x$ lies over $sx$. Its restriction
to the closed point classifies $o_0$ (compose $o_0\to o\to u$), as does $x$, so the two agree by uniqueness. Let $v$ be as in the statement and
$o'=v^*u$. Since $v$ restricts to $x$, the morphism $o_0\to u$ over $x$ and the cartesian morphism $o'\to u$ over $v$ give a morphism $o_0\to o'$
over the closed point, so $o'$ is a deformation of $o_0$. By \cref{pf:ks1:lem-univdef} there is a unique $c$ with a morphism $o'\to o$ over
$\ksoneSp(\mathrm{id}+c)$. Then $v$ and $u_x\circ\ksoneSp(\mathrm{id}+c)$ both lie over $R$ and admit morphisms from $o'$ to $u$, so they are equal by
the uniqueness in the universal property. Conversely, if $v=u_x\circ\ksoneSp(\mathrm{id}+c')$, then the morphism $o'\to u$ over $v$ factors
through the cartesian morphism $o\to u$ over $u_x$, which gives a morphism $o'\to o$ over $\ksoneSp(\mathrm{id}+c')$, and $c'=c$ by the uniqueness in
\cref{pf:ks1:lem-univdef}.
\end{proof}

% ---------------------------------------------------------------------------------------------
\subsection{Hecke towers and their Hodge cohomology}\label{pf:ks1:sec-towers}

\paragraph{Hodge cohomology.} For a scheme $s:Z\to\ksoneSp k$ over a field $k$ let $\Omega^a_{Z/k}$ be the
sheaf of $a$-forms and $\ksoneH{Z}{a}{b}$ its sheaf cohomology. For $f:Z'\to Z$ over $k$ there is a pull-back
$f^*:\ksoneH{Z}{a}{b}\to\ksoneH{Z'}{a}{b}$, functorial in $f$, and for $f$ finite \'etale a trace
$f_*:\ksoneH{Z'}{a}{b}\to\ksoneH{Z}{a}{b}$ induced by the trace of differential forms (these are definitions of
the formal development \cite{LeanSource2026}). If $k$ has characteristic $0$, then (i) $f^*$ is injective when $f$ is a
finite \'etale Galois cover, and (ii) $f_*$ is a trace in the following sense: for every $c:W\to Z$ and every finite
family of sections $\sigma_j:W\to Z'$ of $f$ over $c$ which split $f$ over $W$ (that is, $W\times_ZZ'$ is the disjoint
union of the images of the $\sigma_j$), one has $c^*f_*x=\sum_j\sigma_j^*x$
(\lref{UpperTheory.Gal.G3.pHodgeTr}, \lref{UpperTheory.Gal.G3.pHodgeGood}, and \lref{UpperTheory.Gal.G3.pHodgeInj}).

\paragraph{Correspondences.} Let $Z\xleftarrow{s}Y\xrightarrow{t}Z'$ be morphisms over $k$ with $s$ finite \'etale.
A linear map $A:\ksoneH{Z'}{a}{b}\to\ksoneH{Z}{a}{b}$ is \emph{the operator of $(s,t)$} if there are a finite \'etale
Galois cover $c:W\to Z$ and sections $\sigma_j:W\to Y$ of $s$ over $c$ splitting $s$ over $W$ with
$c^*(Ax)=\sum_j(t\sigma_j)^*x$ for all $x$. By (i) and (ii), if $(s,t)$ has an operator $A$, then
$A=s_*\circ t^*$ (\lref{UpperTheory.Gal.G3.pOpTr}).

\begin{definition}[Towers]\label{pf:ks1:def-tower}
Let $B$ be a scheme and $\Gamma$ a group. A \emph{tower} $\mathcal X$ over $B$ for $\Gamma$ consists of an index set $I$,
subgroups $K_i\subseteq\Gamma$ ($i\in I$) and a submonoid $\Delta\subseteq\Gamma$ containing every $K_i$, schemes
$q_i:Z_i\to B$, and for $h\in\Delta$ with $hK_ih^{-1}\subseteq K_j$ (we say that $h$ is \emph{admissible} for
$(i,j)$) morphisms $\mathrm{tr}(i,j,h):Z_i\to Z_j$ over $B$, such that: for every $j$ and every finite set of pairs
$(i,g)\in I\times\Gamma$ there is $m$ with $K_m\subseteq K_j$ normalised by $K_j$ and $gK_mg^{-1}\subseteq K_i$
for all the pairs; all indices $[K_j:K_i\cap K_j]$ are finite; $\mathrm{tr}(i,i,k)=\mathrm{id}$ for $k\in K_i$ and
$\mathrm{tr}(j,l,h')\mathrm{tr}(i,j,h)=\mathrm{tr}(i,l,h'h)$; $\mathrm{tr}(i,j,1)$ is finite, \'etale and surjective;
if $K_i$ is normalised by $K_j$, two geometric points of $Z_i$ with the same image under $\mathrm{tr}(i,j,1)$ differ by
some $\mathrm{tr}(i,i,k)$, $k\in K_j$; and if $K_i\subseteq K_j$ and $k\in K_j$ fixes a geometric point of $Z_i$ then
$k\in K_i$.
\end{definition}

A function $f:\Gamma\to\C$ is a \emph{Hecke function at level $j$} if it is bi-$K_j$-invariant, $g^{-1}\in\Delta$
whenever $f(g)\neq0$, and $f$ is supported on finitely many cosets $gK_j$. For such $f$ and a good cohomology
theory such as Hodge cohomology over a field of characteristic $0$, the \emph{Hecke operator} $T(f)$ on
$\ksoneH{Z_j}{a}{b}$ is the endomorphism characterised by
\begin{equation}\label{pf:ks1:eq-hecke}
\mathrm{tr}(m,j,1)^*\,T(f)x=\ksonefs_{y\in\Gamma/K_j}f(y)\,\mathrm{tr}(m,j,y^{-1})^*x
\end{equation}
for every $m$ with $K_m\subseteq K_j$ such that $g^{-1}$ is admissible for $(m,j)$ whenever $f(g)\neq0$; it is
unique by (i), and the formal definition sets $T(f)=0$ if no such endomorphism exists. For $h\in\Gamma$ put
$\mathrm{lev}(K,h)=K\cap h^{-1}Kh$.

\begin{lemma}[Double cosets as correspondences]\label{pf:ks1:lem-hopop}
Let $\mathcal X$ be a tower, $k$ a field of characteristic $0$ and $a,b\ge0$. Let $i,j\in I$ and $h\in\Delta$ be such that
$1$ and $h$ are admissible for $(i,j)$. Then $\delta:=[\mathrm{lev}(K_j,h):K_i]$ is finite and nonzero, and
$\delta\cdot T(\mathbf 1_{K_jh^{-1}K_j})$ is the operator of the correspondence
$(\mathrm{tr}(i,j,1),\mathrm{tr}(i,j,h))$ on $\ksoneH{Z_j}{a}{b}$. Consequently
$T(\mathbf 1_{K_jh^{-1}K_j})=\delta^{-1}\,\mathrm{tr}(i,j,1)_*\,\mathrm{tr}(i,j,h)^*$.
\formallabel{UpperTheory.Gal.LocalSmoothExtension.Tw.hopOp}
\formallabel{UpperTheory.Gal.G3.pHopOp}
\end{lemma}

\begin{proof}
Admissibility gives $K_i\subseteq\mathrm{lev}(K_j,h)\subseteq K_j$, and $[K_j:K_i]$ is finite and nonzero, so $\delta$ is
finite and nonzero. Choose $m$ with $K_m\subseteq K_i$ normalised by $K_j$. Then $\mathrm{tr}(m,j,1)$ is a finite \'etale Galois
cover with group $K_j/K_m$ acting through the $\mathrm{tr}(m,m,k)$ (\lref{UpperTheory.Gal.G3.pTwGal}), and the maps
$\sigma_c=\mathrm{tr}(m,i,c^{-1})$, $c\in K_j/K_i$, are sections of $\mathrm{tr}(i,j,1)$ over $\mathrm{tr}(m,j,1)$ which split it
(\lref{UpperTheory.Gal.G3.pTwSplit}). The function $\mathbf 1_{K_jh^{-1}K_j}$ is a Hecke function and its operator
satisfies \eqref{pf:ks1:eq-hecke} (\lref{UpperTheory.Gal.G3.pFnRing} and \lref{UpperTheory.Gal.G3.pHopEx}). Moreover
$\mathrm{tr}(i,j,h)\sigma_c=\mathrm{tr}(m,j,hc^{-1})$. If $c,c'\in K_j$, then $ch^{-1}K_j=c'h^{-1}K_j$ if and only if
$c'^{-1}c\in\mathrm{lev}(K_j,h)$; hence as $c$ runs through $K_j/K_i$ the cosets $ch^{-1}K_j$ run through the left cosets
$yK_j\subseteq K_jh^{-1}K_j$, each exactly $\delta$ times. Since $\mathrm{tr}(m,j,(yk)^{-1})=\mathrm{tr}(j,j,k^{-1})\mathrm{tr}(m,j,y^{-1})=
\mathrm{tr}(m,j,y^{-1})$ for $k\in K_j$, \eqref{pf:ks1:eq-hecke} gives
$\mathrm{tr}(m,j,1)^*(\delta\,T(\mathbf 1_{K_jh^{-1}K_j})x)=\sum_c(\mathrm{tr}(i,j,h)\sigma_c)^*x$, which is the defining property of
the operator of the correspondence. The last assertion follows from \lref{UpperTheory.Gal.G3.pOpTr}.
\end{proof}

\paragraph{Kuga--Sato towers.} Let $D$ be a datum, $G$ a unitary similitude group and $L$ a link between them (see
\cref{pf:ks1:sec-groups}; $L$ identifies $F_D$ with $F_G$ and realises $G$ as the group of $O$-linear similitudes of
$(\Lambda_\Q,E_\Lambda)$). Let $R$ be a field of characteristic $0$, $g_b$ a base for $D$ over $R$ in the
characteristic-$0$ sense (no condition at $p$) and $k\ge0$. A \emph{Kuga--Sato tower} $T$ for $(L,g_b,k)$ is a tower over
$\ksoneSp R$ for $\Gamma=G(\ksoneAfp)$ whose index set is the set of levels $N\ge3$ prime to $p$, with $K_N$ the principal
congruence subgroup of level $N$ and $\Delta$ the monoid of integral elements, together with fine moduli schemes $M_N$ of
level $N$ over $R$ with universal abelian scheme $\mathcal A_N\to M_N$ and isomorphisms
$Z_N\cong\mathcal A_N^k=\mathcal A_N\times_{M_N}\dots\times_{M_N}\mathcal A_N$ over $R$, under which the transition maps are the
$k$-th fibre powers of Hecke isogenies. For $\alpha=(\alpha_1,\dots,\alpha_k)\in\OO_F^k$ let $[\alpha]:Z_N\to Z_N$ be the
endomorphism acting by $\alpha_s$ on the $s$-th factor. A \emph{Hecke datum} of level $N$ is $g=(K_g,h_g,M_g)$ with
$K_N\subseteq K_g\subseteq K_1$, $h_g$ integral, $N\mid M_g$ and $h_g\kappa K_{M_g}(h_g\kappa)^{-1}\subseteq K_N$ for
$\kappa\in K_g$; its Hecke function is $f_g(y)=\#\{\kappa\in K_g/(K_{M_g}\cap K_g):h_g\kappa y\in K_N\}$. For a family
$(g_j)_{j\in J}$ of Hecke data, $\mathcal H_T(N,g)$ is the Hecke tower with base $Z_N$, correspondences
$Z_N\xleftarrow{s_j}Y_j\xrightarrow{t_j}Z_N$ where $Y_j=\coprod_{\kappa\in K_{g_j}/(K_{M_{g_j}}\cap K_{g_j})}Z_{M_{g_j}}$,
$s_j$ is $\mathrm{tr}(M_{g_j},N,1)$ on each copy and $t_j$ is $\mathrm{tr}(M_{g_j},N,h_{g_j}\kappa)$ on the copy indexed by
$\kappa$. For a Hecke tower $H=(Z,(Y_j,s_j,t_j)_j)$ over $k$, $\theta:J\to k$ and $a\ge0$, \emph{Hodge occurrence} $\mathrm{HO}(H,\theta,a)$ means:
there are $b$ and $0\neq x\in\ksoneH{Z}{a}{b}$ such that every $s_j$ is finite \'etale and $(s_j)_*t_j^*x=\theta_jx$ for all
$j$. For a field embedding $\tau:k\to K$, $\mathrm{bc}_\tau(H)$ denotes the base change.

\begin{definition}[Holomorphic Lefschetz formula]\label{pf:ks1:def-hlf}
$\sfname{HLF}$ is the following statement. Let $s:Z\to\ksoneSp\C$ be projective and smooth of relative dimension $d$,
$Y\to\ksoneSp\C$ a scheme and $c_1,c_2:Y\to Z$ finite \'etale morphisms over $\C$. Let $\mathcal F$ be the finite set
of $y\in Y(\C)$ with $c_1(y)=c_2(y)$ and let $A_y\in M_d(\C)$ ($y\in\mathcal F$) be \emph{tangent matrices}: in local
coordinates $z_1,\dots,z_d$ of $Z$ at $c_1(y)$, the differentials $d(z_l\circ c_1)$ form a basis of the cotangent space of $Y$
at $y$ and $d(z_j\circ c_2)=\sum_l(A_y)_{jl}\,d(z_l\circ c_1)$; assume $\det(1-A_y)\neq0$. Then for every $a$ the spaces
$\ksoneH{Z}{a}{b}$ are finite-dimensional and
\begin{equation}\label{pf:ks1:eq-hlf}
\sum_{b=0}^{d}(-1)^b\tr\bigl((c_1)_*c_2^*\mid\ksoneH{Z}{a}{b}\bigr)=\sum_{y\in\mathcal F}\frac{[X^a]\det(1+XA_y)}{\det(1-A_y)}.
\end{equation}
\end{definition}

This is the holomorphic Lefschetz formula of Atiyah--Bott \cite{AtiyahBott1968} for finite \'etale correspondences with
transversal fixed points, in the form of \cref{thm:galois:hlf} (note $\tr\bigwedge^aA_y=[X^a]\det(1+XA_y)$); in the formal development it is proved
by \lref{UpperTheory.Gal.Int.Geo.weight_one_geometric_realization}, which establishes exactly this statement. For $d\ge0$, $A\in M_d(\C)$ and $y\in\C$ we put
$\mathrm{lef}(A,y)=\det(1-yA)/\det(1-A)$, and for $s:Z\to\ksoneSp\C$ and a family $u=(u_{a,b})$ of endomorphisms of the
$\ksoneH{Z}{a}{b}$,
\begin{equation}\label{pf:ks1:eq-alttr}
\mathcal L_d(u)(y)=\sum_{a=0}^d(-y)^a\sum_{b=0}^d(-1)^b\tr(u_{a,b}\mid\ksoneH{Z}{a}{b}).
\end{equation}

\begin{lemma}\label{pf:ks1:lem-lefsum}
For $A\in M_d(\C)$ with $\det(1-A)\neq0$ and $y\in\C$,
$\sum_{a=0}^d(-y)^a\,[X^a]\det(1+XA)/\det(1-A)=\mathrm{lef}(A,y)$.
\formallabel{UpperTheory.Gal.LocalSmoothExtension.lefSum_of}
\end{lemma}

\begin{proof}
$\det(1+XA)$ is a polynomial of degree at most $d$ in $X$, and its value at $X=-y$ is $\det(1-yA)$.
\end{proof}

\begin{lemma}\label{pf:ks1:lem-hodgefin}
$\sfname{HLF}$ implies: if $s:Z\to\ksoneSp\C$ is projective and smooth of some relative dimension, then every
$\ksoneH{Z}{a}{b}$ is finite-dimensional.
\formallabel{UpperTheory.Gal.LocalSmoothExtension.hodgeFin_of}
\end{lemma}

\begin{proof}
Apply $\sfname{HLF}$ with $Y=\emptyset$ (so $\mathcal F=\emptyset$); the finite-dimensionality is part of its conclusion.
\end{proof}


% ---------------------------------------------------------------------------------------------
\subsection{Unitary similitude groups, fixed point sums and definite traces}\label{pf:ks1:sec-groups}

\begin{definition}[Unitary similitude groups]\label{pf:ks1:def-htgroup}
Let $p$ be a prime. A \emph{unitary similitude group of Harris--Taylor type} $G$ for $p$ consists of: a CM field $F$;
an imaginary quadratic subfield $K_q\subseteq F$ with an embedding $u:K_q\to\Q_p$; an integer $n\ge1$; a central
simple $F$-algebra $C$ of dimension $n^2$ with an involution $x\mapsto x^*$ of the second kind (a ring
anti-automorphism of order $2$ inducing complex conjugation on $F$); for each infinite place $v$ of $F$ a homomorphism
$\psi_v:C\to M_n(\C)$ which is the scalar embedding given by $v$ on $F$, and an invertible hermitian matrix
$\mathrm{gr}_v$ with $\psi_v(x^*)=\mathrm{gr}_v^{-1}\psi_v(x)^\dagger\mathrm{gr}_v$, of signature $(1,n-1)$ or $(0,n)$;
embeddings $\epsilon_i:F\to\Qbar_p$ ($i<m$) lying over $u$ and representing each $\Q_p$-conjugacy class of such
embeddings exactly once, with homomorphisms $\rho_i:C\to M_n(F_{w_i})$ extending $\epsilon_i$, where
$F_{w_i}=\Q_p(\epsilon_i(F))$; and embeddings of $\Qbar$ into $\Qbar_p$ and into $\C$. We write $a_G$ for the number of
infinite places of signature $(1,n-1)$.
\end{definition}

Put $G(\ksoneAfp)=\{(g,\nu)\in(\ksoneAfp\otimes_\Q C)^\times\times(\ksoneAfp)^\times: gg^*=\nu\}$,
$G_p=\Q_p^\times\times\prod_{i<m}\GL_n(F_{w_i})$ and $K_p=\Z_p^\times\times\prod_{i<m}\GL_n(\OO_{F_{w_i}})$, and let
$\mathbf 1_{K_p}$ be the indicator function of $K_p$. A \emph{level} is a compact open subgroup of $G(\ksoneAfp)$. We
write $(C^\times\times\Q^\times)\to G_p$, $(y,s)\mapsto(s,(\rho_i(y))_i)$. A function $f$ on $G(\ksoneAfp)$ is a \emph{test
function} for a level $K$ if $f(k^{-1}gk)=f(g)$ for $k\in K$ and $f$ is supported on finitely many double cosets $KgK$;
a function $f'$ on $G_p$ is a test function for a subgroup $K'$ commensurable with $K_p$ if it is invariant under
conjugation by $K'$ and supported on finitely many double cosets $K_pgK_p$. An element $g\in G(\ksoneAfp)$ is
\emph{regular} if for some prime $\ell\neq p$ the centraliser of its component in $\Q_\ell\otimes C$ is commutative. An element
$x\in C$ is \emph{regular irreducible} if its minimal polynomial $\chi_x$ over $F$ has degree $n$ and is separable and
irreducible.

A \emph{link} $L$ between a datum $D$ and $G$ identifies $F_D\cong F_G$, $n_G=n_D$, realises $C$ as the centraliser of $O$ in
$\End_\Q(\Lambda_\Q)$ compatibly with the involutions, matches the places of signature $(1,n-1)$ with the embeddings
$\tau_i$ of $D$, and fixes compatible embeddings at $p$ and an orientation. Through $L$ the group $G(\ksoneAfp)$ acts on
$\ksoneAfp\otimes\Lambda$; $K_N$ denotes the principal congruence subgroup of level $N$ (the $g$ such that $g$ and $g^{-1}$
are $\equiv1$ modulo $N$ on $\hat\Z^p\otimes\Lambda$), and $h$ is \emph{integral} if it preserves $\hat\Z^p\otimes\Lambda$.

\begin{definition}[Forms and fixed point sums]\label{pf:ks1:def-forms}
Let $G$ be as above.
\begin{enumerate}
\item A \emph{form} is $\beta\in C^\times$ with $\beta^*=\beta$ for which there are $x\in(\ksoneAfp\otimes C)^\times$,
  $t\in(\ksoneAfp)^\times$ with $x\beta x^*=t$ and $z\in(\R\otimes C)^\times$ with $z\beta z^*=1$. Forms $\beta,\beta'$ are
  \emph{equivalent} if $\beta'=s\,y\beta y^*$ with $y\in C^\times$ and $s\in\Q_{>0}$; $\mathrm{Cl}_{\mathrm{form}}$ is the set of classes.
\item For a form $\beta$ let $\Gamma_\beta=\{(y,s)\in C^\times\times\Q^\times: y\beta y^*=s\beta,\ s>0\}$. For a choice of
  $(x,t)$ as in (1) let $\iota_f:\Gamma_\beta\to G(\ksoneAfp)$, $(y,s)\mapsto(xyx^{-1},s)$, and $\iota_p:\Gamma_\beta\to G_p$
  the map above; let $\iota_v(y)=\psi_v(y)$. Put $\mathrm{herm}_v(\beta)=\mathrm{gr}_v\psi_v(\beta^{-1})$ and
  $\mathcal P_\beta=\{z\in\C^n:\Re(z^\dagger\mathrm{herm}_v(\beta)z)>0\}$.
\item Let $K$ be a level and $K'\subseteq G_p$ a subgroup. $\mathcal X_\beta$ is the quotient of
  $(G(\ksoneAfp)\times G_p)\times\mathcal P_\beta$ by the relation $(g,z)\sim((\iota_f(y),\iota_p(y))\,g\,k,\ c\,\iota_v(y)z)$
  for $y\in\Gamma_\beta$, $k\in K\times K'$, $c\in\C^\times$. For a representative $(g,z)$ the \emph{stabiliser} is the set of
  $y\in\Gamma_\beta$ with $g_f^{-1}\iota_f(y)g_f\in K$, $g_p^{-1}\iota_p(y)g_p\in K'$ and $\iota_v(y)z\in\C z$.
\item For $c:F[X]\times\C\to\C$, a function $f$ on $G(\ksoneAfp)$ and $f'$ on $G_p$, put
  \[\Sigma_\beta(c;K,f;K',f')=\ksonefs_{q\in\mathcal X_\beta}\frac{1}{|\mathrm{Stab}(g,z)|}\ksonefs_{y}
  f(g_f^{-1}\iota_f(y)g_f)\,f'(g_p^{-1}\iota_p(y)g_p)\,c\bigl(\chi_y,\mu_z(y)\bigr),\]
  where $(g,z)$ is a chosen representative of $q$, $y$ runs over the elements of $\Gamma_\beta$ with $\iota_v(y)z=\mu_z(y)z$
  for some (necessarily unique) $\mu_z(y)\in\C$, and $\chi_y$ is the minimal polynomial of $y$ over $F$; and
  \[\Sigma(c;K,f;K',f')=\ksonefs_{[\beta]\in\mathrm{Cl}_{\mathrm{form}}}\Sigma_\beta(c;K,f;K',f')\]
  with chosen representatives $\beta$ and chosen $(x,t)$. The \emph{finiteness condition} $\mathrm{Fin}(K,f)$ means that, for $K'=K_p$, the set of triples
  $([\beta],q,y)$ with $y$ as above and
  \[f(g_f^{-1}\iota_f(y)g_f)\,\mathbf 1_{K_p}(g_p^{-1}\iota_p(y)g_p)\neq0\]
  is finite.
\end{enumerate}
Here $v$ is a fixed infinite place; in all applications it is the unique place of signature $(1,n-1)$.
\end{definition}

When $v$ has signature $(1,n-1)$ and $\beta$ is a form, $\mathcal P_\beta/\C^\times$ is a complex ball and
$\mathcal X_\beta$ is the set of complex points of the corresponding component of a unitary Shimura variety; $y$ runs
over the elements fixing the point, $\mu_z(y)$ is the eigenvalue of $y$ on the fixed line, and $\Sigma$ is a weighted
count of fixed points of Hecke correspondences. These interpretations are not used in the proofs below.

\begin{definition}[Pairs and terms]\label{pf:ks1:def-pairs}
\begin{enumerate}
\item A \emph{pair} is $P=(\beta,\gamma,c)$ with $\beta$ a form, $\gamma\in C^\times$, $c\in\Q^\times$, $\gamma\beta\gamma^*=c\beta$,
  and $\chi_P:=\chi_\gamma$ of degree $n$ and separable. A \emph{definite pair} is the same with the condition at infinity
  replaced by: for every infinite place $w$ the hermitian form $z\mapsto z^\dagger\mathrm{herm}_w(\beta)z$ is positive
  definite if $w$ has signature $(1,n-1)$ and negative definite otherwise. Pairs $P,P'$ are equivalent if
  $\beta'=s\,y\beta y^*$ and $\gamma'=y\gamma y^{-1}$ with $y\in C^\times$, $s\in\Q_{>0}$; $\mathrm{Cl}_{\mathrm{pair}}$ and $\mathrm{Cl}^{\mathrm{def}}_{\mathrm{pair}}$
  are the sets of classes of pairs and of definite pairs. A \emph{definite form} is $\beta$ as in a definite pair;
  $\mathrm{Cl}^{\mathrm{def}}_{\mathrm{form}}$ is the set of classes of definite forms.
\item A pair is \emph{marked} by $\mu\in\C$ if for some place $w$ of signature $(1,n-1)$ there is $z\in\C^n$ with
  $\psi_w(\gamma)z=\mu z$ and $\Re(z^\dagger\mathrm{herm}_w(\beta)z)>0$.
\item For a (definite) pair $P$ let $\mathrm{Cent}_P=\{(y,s)\in\Gamma_\beta: y\gamma=\gamma y\}$, embedded into
  $G(\ksoneAfp)\times G_p$ by $(\iota_f,\iota_p)$ (with the $p$-component conjugated by a chosen $z_p\in G_p$). The
  \emph{term} of $P$ is the orbital sum
  \[\mathrm{T}_P(K,f;K',f')=\ksonefs_{q}\frac{1}{|\mathrm{stab}(q)|}\,f(q_f^{-1}\gamma_fq_f)\,f'(q_p^{-1}\delta q_p),\]
  where $q$ runs over $\mathrm{Cent}_P\backslash(G(\ksoneAfp)\times G_p)/(K\times K')$ with chosen representatives,
  $\mathrm{stab}(q)=\{y\in\mathrm{Cent}_P:\ q^{-1}yq\in K\times K'\}$, $\gamma_f=\iota_f(\gamma,c)$ and $\delta=z_p^{-1}\iota_p(\gamma,c)z_p$.
\item For $c:F[X]\times\C\to\C$ put
  $\Sigma^{\mathrm{mk}}(c;K,f;K',f')=\ksonefs_{[P]\in\mathrm{Cl}_{\mathrm{pair}}}c(\chi_P,\mu_P)\,\mathrm T_P(K,f;K',f')$, the summand being $0$
  unless $P$ is marked, $\mu_P$ a chosen marking; for $\mathrm{wt}:F[X]\to\C$ and a root selection $\mathrm{sel}$ (a map
  $F[X]\to\C$ such that $\mathrm{sel}(Q)$ is a root of $Q$ at every place of signature $(1,n-1)$, for $Q$ monic of degree $n$)
  put $\Sigma^{\mathrm{mk}}_{\mathrm{sel}}(\mathrm{wt})=\ksonefs_{[P]}\mathrm{wt}(\chi_P)\,\mathrm T_P$ over the classes
  with $\chi_P$ irreducible and $P$ marked by $\mathrm{sel}(\chi_P)$; and
  $\Sigma^{\mathrm{def}}(\mathrm{wt};K,f;K',f')=\ksonefs_{[P]\in\mathrm{Cl}^{\mathrm{def}}_{\mathrm{pair}},\ \chi_P\text{ irreducible}}\mathrm{wt}(\chi_P)\,\mathrm T_P(K,f;K',f')$.
\end{enumerate}
\end{definition}

\begin{definition}[Algebraic modular forms and definite traces]\label{pf:ks1:def-dtrace}
Let $\iota:\Gamma\to H$ be a homomorphism of groups, $K\subseteq H$ a subgroup, $\xi$ a representation of $\Gamma$ on a
complex vector space $V$, and $F:H\to\C$. The space of \emph{forms} is
$\mathcal M(\iota,K,\xi)=\{\phi:H\to V:\ \phi(\iota(\gamma)h)=\xi(\gamma)\phi(h),\ \phi(hk)=\phi(h)\}$, and
$\mathrm{htr}(\iota,K,\xi,F)$ is the trace of an endomorphism $T$ of $\mathcal M$ with
$(T\phi)(g)=\ksonefs_{h\in H/K}F(h)\phi(gh)$, if such an endomorphism exists (with a chosen one), and $0$ otherwise. For a
definite form $\beta$, a finite-dimensional representation $\xi$ of $C^\times\times\Q^\times$, a level $K$, $K'\subseteq G_p$
and $f,f'$, put $\mathrm{tr}_\beta(\xi;K,f;K',f')=\mathrm{htr}((\iota_f,\iota_p),K\times K',\xi|_{\Gamma_\beta},f\otimes f')$
and $\mathrm{Tr}^{\mathrm{def}}(\xi;K,f;K',f')=\ksonefs_{[\beta]\in\mathrm{Cl}^{\mathrm{def}}_{\mathrm{form}}}\mathrm{tr}_\beta(\xi;K,f;K',f')$.
\end{definition}

These are the algebraic modular forms of Gross \cite{Gross1999} on the definite inner forms $\Gamma_\beta$, with values in
$\xi$; $\mathrm{Tr}^{\mathrm{def}}$ is the trace of the Hecke operator of $f\otimes f'$ summed over the classes of definite forms.

\begin{definition}[Weights, good functions and the spectral condition]\label{pf:ks1:def-spec}
Let $D$, $G$, $L$ be as above and $\lambda:\Hom(F,\C)\to\N^n$.
\begin{enumerate}
\item For a monic $Q\in F[X]$ with $Q(0)\neq0$ let $Q^\flat=Q(0)^{-1}X^{\deg Q}Q(1/X)$ (the monic polynomial with inverse
  roots), and for $Q\in F_G[X]$ let ${Q}^\vee=(Q\circ e^{-1})^\flat\in F_D[X]$, $e:F_D\cong F_G$ the identification
  of $L$. For $\chi\in F[X]$ of degree $n$ let
  $\mathrm{S}_\lambda(\chi)=\prod_{\tau:F\to\C}s_{\lambda(\tau)}(r_{\tau,1},\dots,r_{\tau,n})$, where
  $r_{\tau,1},\dots,r_{\tau,n}$ are the roots of $\tau(\chi)$ with multiplicity and $s_\mu$ is the Schur polynomial given by
  the Jacobi--Trudi determinant $\det(h_{\mu_i-i+j})$ (it is symmetric, so the order of the roots is irrelevant).
\item A finite-dimensional representation $\xi$ of $C^\times\times\Q^\times$ has \emph{weight $\lambda$} if
  $\tr\xi(y)=\mathrm{S}_\lambda({\chi_y}^\vee)$ for every definite form $\beta$ and every $(y,s)\in\Gamma_\beta$
  with $\chi_y$ of degree $n$ and separable.
\item A function $f$ on $G(\ksoneAfp)$ is \emph{good} for a level $K$ if: $f$ is supported on regular elements;
  $\mathrm T_P(K,f;K_p,\mathbf 1_{K_p})=0$ for every marked pair $P$ with $\chi_P$ reducible; and for every definite form $\beta$
  (with any choice of $(x,t)$) and every $y\in\Gamma_\beta$ which is not regular irreducible, the orbital integral
  \[\mathrm{orb}_K(y)=\ksonefs_{q\in Z(y)\backslash(G(\ksoneAfp)\times G_p)/(K\times K_p)}\mathrm{vol}\bigl(Z(y)\cap qKq^{-1}\bigr)^{-1}
  (f\otimes\mathbf 1_{K_p})(q^{-1}yq)\]
  vanishes; here $Z(y)$ is the centraliser of $(\iota_f(y),\iota_p(y))$ in $G(\ksoneAfp)\times G_p$, $K$ stands for
  $K\times K_p$, and $\mathrm{vol}(U)=[U:U\cap U_0]/[U_0:U_0\cap U]$ with $U_0=Z(y)\cap K$.
\item For $\psi:J\to(G(\ksoneAfp)\to\C)$, a word $w=(j_1,\dots,j_r)$ in $J$ and $f$ put
  $\psi_w*f=\psi_{j_1}*(\psi_{j_2}*\cdots(\psi_{j_r}*f))$ (convolution for $K$), and for a noncommutative polynomial
  $s=\sum_ws_w\,w$ put $s(\psi)*f=\sum_ws_w\,\psi_w*f$ and $s(\theta)=\sum_ws_w\prod_{j\in w}\theta_j$ for $\theta:J\to\C$.
  The \emph{spectral condition} $\mathcal S(L,\lambda,K,\psi,\theta,f)$ means: $\psi_j*\psi_{j'}*\psi_w*f=\psi_{j'}*\psi_j*\psi_w*f$ for all $j,j',w$; and
  there is a finite-dimensional $\xi$ of weight $\lambda$ such that for every $s$ with $s(\theta)\neq0$ there is $a$ such that
  $(as)(\psi)*f$ is good for $K$ and $\mathrm{Tr}^{\mathrm{def}}(\xi;K,(as)(\psi)*f;K_p,\mathbf 1_{K_p})\neq0$.
\end{enumerate}
\end{definition}

In words: \emph{the eigensystem $\theta$ of the operators $\psi_j$ occurs in algebraic modular forms of weight $\lambda$ on the
definite inner forms, and this is detected by traces of good test functions}.

\begin{definition}[Hodge exponents and Lefschetz weights]\label{pf:ks1:def-hexp}
Let $F$ be a number field, $t_0:F\to\C$ with $\bar t_0\neq t_0$, $\varepsilon:\Hom(F,\C)\to\N$ and
$\lambda:\Hom(F,\C)\to\N^n$; write $|\lambda(\tau)|=\sum_j\lambda(\tau)_j$ and $\Sigma_\varepsilon=\sum_{\tau\neq t_0,\bar t_0}
\varepsilon(\tau)|\lambda(\tau)|$. A vector $h\in\N^n$ is a \emph{Hodge exponent vector} for $(t_0,\varepsilon,\lambda)$ if
either $\lambda(\bar t_0)$ is constant equal to some $m$ and $h_i=(n-1)m+\Sigma_\varepsilon+\lambda(t_0)_i+(n-1-i)$ for all
$0\le i<n$, or $\lambda(t_0)$ is constant equal to some $m$ and $h_i=m+\Sigma_\varepsilon+\sum_{j\neq i}\lambda(\bar t_0)_j+i$ for all
$i$. For $\chi\in F[X]$, $\mu\in\C$, $x\in F$ and $y\in\C$, with $r$ running over the roots of $t_0(\chi)$ (resp.\ $\tau(\chi)$)
with multiplicity and one copy of $\mu$ removed where indicated, put
\begin{align*}
\mathrm{ball}(\chi,\mu,y)&=\prod_{r\neq\mu}\frac{r-y\mu}{r-\mu},\\
\mathrm{fib}(x,\chi,\mu,y)&=\Bigl[(1-t_0(x)y\mu)\prod_{r\neq\mu}(1-t_0(x)r)\Bigr]\Bigl[(1-\overline{t_0(x)}\,\bar\mu)
  \prod_{r\neq\mu}(1-\overline{t_0(x)}\,y\bar r)\Bigr]\\
  &\qquad\times\prod_{\tau\neq t_0,\bar t_0}\ \prod_{r}(1-\tau(x)y^{\varepsilon(\tau)}r),\\
\mathrm{loc}(\alpha,\chi,\mu,y)&=\mathrm{ball}(\chi,\mu,y)\prod_{s=1}^k\mathrm{fib}(\alpha_s,\chi,\mu,y)\qquad(\alpha\in\OO_F^k).
\end{align*}
For a datum $D$ put $\varepsilon_D(\tau)=1$ if $\tau(\omega)=\bar\rho_c$ and $0$ otherwise. For a link $L$, $i_0<a$ and an infinite
place $v$ of $F_G$ put $\mathrm{mk}(\mu)=\mu^{-1}$ if $v(e(\omega))=\rho_c$ and $\mathrm{mk}(\mu)=\bar\mu^{-1}$ otherwise, and
\[c_W(\alpha,y)(Q,\mu)=\mathrm{loc}\bigl(\alpha,{Q}^\vee,\mathrm{mk}(\mu),y\bigr)\quad\text{with }t_0=\tau_{i_0},\ \varepsilon=\varepsilon_D.\]
For $\sigma\in\Aut(\C)$ let $(\sigma\lambda)(\tau)=\lambda(\sigma^{-1}\circ\tau)$.
\end{definition}

These are the Hodge degrees and the local holomorphic Lefschetz terms at complex multiplication points of the $k$-fold
Kuga--Sato variety over a unitary Shimura variety of signature $(1,n-1)$; the formal proofs use only the formulas.

\begin{definition}[$\sfname{FixK}$]\label{pf:ks1:def-fixk}
For a datum $D$ and a group $G$ let $\mathrm{SC}(D,G)$ mean: $a=1$, $m=1$, $a_G=1$, $n$ is odd and $D$ is of compact type. A
complex base $g_b$ is \emph{standard} if $\rho=\rho_c$ and $t_i=\tau_i|_{\OO_F}$. $\sfname{FixK}$ is the statement: let
$\mathrm{SC}(D,G)$ hold, $L$ a link, $g_b$ a standard complex base, $i_0<a$, $v$ the place of signature $(1,n-1)$, $k\ge0$, $T$
a Kuga--Sato tower for $(L,g_b,k)$, $N\mid N'$ levels with $Z_N$ smooth of relative dimension $d$, $h\in G(\ksoneAfp)$ such
that $1$ and $h$ are admissible for $(N',N)$ and $\mathbf 1_{K_Nh^{-1}K_N}$ is supported on regular elements, and
$\alpha\in\OO_F^k$ with every $\alpha_s$ nonzero and divisible by $2$. Then $\mathrm{Fin}(K_N,\mathbf 1_{K_Nh^{-1}K_N})$
holds, $c_2:=[\alpha]\circ\mathrm{tr}(N',N,h)$ is finite \'etale, and with $c_1=\mathrm{tr}(N',N,1)$ there are a finite
set $\mathcal F$ and tangent matrices $A_x$ as in \cref{pf:ks1:def-hlf} for $(c_1,c_2)$ on $Z_{N'}$, with
$\det(1-A_x)\neq0$ and
\begin{equation}\label{pf:ks1:eq-fixk}
\sum_{x\in\mathcal F}\mathrm{lef}(A_x,y)=[\mathrm{lev}(K_N,h):K_{N'}]\cdot\Sigma\bigl(c_W(\alpha,y);K_N,\mathbf 1_{K_Nh^{-1}K_N};K_p,\mathbf 1_{K_p}\bigr)
\qquad(y\in\C).
\end{equation}
\end{definition}

$\sfname{FixK}$ is proved in \lref{UpperTheory.Gal.FixZ.FixK_holds}.

\begin{definition}[$\sfname{KSHodge}$]\label{pf:ks1:def-kshodge}
$\sfname{KSHodge}$ is the statement: for every prime $p$, datum $D$, group $G$ and link $L$ with $\mathrm{SC}(D,G)$, and every
$\lambda:\Hom(F,\C)\to\N^n$ with each $\lambda(\tau)$ non-increasing, there is $k\ge0$ such that the following holds. Let $R$ be a
field of characteristic $0$, $g_b$ a base over $R$, $T$ a Kuga--Sato tower for $(L,g_b,k)$, $N$ a level, $(g_j)_{j\in J}$ Hecke data
of level $N$, $\theta:J\to\C$, and $f$ a function in the Hecke algebra of $K_N$ supported on elements with integral inverse, with
$\mathcal S(L,\lambda,K_N,(f_{g_j})_j,\theta,f)$. Let $\tau:R\to\C$ with $\tau(\rho)=\rho_c$ and $\tau\circ t_i=\tau_i$ on $\OO_F$,
let $\sigma\in\Aut(\C)$, $i_0<a$, and $h$ a Hodge exponent vector for $(\tau_{i_0},\varepsilon_D,\sigma\lambda)$. Then
$\mathrm{HO}(\mathrm{bc}_\tau(\mathcal H_T(N,g)),\sigma\circ\theta,h_i)$ holds for every $i$.
\end{definition}

\begin{theorem}[Hodge occurrence]\label{pf:ks1:thm-hodge}
$\sfname{HLF}$ and $\sfname{FixK}$ imply $\sfname{KSHodge}$.
\formallabel{UpperTheory.Gal.LocalSmoothExtension.ksHodge_of_geo}
\end{theorem}

The proof is given in \cref{pf:ks1:sec-proof}. In the formal development \cref{pf:ks1:thm-hodge} is applied with
$\sfname{HLF}$ proved in \lref{UpperTheory.Gal.Int.Geo.weight_one_geometric_realization} and $\sfname{FixK}$ proved in
\lref{UpperTheory.Gal.FixZ.FixK_holds}, inside \lref{UpperTheory.Gal.Int.Geo.kodaira_spencer_of_weight_one}.


% ---------------------------------------------------------------------------------------------
\subsection{The two counting identities}\label{pf:ks1:sec-counting}

Throughout this subsection $G$ is a unitary similitude group with $a_G=1$ and $v$ is its unique infinite place of
signature $(1,n-1)$.

\begin{lemma}[Unfolding orbit sums]\label{pf:ks1:lem-unfold}
Let a group $\Gamma$ act on a set $Y$ and let $\Phi:Y\times\Gamma\to\C$ satisfy $\Phi(\eta y,\eta\gamma\eta^{-1})=\Phi(y,\gamma)$.
Assume that $\Phi(y,\cdot)$ has finite support for every $y$ and that only finitely many orbits $o\in\Gamma\backslash Y$ have
$\Phi(y_o,\cdot)\neq0$, where $y_o$ is a chosen representative. Then
\[\ksonefs_{o\in\Gamma\backslash Y}\frac{1}{|\Gamma_{y_o}|}\ksonefs_{\gamma\in\Gamma}\Phi(y_o,\gamma)=
\ksonefs_{[\gamma]}\ \ksonefs_{o'\in Z(\gamma)\backslash Y}\frac{1}{|\Gamma_{y_{o'}}\cap Z(\gamma)|}\Phi(y_{o'},\gamma),\]
where $[\gamma]$ runs over the conjugacy classes of $\Gamma$ with chosen representatives $\gamma$, $Z(\gamma)$ is the centraliser,
$\Gamma_y$ the stabiliser, and $|S|^{-1}=0$ for infinite $S$.
\formallabel{GlobalTheory.T.orbit_sum_conj_classes}
\end{lemma}

\begin{proof}
Let $\Gamma$ act on $Y\times\Gamma$ by $\eta\cdot(y,\gamma)=(\eta y,\eta\gamma\eta^{-1})$; $\Phi$ is invariant, and both sides
are sums over the orbits of this action, grouped by the orbit of $y$ (left side) or by the class of $\gamma$ (right side),
of the quantity $|\Gamma_{(y,\gamma)}|^{-1}\Phi(y,\gamma)=|\Gamma_y\cap Z(\gamma)|^{-1}\Phi(y,\gamma)$. Indeed, the orbits lying over
the class $[\gamma]$ are in bijection with $Z(\gamma)\backslash Y$, which gives the right side. The orbits lying over the
orbit $\Gamma y$ are in bijection with the $\Gamma_y$-conjugacy classes in $\Gamma$; if $\Phi(y,\gamma)\neq0$ then the
$\Gamma_y$-conjugacy class of $\gamma$ lies in the finite support of $\Phi(y,\cdot)$, so $\Gamma_y\cap Z(\gamma)$ has finite
index in $\Gamma_y$, and the orbit--stabiliser formula gives
$\sum_{\gamma'\in\Gamma_y\cdot\gamma}\Phi(y,\gamma')=[\Gamma_y:\Gamma_y\cap Z(\gamma)]\,\Phi(y,\gamma)$. Dividing by $|\Gamma_y|$ (both
sides vanish if $\Gamma_y$ is infinite, since then $\Gamma_y\cap Z(\gamma)$ is infinite) and summing over the
$\Gamma_y$-classes gives $|\Gamma_y|^{-1}\sum_{\gamma}\Phi(y,\gamma)$, which is the left side. All sums are finite by the hypotheses.
\end{proof}

\begin{lemma}[Positive eigenlines]\label{pf:ks1:lem-markuniq}
Let $\beta$ be a form, $\gamma\in C^\times$ and $c\in\Q^\times$ with $\gamma\beta\gamma^*=c\beta$, and let $z,z'\in\C^n$ be eigenvectors of
$\psi_v(\gamma)$, with eigenvalues $\mu,\mu'$, which are positive for $\mathrm{herm}_v(\beta)$.
\begin{enumerate}
\item $\mu=\mu'$; in particular a pair is marked by at most one $\mu$.
\item If $\chi_\gamma$ has degree $n$ and is separable, then $z'\in\C z$.
\end{enumerate}
\formallabel{GlobalTheory.LKMid.Pair.mark_unique}
\formallabel{GlobalTheory.T.separable_formGroup_element_fixes_unique_positive_line}
\end{lemma}

\begin{proof}
Since $\beta$ is trivial at $\R$, the hermitian form $\langle z,w\rangle=z^\dagger\mathrm{herm}_v(\beta)w$ is equivalent to $\mathrm{gr}_v$, so it has signature
$(1,n-1)$. From $\gamma\beta\gamma^*=c\beta$ one gets $\langle\psi_v(\gamma)z,\psi_v(\gamma)w\rangle=c\langle z,w\rangle$; hence $|\mu|^2=c=|\mu'|^2$ and
$\bar\mu\mu'\langle z,z'\rangle=c\langle z,z'\rangle$. (1) If $\mu\neq\mu'$, then $\bar\mu\mu'\neq\bar\mu\mu=c$ (as $\mu\neq0$), so $\langle z,z'\rangle=0$; but the
orthogonal complement of a positive vector for a form of signature $(1,n-1)$ is negative definite, contradicting $\langle z',z'\rangle>0$. A marking of
a pair is the eigenvalue of a positive eigenvector at the unique place $v$ of signature $(1,n-1)$, so it is unique. (2) Every root of $v(\chi_\gamma)$ is an
eigenvalue of $\psi_v(\gamma)$ (\lref{GlobalTheory.LKCx.PolyArch.roots_eig_any}: every root of $v(\chi_\gamma)$ is an eigenvalue of $\psi_v(\gamma)$);
as $v(\chi_\gamma)$ has $n$ distinct roots, $\psi_v(\gamma)\in M_n(\C)$ has $n$ distinct eigenvalues and its eigenspaces are lines. By (1), $z$ and $z'$ lie
in the same eigenspace.
\end{proof}

\begin{lemma}[Finitely many contributing points]\label{pf:ks1:lem-finpts}
Let $K$ be a level, $f$ a test function for $K$ supported on regular elements, $K'\subseteq G_p$ commensurable with $K_p$ and $f'$ a test function
for $K'$. Let $\beta$ be a form with chosen $(x,t)$. Then the set of points $q\in\mathcal X_\beta$, with chosen representative $(g,z)$, for which some
$y\in\Gamma_\beta$ satisfies $\iota_v(y)z\in\C z$ and $f(g_f^{-1}\iota_f(y)g_f)\,f'(g_p^{-1}\iota_p(y)g_p)\neq0$ is finite.
\formallabel{GlobalTheory.LKCx.CxReg.finite_pts_reg}
\end{lemma}

\begin{proof}
Say that $y\in\Gamma_\beta$ \emph{contributes at} $q$ if it satisfies the two conditions of the statement for the representative of $q$. If $y$
contributes at some point, then $\iota_f(y)$ is conjugate to an element of the support of $f$, hence regular, and $y$ fixes the line of a vector of
$\mathcal P_\beta$; therefore $\chi_y$ has degree $n$ and is separable (\lref{GlobalTheory.LKCx.CxReg.sep_of_regAf}: regular elements of
$\Gamma_\beta$ fixing a positive line have separable minimal polynomial of degree $n$). We use three finiteness statements.
(a) The minimal polynomials of the elements contributing at some point form a finite set
(\lref{GlobalTheory.T.orbital_term_real_trivialisation}).
(b) For every polynomial $P$ there is a finite set $Y_P\subseteq\Gamma_\beta$ such that every $y$ with $\chi_y=P$ contributing at some point is
$\Gamma_\beta$-conjugate to an element of $Y_P$ (\lref{GlobalTheory.T.orbital_integral_test_function_first}).
(c) For every $y_0\in\Gamma_\beta$ only finitely many points $q$ admit a conjugate $\gamma y_0\gamma^{-1}$, $\gamma\in\Gamma_\beta$, contributing at $q$
(\lref{GlobalTheory.T.orbital_integral_test_function_second}).
Statements (b) and (c) take the separability shown above as a hypothesis. If $y$ contributes at $q$, then $\chi_y$ is one of the finitely many
polynomials $P$ of (a), $y=\gamma y_0\gamma^{-1}$ with $y_0\in Y_P$ by (b), and $q$ lies in the finite set of (c) for $y_0$. Hence the set in question
is contained in a finite union of finite sets.
\end{proof}

\begin{proposition}[Fixed points and marked pairs]\label{pf:ks1:prop-cxcount}
Let $c:F[X]\times\C\to\C$, $K$ a level, $f$ a test function for $K$ supported on regular elements, $K'\subseteq G_p$ commensurable
with $K_p$ and $f'$ a test function for $K'$. Then
$\Sigma(c;K,f;K',f')=\Sigma^{\mathrm{mk}}(c;K,f;K',f')$.
\formallabel{UpperTheory.Cb.complex_count_regular_marked_points}
\end{proposition}

\begin{proof}
Fix a class $[\beta]\in\mathrm{Cl}_{\mathrm{form}}$ with its chosen representative $\beta$ and $(x,t)$. For $y\in\Gamma_\beta$ put
\[\mathrm{PS}(y)=\ksonefs_{q\in Z(y)\backslash(G(\ksoneAfp)\times G_p)/(K\times K')}\frac{1}{|\mathrm{stab}_{Z(y)}(q)|}\,
f(q_f^{-1}\iota_f(y)q_f)\,f'(q_p^{-1}\iota_p(y)q_p),\]
where $Z(y)\subseteq\Gamma_\beta$ is the centraliser acting through $(\iota_f,\iota_p)$, and let $w(y)=c(\chi_y,\mu(y))\,\mathrm{PS}(y)$ if
$y$ fixes some line $\C z$ with $z\in\mathcal P_\beta$, $\mu(y)$ being its eigenvalue, and $w(y)=0$ otherwise. We show that both
sides equal $\ksonefs_{[\beta]}\ksonefs_{[y]}w(y)$, the inner sum over conjugacy classes of $\Gamma_\beta$.

(a) The group $\Gamma_\beta$ acts on $Y=(G(\ksoneAfp)\times G_p)/(K\times K')\times\mathcal P_\beta/\C^\times$ by left translation
through $(\iota_f,\iota_p)$ and through $\iota_v$; $\mathcal X_\beta=\Gamma_\beta\backslash Y$, the stabiliser of a point is the set
$\mathrm{Stab}$ of \cref{pf:ks1:def-forms}, and $\Sigma_\beta$ is the left side of \cref{pf:ks1:lem-unfold} for the function
$\Phi((gK,[z]),y)=f(g_f^{-1}\iota_f(y)g_f)f'(g_p^{-1}\iota_p(y)g_p)c(\chi_y,\mu_z(y))$ (zero unless $y$ fixes $[z]$), which is
invariant under simultaneous conjugation since $\chi$ and eigenvalues are conjugation invariant and $f,f'$ are invariant under
conjugation by $K$, $K'$ (\lref{UpperTheory.Cb.Mark.cxSumAtM_orbit}). The finiteness hypotheses of \cref{pf:ks1:lem-unfold}
hold: for fixed $(g,z)$ only finitely many $y$ fixing $[z]$ have $f(\dots)f'(\dots)\neq0$
(\lref{GlobalTheory.FinFix.finite_fix_core}; an integrality argument at the finite places and a boundedness argument at the
infinite places), and only finitely many points of $\mathcal X_\beta$
contribute (\cref{pf:ks1:lem-finpts}). By \cref{pf:ks1:lem-unfold},
$\Sigma_\beta=\ksonefs_{[y]}\ksonefs_{o\in Z(y)\backslash Y}|\mathrm{Stab}(y_o)\cap Z(y)|^{-1}\Phi(y_o,y)$. Fix $y$ and suppose some
$\Phi(\cdot,y)$ is nonzero. Then $\iota_f(y)$ is conjugate to an element of the support of $f$, hence regular, and $y$ fixes a positive
line; therefore $\chi_y$ has degree $n$ and is separable (\lref{GlobalTheory.LKCx.CxReg.sep_of_regAf}), and then $y$ fixes
exactly one line in $\mathcal P_\beta/\C^\times$ (\cref{pf:ks1:lem-markuniq}(2)).
Hence the $Z(y)$-orbits on $\{(gK,[z]):\Phi((gK,[z]),y)\neq0\}$ are the $Z(y)$-orbits of cosets $gK$, with $[z]$ the unique fixed line,
the stabilisers agree, and the inner sum is $c(\chi_y,\mu(y))\,\mathrm{PS}(y)=w(y)$
(\lref{GlobalTheory.T.conjugacyClassPart_eq_weight_mul_kottwitzPointSum}: for $y$ with a unique fixed positive line, the sum over
$Z(y)$-orbits of pairs (coset, fixed line) is the weight times the point sum of $y$).
Summing over $[\beta]$ gives $\Sigma=\ksonefs_{[\beta]}\ksonefs_{[y]}w(y)$.

(b) Let $P=(\beta',\gamma,c')$ be a marked pair. Then $c'>0$ ($|\mu|^2=c'$ as in the proof of \cref{pf:ks1:lem-markuniq}) and
$\beta'$ is a form; if $\beta$ is the chosen representative of its class, $\beta'=s\,g\beta g^*$, then
$y_P:=(g^{-1}\gamma g,c')\in\Gamma_\beta$, as $(g^{-1}\gamma g)\beta(g^{-1}\gamma g)^*=g^{-1}\gamma(s^{-1}\beta')\gamma^*g^{-*}=c'\beta$. The
conjugacy class of $y_P$ in $\Gamma_\beta$ depends only on the class of $P$, and the map $[P]\mapsto([\beta'],[y_P])$ is injective: if
$\eta y_P\eta^{-1}=y_{P'}$ with $\eta\in\Gamma_\beta$, $\beta'=s\,g\beta g^*$ and $\beta''=s'\,g'\beta g'^*$ for $P'=(\beta'',\gamma',c')$, then
$u=g'\eta g^{-1}$ satisfies $\gamma'=u\gamma u^{-1}$ and $\beta''=r\,u\beta'u^*$ for some $r\in\Q_{>0}$, so $P'$ is equivalent to $P$. Its image contains every class $[y]$ with $w(y)\neq0$ (if
$w(y)\neq0$ then, as in (a), $\chi_y$ has degree $n$ and is separable, so $(\beta,y,s)$ is a pair, marked by $\mu(y)$); moreover
$c(\chi_P,\mu_P)\,\mathrm T_P=w(y_P)$: the marking of $P$ is the eigenvalue on the positive fixed
line by \cref{pf:ks1:lem-markuniq}, and the term of $P$ is the point sum of its centraliser
(\lref{UpperTheory.Cb.Mark.clsW_regroup_of_conj_invariant}). Only finitely many marked classes have nonzero term
(\lref{GlobalTheory.T.finite_markedPairClasses_degreeOne_term_ne_zero}). Regrouping the finite sum $\Sigma^{\mathrm{mk}}$ along
this map gives $\Sigma^{\mathrm{mk}}=\ksonefs_{[\beta]}\ksonefs_{[y]}w(y)$.
\end{proof}

\begin{lemma}[Signs and local norms]\label{pf:ks1:lem-normsign}
Let $L$ be a number field with an automorphism $\sigma\neq\mathrm{id}$ such that $\phi\circ\sigma=\overline{\phi}$ for every embedding $\phi:L\to\C$, let
$S$ be a set of embeddings $L\to\C$ stable under complex conjugation with $4\mid|S|$, and let $p$ be a prime. Then there is $a\in L^\times$ with
$\sigma(a)=a$, $\phi(a)<0$ for $\phi\in S$, $\phi(a)>0$ for $\phi\notin S$, and $a=u\,\sigma(u)$ for some $u\in(\ksoneAfp\otimes L)^\times$.
\formallabel{GlobalTheory.T.exists_totallyReal_even_neg_places_local_norm}
\formallabel{GlobalTheory.NormSign.exists_neg_even_norm_away}
\end{lemma}

\begin{proof}
No embedding of $L$ is real: if $\phi$ were real, then $\phi\circ\sigma=\phi$ and $\sigma=\mathrm{id}$. We argue by induction on $|S|/4$. If $S=\emptyset$, take
$a=1$. Otherwise choose $\varepsilon\in S$ and $\varepsilon'\in S\smallsetminus\{\varepsilon,\bar\varepsilon\}$ (possible since $|S|\ge4$); then
$S'=S\smallsetminus\{\varepsilon,\bar\varepsilon,\varepsilon',\bar\varepsilon'\}$ is stable under conjugation and $|S'|=|S|-4$. By induction there is $b$ for $S'$.
There is $c\in L^\times$ with $\sigma(c)=c$, negative at $\varepsilon$ and $\varepsilon'$ (hence at their conjugates, as $c$ is fixed by $\sigma$), positive at all
other embeddings, and $c=u_c\sigma(u_c)$ with $u_c\in(\ksoneAfp\otimes L)^\times$
(\lref{GlobalTheory.T.exists_realSubfield_negative_twoPlaces_normEverywhere}: such elements exist for any two embeddings which are neither equal
nor conjugate). Then $a=bc$ has the required signs, is fixed by $\sigma$, and $a=(u_bu_c)\sigma(u_bu_c)$.
\end{proof}

\begin{lemma}[Moving definite pairs to marked pairs]\label{pf:ks1:lem-move}
Assume $n$ odd. Let $Q_0\in F[X]$ be irreducible and let $\mu$ be a root of $v(Q_0)$. There is $A\in F[X]$ with the following
properties.
\begin{enumerate}[label=(\roman*)]
\item For every definite pair $(\beta,\gamma,c)$ with $\chi_\gamma=Q_0$, the triple $(A(\gamma)\beta,\gamma,c)$ is a pair marked by $\mu$.
\item Every pair $P$ with $\chi_P=Q_0$ which is marked by $\mu$ is of the form $(A(\gamma)\beta,\gamma,c)$ for a definite pair $(\beta,\gamma,c)$.
\end{enumerate}
In both cases there is $U\in(\ksoneAfp\otimes C)^\times$ commuting with $\gamma$ with $U\beta U^*=A(\gamma)\beta$.
\formallabel{GlobalTheory.T.exists_polynomial_move_definiteSimilitudePair_marked}
\formallabel{GlobalTheory.LKDef.move_definite}
\end{lemma}

\begin{proof}
If there is neither a definite pair nor a pair marked by $\mu$ with polynomial $Q_0$, put $A=1$. Otherwise $Q_0$ is the
polynomial of such a pair, hence monic of degree $n$ and separable, and all complex roots of $Q_0$ (under all embeddings of $F$) have absolute value $c^{1/2}$, $c\in\Q_{>0}$ the similitude factor
of such a pair: for a definite pair this is the argument of \cref{pf:ks1:lem-roots}(1) at every infinite place, the form being definite there; for a
pair marked by $\mu$ the form is definite at the places other than $v$, and at $v$ the orthogonal complement of the positive eigenline is a negative
definite subspace stable under $\psi_v(\gamma)$, on which the same argument applies
(\lref{GlobalTheory.LKMid.MarkEll.mark_ell}: the roots of the polynomial of a marked pair have absolute value
$c^{1/2}$ under every embedding). Let $L=F[X]/(Q_0)$, a number field of degree $n$ over $F$; since
$|z|^2=c$ for every root $z$, complex conjugation of $F$ extends to an involution $\sigma$ of $L$ with $\sigma(X)=c/X$ which induces complex
conjugation under every embedding $L\to\C$ (\lref{GlobalTheory.T.cmField_conj_extends_of_adjoin_absSq}). Let $\varepsilon:L\to\C$ be the
embedding over $v$ with $\varepsilon(X)=\mu$, and let $S$ be the set of embeddings $\phi:L\to\C$ such that either $\phi$ restricts to $v$ on $F$
and $\phi\neq\varepsilon$, or $\phi$ restricts to $\bar v$ and $\phi\neq\bar\varepsilon$. Then $|S|=2(n-1)$, which is divisible by $4$ because
$n$ is odd (this is the only place where $n$ odd is used). Let $a\in L^\times$ and $u\in(\ksoneAfp\otimes L)^\times$ with $a=u\,\sigma(u)$ be given by \cref{pf:ks1:lem-normsign} for $(L,\sigma,S,p)$, and
let $A\in F[X]$ represent $a$. We show that $A$ has the stated properties.

Let $Q=(\beta,\gamma,c)$ be a definite pair with $\chi_Q=Q_0$, and let $j:L\to C$ be the $F$-embedding with $j(X)=\gamma$, so
$j(a)=A(\gamma)$. The element $j(a)\beta$ is again hermitian and $\gamma$ is a similitude of it with factor $c$, because $j(a)$ commutes with $\gamma$
and $\beta j(y)^*=j(\sigma y)\beta$ for $y\in L$. At the infinite places, the sign conditions on $a$ turn the definite form $\beta$ into a form
which is trivial at $\R$ and has, at $v$, a positive vector which is an eigenvector of $\psi_v(\gamma)$ with eigenvalue $\varepsilon(X)=\mu$
(\lref{GlobalTheory.LKDef.DMove.dreal_move}: after diagonalising $\gamma$ at each infinite place, multiplication by $j(a)$ multiplies the
diagonal entries of the hermitian form by the values $\phi(a)$). At the finite places away from $p$, the element $u$ gives
$U\in(\ksoneAfp\otimes C)^\times$ commuting with $j(L)$, hence with $\gamma$, and with $U\beta U^*=j(a)\beta$
(\lref{GlobalTheory.T.exists_adelicUnit_carrying_form_to_normMultiple}: a norm $a=u\sigma(u)$ from $\ksoneAfp\otimes L$ yields such a
$U$); in particular $j(a)\beta$ is trivial at the finite places away from $p$, since $\beta$ is. Thus $(A(\gamma)\beta,\gamma,c)$ is a pair marked by
$\mu$, which proves (i). Statement (ii) is proved in the same way, starting from $P$ and putting $\beta=j(a)^{-1}\beta_P$, with the inverse sign computation
(\lref{GlobalTheory.LKDef.DMove.real_back}: if $P$ is marked by $\mu$, then the hermitian form attached to $j(a)^{-1}\beta_P$ by
$\mathrm{herm}_w$ is positive definite at $v$ and negative definite at the other infinite places).
\end{proof}

\begin{proposition}[Marked pairs and definite pairs]\label{pf:ks1:prop-move}
Assume $n$ odd. Let $\mathrm{sel}$ be a root selection. There is a bijection $[P]\mapsto[Q]$ from the set of classes of pairs $P$
with $\chi_P$ irreducible and marked by $\mathrm{sel}(\chi_P)$ onto the set of classes of definite pairs $Q$ with $\chi_Q$ irreducible,
such that $\chi_Q=\chi_P$ and $\mathrm T_Q=\mathrm T_P$ (for every $K,f$ with $f$ invariant under conjugation by $K$, every $r>0$ and every
$K',\varphi$ with $\varphi$ invariant under twisted conjugation by $K'$; we only use $r=1$). Consequently, for every
$\mathrm{wt}:F[X]\to\C$, $\Sigma^{\mathrm{mk}}_{\mathrm{sel}}(\mathrm{wt};K,f;K',f')=\Sigma^{\mathrm{def}}(\mathrm{wt};K,f;K',f')$.
\formallabel{GlobalTheory.LKDef.cls_equiv_dcls_irr}
\formallabel{GlobalTheory.LKDef.marked_field_pairs_sum_eq_definite_pairs_sum}
\end{proposition}

\begin{proof}
For a monic irreducible $Q_0\in F[X]$ of degree $n$ let $A=A_{Q_0}$ be the polynomial of \cref{pf:ks1:lem-move} for $Q_0$ and
$\mu=\mathrm{sel}(Q_0)$, and let $U$ be as there. If $(\beta',\gamma')=(s\,y\beta y^*,y\gamma y^{-1})$, then $A(\gamma')\beta'=s\,y(A(\gamma)\beta)y^*$, because
$A(y\gamma y^{-1})=yA(\gamma)y^{-1}$; hence (i) and (ii) of \cref{pf:ks1:lem-move} are compatible with the equivalence relations and define mutually inverse maps on classes.
Equivalent pairs with irreducible polynomial have the same polynomial, so $\chi$ is preserved. Finally $\mathrm T_P$ depends on $\beta$ only through the
embedding $\iota_f$ of the centraliser, and $U$ identifies the embeddings for $P$ and $Q$ compatibly with $\gamma$, so $\mathrm T_P=\mathrm T_Q$
(\lref{GlobalTheory.T.kottwitzTerm_eq_definitePair_of_local_isometry}: terms of pairs with the same $(\gamma,c)$ whose forms are related by an adelic
unit away from $p$ commuting with $\gamma$ coincide). The last assertion follows by transporting the (finitely supported) sum along the bijection.
\end{proof}

\begin{proposition}[Fixed point sums and definite pairs]\label{pf:ks1:prop-h1}
Assume $n$ odd. Let $c:F[X]\times\C\to\C$, $K$ a level, $f$ a test function for $K$ supported on regular elements, $K'$ and $f'$ as in
\cref{pf:ks1:prop-cxcount}, and assume $\mathrm T_P(K,f;K',f')=0$ for every marked pair with reducible $\chi_P$. Then
\[\Sigma(c;K,f;K',f')=\Sigma^{\mathrm{def}}(\mathrm{rs}(c);K,f;K',f'),\qquad \mathrm{rs}(c)(Q)=\sum_{\mu}c(Q,\mu),\]
the sum over the roots $\mu$ of $v(Q)$ with multiplicity.
\formallabel{UpperTheory.Gal.Cnt.weighted_count_definite_unitary}
\formallabel{GlobalTheory.RootSumWeights.markWeightedFixedPoint_definitePair_odd_of_counts}
\end{proposition}

\begin{proof}
By \cref{pf:ks1:prop-cxcount}, $\Sigma(c)=\Sigma^{\mathrm{mk}}(c)$. Choose maps $\mathrm{sel}_0,\dots,\mathrm{sel}_{n-1}:F[X]\to\C$ such that
for every monic $Q$ of degree $n$ the roots of $v(Q)$, with multiplicity, are $\mathrm{sel}_0(Q),\dots,\mathrm{sel}_{n-1}(Q)$; each
$\mathrm{sel}_j$ is a root selection since $v$ is the only place of signature $(1,n-1)$. Let $P$ be a marked pair with $\chi_P$
irreducible. Its marking $\mu_P$ is a root of $v(\chi_P)$ and is unique (\cref{pf:ks1:lem-markuniq}); as $\chi_P$ is separable, there is
exactly one $j$ with $\mathrm{sel}_j(\chi_P)=\mu_P$, and $P$ is marked by $\mathrm{sel}_j(\chi_P)$ for that $j$ only. Marked pairs with
reducible $\chi_P$ contribute $0$ by hypothesis. Hence, the sums being finite
(\lref{GlobalTheory.RootSumWeights.finitely_many_marked_classes_term_degree_one}),
\[\Sigma^{\mathrm{mk}}(c)=\sum_{j<n}\Sigma^{\mathrm{mk}}_{\mathrm{sel}_j}\bigl(Q\mapsto c(Q,\mathrm{sel}_j(Q))\bigr)=
\sum_{j<n}\Sigma^{\mathrm{def}}\bigl(Q\mapsto c(Q,\mathrm{sel}_j(Q))\bigr)=\Sigma^{\mathrm{def}}(\mathrm{rs}(c)),\]
by \cref{pf:ks1:prop-move} (test functions for $K$ and $K'$ satisfy its invariance hypotheses) and since
$\sum_jc(Q,\mathrm{sel}_j(Q))=\mathrm{rs}(c)(Q)$ for the monic polynomials $\chi_Q$ of degree $n$; the finiteness of the set of definite
classes with nonzero summand, needed to exchange the sums, follows from that of marked classes via the bijection of
\cref{pf:ks1:prop-move}.
\end{proof}

\begin{lemma}[Trace formula for finite double coset spaces]\label{pf:ks1:lem-finitetrace}
Let $\iota:\Gamma\to H$ be an injective homomorphism, $K\subseteq H$ a subgroup, $\xi$ a finite-dimensional representation of $\Gamma$
and $F:H\to\C$ bi-$K$-invariant with support in finitely many cosets $hK$. Assume $\iota(\Gamma)\backslash H/K$ is finite and
$\Gamma_x=\iota^{-1}(xKx^{-1})$ is finite for every $x\in H$. Then
\[\mathrm{htr}(\iota,K,\xi,F)=\ksonefs_{[\gamma]}\tr\xi(\gamma)\ksonefs_{q\in\iota(Z(\gamma))\backslash H/K}
\frac{1}{|\iota(Z(\gamma))\cap qKq^{-1}|}F(q^{-1}\iota(\gamma)q).\]
\formallabel{GlobalTheory.HeckeTrace.heckeTrace_eq_orbital}
\end{lemma}

\begin{proof}
Let $x_1,\dots,x_r$ represent $\iota(\Gamma)\backslash H/K$ and $\Gamma_i=\Gamma_{x_i}$. A form $\phi$ is determined by the vectors
$\phi(x_i)$, which lie in $V^{\Gamma_i}$, and every choice of such vectors defines a form; so $\mathcal M\cong\bigoplus_iV^{\Gamma_i}$ is
finite-dimensional. The formula $(T\phi)(g)=\sum_{h\in H/K}F(h)\phi(gh)$ is a finite sum, independent of representatives, and defines
an endomorphism of $\mathcal M$ (right $K$-invariance follows from the left $K$-invariance of $F$ by reindexing $h\mapsto kh$); so
$\mathrm{htr}=\tr T$. The diagonal block of $T$ at $x_i$ is $v\mapsto\sum_hF(h)\xi(\gamma_h)v$, the sum over $hK$ with
$x_ih\in\iota(\gamma_h)x_iK$; such $\gamma_h$ form a coset $\gamma_h\Gamma_i$, and $h\mapsto\gamma$ identifies the pairs
$(hK,\gamma\in\gamma_h\Gamma_i)$ with the $\gamma\in\Gamma$, $hK=x_i^{-1}\iota(\gamma)x_iK$. With the projector
$e_i=|\Gamma_i|^{-1}\sum_{\delta\in\Gamma_i}\xi(\delta)$ onto $V^{\Gamma_i}$ the block is therefore
$|\Gamma_i|^{-1}\sum_{\gamma\in\Gamma}F(x_i^{-1}\iota(\gamma)x_i)\xi(\gamma)e_i$, and its trace is
$|\Gamma_i|^{-1}\sum_\gamma F(x_i^{-1}\iota(\gamma)x_i)\tr\xi(\gamma)$, because $F(x_i^{-1}\iota(\gamma\delta)x_i)=F(x_i^{-1}\iota(\gamma)x_i)$
for $\delta\in\Gamma_i$. Thus $\mathrm{htr}=\sum_{x}|\Gamma_x|^{-1}\sum_{\gamma}F(x^{-1}\iota(\gamma)x)\tr\xi(\gamma)$, the sum over
$x\in\iota(\Gamma)\backslash H/K$. Apply \cref{pf:ks1:lem-unfold} to $\Gamma$ acting on $H/K$ through $\iota$ and
$\Phi(xK,\gamma)=F(x^{-1}\iota(\gamma)x)\tr\xi(\gamma)$, which is invariant, has finite support in $\gamma$ (the support of $F$ meets
finitely many cosets, each giving a coset of the finite group $\Gamma_x$) and finitely many relevant orbits; the stabiliser of $xK$ is
$\Gamma_x$, and $Z(\gamma)$ corresponds to $\iota(Z(\gamma))$ since $\iota$ is injective.
\end{proof}

For a definite form $\beta$ (with its chosen $(x,t)$), a level $K$, a subgroup $K'\subseteq G_p$, functions $f$ on $G(\ksoneAfp)$ and $f'$ on $G_p$,
and $y\in\Gamma_\beta$, put
\[\mathrm{CS}_\beta(y)=\ksonefs_{q}\frac{1}{|\mathrm{stab}(q)\cap Z(y)|}\,f(q_f^{-1}\iota_f(y)q_f)\,f'(q_p^{-1}\iota_p(y)q_p),\]
the sum over $q\in\iota(Z(y))\backslash(G(\ksoneAfp)\times G_p)/(K\times K')$ with chosen representatives, where $\iota=(\iota_f,\iota_p)$, $Z(y)$ is the
centraliser of $y$ in $\Gamma_\beta$ and $\mathrm{stab}(q)=\{y'\in\Gamma_\beta:q^{-1}\iota(y')q\in K\times K'\}$.

\begin{lemma}[Finitely many definite classes contribute]\label{pf:ks1:lem-dfin}
Let $K$ be a level, $f$ a bi-$K$-invariant test function for $K$, $K'\subseteq G_p$ commensurable with $K_p$ and $f'$ a bi-$K'$-invariant test
function for $K'$. Then only finitely many classes in $\mathrm{Cl}^{\mathrm{def}}_{\mathrm{pair}}$ have a chosen representative $Q$ with $\chi_Q$
irreducible and $\mathrm T_Q(K,f;K',f')\neq0$.
\formallabel{GlobalTheory.T.bi_invariant_conjugacy_classes_finite}
\formallabel{GlobalTheory.LKDef.FinPairs.finite_classes_irr_term_ne_zero}
\end{lemma}

\begin{proof}
The polynomials $\chi_Q$ of definite pairs $Q$ with $\chi_Q$ irreducible and $\mathrm T_Q\neq0$ form a finite set
(\lref{GlobalTheory.T.finite_minpoly_definitePairs_term_ne_zero}), and for each polynomial $P_0$ only finitely many classes have a chosen
representative $Q$ with $\chi_Q=P_0$ irreducible and $\mathrm T_Q\neq0$ (\lref{GlobalTheory.T.finite_definitePair_classes_fixed_minpoly_degree_one}).
The set in question is contained in the union of these finite sets over the finitely many $P_0$.
\end{proof}

\begin{lemma}[Regrouping definite pairs]\label{pf:ks1:lem-regroup}
Let $\mathrm{wt}:F[X]\to\C$ and let $K,f,K',f'$ be as in \cref{pf:ks1:lem-dfin}. Then
\[\Sigma^{\mathrm{def}}(\mathrm{wt};K,f;K',f')=\ksonefs_{[\beta]\in\mathrm{Cl}^{\mathrm{def}}_{\mathrm{form}}}\ \ksonefs_{[y]}\mathrm{wt}(\chi_y)\,\mathrm{CS}_\beta(y),\]
the inner sum over the conjugacy classes of $\Gamma_\beta$ whose chosen representative $y$ is regular irreducible.
\formallabel{GlobalTheory.T.bi_invariant_test_function_sum}
\formallabel{GlobalTheory.LKDef.Regroup.regroup}
\formallabel{GlobalTheory.LKDef.Regroup.regroup_of}
\end{lemma}

\begin{proof}
Let $\mathcal I$ be the set of pairs $([\beta],[y])$ with $[\beta]\in\mathrm{Cl}^{\mathrm{def}}_{\mathrm{form}}$, with chosen representative $\beta$,
and $[y]$ a conjugacy class of $\Gamma_\beta$ whose chosen representative $y=(y,s)$ is regular irreducible. Then $(\beta,y,s)$ is a definite pair;
let $\kappa([\beta],[y])\in\mathrm{Cl}^{\mathrm{def}}_{\mathrm{pair}}$ be its class.

(a) $\kappa$ is injective. Let $\kappa([\beta_1],[y_1])=\kappa([\beta_2],[y_2])$, so $\beta_2=r\,u\beta_1u^*$ and $y_2=uy_1u^{-1}$ with $u\in C^\times$ and
$r\in\Q_{>0}$. Then $s_2\beta_2=y_2\beta_2y_2^*=r\,u(y_1\beta_1y_1^*)u^*=s_1\beta_2$, so $s_2=s_1$. The forms $\beta_1$, $\beta_2$ are
equivalent chosen representatives, so $\beta_1=\beta_2=:\beta$, and $(u,r^{-1})\in\Gamma_\beta$ conjugates $(y_1,s_1)$ to $(y_2,s_2)$; so $[y_1]=[y_2]$.

(b) The image of $\kappa$ contains every class of definite pairs with irreducible polynomial. Let $Q=(\beta',\gamma,c)$ be a definite pair with
$\chi_\gamma$ irreducible. First, $c>0$: the relation $\psi_w(x^*)=\mathrm{gr}_w^{-1}\psi_w(x)^\dagger\mathrm{gr}_w$ gives
$\mathrm{herm}_w(g\beta'g^*)=\psi_w(g^{-1})^\dagger\,\mathrm{herm}_w(\beta')\,\psi_w(g^{-1})$ for $g\in C^\times$, while
$\mathrm{herm}_w(c\beta')=c^{-1}\mathrm{herm}_w(\beta')$; for $g=\gamma$ and an infinite place $w$, at which $\mathrm{herm}_w(\beta')$ is definite,
this forces $c>0$. Let $\beta$ be the chosen representative of the class of $\beta'$, $\beta'=r\,u\beta u^*$. Then $y=(u^{-1}\gamma u,c)\in\Gamma_\beta$,
as $u^{-1}\gamma u\,\beta\,(u^{-1}\gamma u)^*=r^{-1}u^{-1}\gamma\beta'\gamma^*u^{-*}=c\beta$, and $y$ is regular irreducible. If $y_0=\eta y\eta^{-1}$, with
$\eta=(h,r')\in\Gamma_\beta$, is the chosen representative of $[y]$, then $\beta'=rr'\,w\beta w^*$ and $\gamma=wy_0w^{-1}$ for $w=uh^{-1}$, so $(\beta,y_0,c)$
is equivalent to $Q$ and $[Q]=\kappa([\beta],[y])$.

(c) For $([\beta],[y])\in\mathcal I$ the summand of $\Sigma^{\mathrm{def}}$ at $\kappa([\beta],[y])$ is $\mathrm{wt}(\chi_y)\,\mathrm{CS}_\beta(y)$: the
polynomial is $\chi_y$, the term of a definite pair does not change under equivalence (as $f$, $f'$ are invariant under conjugation by $K$, $K'$),
and the term of $(\beta,y,s)$ is $\mathrm{CS}_\beta(y)$ (\lref{GlobalTheory.T.definitePair_term_eq_centraliser_orbit_sum}: the term of a definite
pair $(\beta,y,s)$ equals the centraliser orbit sum of $y$ in $\Gamma_\beta$).

By \cref{pf:ks1:lem-dfin} the summand of $\Sigma^{\mathrm{def}}$ has finite support, which lies in the image of $\kappa$ by (b). By (a) and (c) the
finite sum can be reindexed along $\kappa$, which gives the sum of $\mathrm{wt}(\chi_y)\mathrm{CS}_\beta(y)$ over $\mathcal I$, a sum with finite
support; grouping it by $[\beta]$ gives the right-hand side.
\end{proof}

\begin{lemma}[Hecke traces on definite forms]\label{pf:ks1:lem-dtrace}
Let $\xi$ be a finite-dimensional representation of $C^\times\times\Q^\times$ and let $K,f,K',f'$ be as in \cref{pf:ks1:lem-dfin}. For every definite form
$\beta$ (with its chosen $(x,t)$),
\[\mathrm{tr}_\beta(\xi;K,f;K',f')=\ksonefs_{[y]}\tr\xi(y)\,\mathrm{CS}_\beta(y),\]
the sum over all conjugacy classes of $\Gamma_\beta$, with chosen representatives $y$.
\formallabel{GlobalTheory.T.heckeTrace_definiteForm_eq_sum_conjClasses}
\formallabel{GlobalTheory.HeckeTrace.dform_trace_eq}
\end{lemma}

\begin{proof}
Apply \cref{pf:ks1:lem-finitetrace} to $\Gamma=\Gamma_\beta$, $H=G(\ksoneAfp)\times G_p$, $\iota=(\iota_f,\iota_p)$, the subgroup $K\times K'$ and
$F=f\otimes f'$. The homomorphism $\iota$ is injective, since already $(y,s)\mapsto(xyx^{-1},s)$ is. The function $F$ is bi-$(K\times K')$-invariant
because $f$ and $f'$ are bi-invariant. The finiteness hypotheses hold: $F$ is supported on finitely many cosets $h(K\times K')$
(\lref{GlobalTheory.T.heckeTestProduct_support_finite_leftCosets}), the double coset space $\iota(\Gamma_\beta)\backslash H/(K\times K')$ is finite
(\lref{GlobalTheory.T.finite_doubleCoset_definiteHermitianForm_level}, finiteness of the class number of the definite group), and the groups
$\iota^{-1}(h(K\times K')h^{-1})$ are finite (\lref{GlobalTheory.T.definiteForm_stabilizer_finite_away_from_p}). The inner sums of
\cref{pf:ks1:lem-finitetrace} are the sums $\mathrm{CS}_\beta(y)$ (\lref{GlobalTheory.HeckeTrace.orbital_eq_centSum}: both are the same sum over
$\iota(Z(y))\backslash H/(K\times K')$).
\end{proof}

\begin{proposition}[Definite pairs and definite traces]\label{pf:ks1:prop-h2}
Let $\xi$ be a finite-dimensional representation of $C^\times\times\Q^\times$ and $\mathrm{wt}:F[X]\to\C$ such that
$\tr\xi(y)=\mathrm{wt}(\chi_y)$ for every definite form $\beta$ and every $(y,s)\in\Gamma_\beta$ with $\chi_y$ of degree $n$ and
separable. Let $K$ be a level, $f$ a bi-$K$-invariant test function for $K$, $K'$ commensurable with $K_p$, $f'$ a bi-$K'$-invariant
test function for $K'$, and assume that the orbital integral of $f\otimes f'$ (for $K\times K'$) at $(\iota_f(y),\iota_p(y))$ vanishes
for every definite form $\beta$ and every $y\in\Gamma_\beta$ which is not regular irreducible. Then
$\Sigma^{\mathrm{def}}(\mathrm{wt};K,f;K',f')=\mathrm{Tr}^{\mathrm{def}}(\xi;K,f;K',f')$.
\formallabel{UpperTheory.Gal.Cnt.trace_definite_forms_identity}
\formallabel{GlobalTheory.LKDef.trace_definite_forms_identity}
\end{proposition}

\begin{proof}
Let $\mathrm{CS}_\beta(y)$ be the sums defined before \cref{pf:ks1:lem-dfin}.
\begin{enumerate}
\item $\Sigma^{\mathrm{def}}(\mathrm{wt})=\ksonefs_{[\beta]\in\mathrm{Cl}^{\mathrm{def}}_{\mathrm{form}}}\ksonefs_{[y]}\mathrm{wt}(\chi_y)\,\mathrm{CS}_\beta(y)$, the inner sum
  over the conjugacy classes of $\Gamma_\beta$ consisting of regular irreducible elements, by \cref{pf:ks1:lem-regroup}.
\item $\mathrm{tr}_\beta(\xi;K,f;K',f')=\ksonefs_{[y]}\tr\xi(y)\,\mathrm{CS}_\beta(y)$, the sum over all conjugacy classes of $\Gamma_\beta$, by
  \cref{pf:ks1:lem-dtrace}.
\item If $y$ is not regular irreducible, $\mathrm{CS}_\beta(y)=0$: $\mathrm{CS}_\beta(y)$ is a rational multiple of the orbital integral of
  $f\otimes f'$ at $(\iota_f(y),\iota_p(y))$, which vanishes by hypothesis
  (\lref{GlobalTheory.T.centSum_eq_zero_of_orbitalIntegral_eq_zero}).
\item If $y$ is regular irreducible, then $\chi_y$ has degree $n$ and is separable, so $\tr\xi(y)=\mathrm{wt}(\chi_y)$.
\end{enumerate}
By (2)--(4), $\mathrm{tr}_\beta=\ksonefs_{[y]\ \mathrm{reg.\ irr.}}\mathrm{wt}(\chi_y)\mathrm{CS}_\beta(y)$; summing over $[\beta]$ and comparing
with (1) gives the assertion.
\end{proof}

\begin{remark}\label{pf:ks1:rem-counting}
\Cref{pf:ks1:prop-h1} is a formal counterpart of the comparison, going back to Kottwitz \cite{Kottwitz1990,Kottwitz1992}, between
fixed points of Hecke correspondences on unitary Shimura varieties and orbital integrals on an inner form which is definite at
infinity; \cref{pf:ks1:prop-h2} is the trace formula for a group which is compact at infinity, i.e.\ for the algebraic modular forms of
\cite{Gross1999}. Both are proved in the development from scratch, for the specific sums of \cref{pf:ks1:def-forms,pf:ks1:def-pairs}.
\end{remark}


% ---------------------------------------------------------------------------------------------
\subsection{Lefschetz weights and Schur polynomials}\label{pf:ks1:sec-schur}

Let $R$ be a commutative $\Q$-algebra. For $g\in R^n$, $y\in R$ and $i<n$ let $g^{(i,y)}$ be $g$ with $g_i$ replaced by $yg_i$, and
$g^{[i,y]}$ be $g$ with $g_k$ replaced by $yg_k$ for all $k\neq i$. Put $V(g)=\prod_{i<j}(g_i-g_j)$ and
$\det\nolimits_1(x,g)=\prod_j(1-xg_j)$. For $\mu\in\N^n$, $s_\mu$ is the Jacobi--Trudi determinant $\det(h_{\mu_i-i+j})$ of complete
symmetric polynomials (indices from $0$). We use the following standard facts \cite[Ch.~I, \S3]{Macdonald1995}, all proved in the
development: $s_\mu$ is symmetric (\lref{GroundTheory.Heights₂.Heights.RestTerms.Schur.schur_isSymmetric}); $s_0=1$ and
$s_{\mu+(m,\dots,m)}(g)=(\prod_kg_k)^m s_\mu(g)$ (\lref{UpperTheory.Gal.Schur.SchurDet_holds}); $s_\mu(yg)=y^{|\mu|}s_\mu(g)$ with
$|\mu|=\sum\mu_k$ (\lref{UpperTheory.Gal.Schur.schurHomog}); and the bialternant formula
\begin{equation}\label{pf:ks1:eq-bialt}
s_\mu(g)\,V(g)=\det\bigl(g_j^{\,\mu_i+n-1-i}\bigr)_{i,j<n}
\end{equation}
(\lref{UpperTheory.Gal.Schur.tie}). For a finite set $I$, $t_0\neq t_1\in I$, $\varepsilon:I\to\N$, a family $g=(g_t)_{t\in I}$ of
vectors in $R^n$, $y\in R$ and $i<n$ define the \emph{twisted family} $\mathrm{tw}_i(g)$ by $\mathrm{tw}_i(g)_{t_0}=g_{t_0}^{(i,y)}$,
$\mathrm{tw}_i(g)_{t_1}=g_{t_1}^{[i,y]}$ and $\mathrm{tw}_i(g)_t=y^{\varepsilon(t)}g_t$ otherwise.

\begin{lemma}[Determinants with scaled columns]\label{pf:ks1:lem-mixdet}
Let $M\in M_n(R)$ and $d,d'\in R^n$, and for $i<n$ let $M^{(i)}$ be the matrix whose $i$-th column is $\mathrm{diag}(d)M_{\cdot i}$ and whose
$j$-th column is $\mathrm{diag}(d')M_{\cdot j}$ for $j\neq i$. Then $\sum_i\det M^{(i)}=\bigl(\sum_ad_a\prod_{b\neq a}d'_b\bigr)\det M$.
\formallabel{UpperTheory.Gal.Schur.mixDet}
\end{lemma}

\begin{proof}
By the Leibniz formula, $\sum_i\det M^{(i)}=\sum_\sigma\mathrm{sgn}(\sigma)\prod_xM_{\sigma(x),x}\cdot\sum_id_{\sigma(i)}\prod_{x\neq i}d'_{\sigma(x)}$, and
the inner sum equals $\sum_ad_a\prod_{b\neq a}d'_b$ for every permutation $\sigma$ (put $a=\sigma(i)$).
\end{proof}

\begin{lemma}[Two Schur identities]\label{pf:ks1:lem-schury}
Let $g\in R^n$ be such that $g_k-g_i$ is a unit for $k\neq i$, let $\mu\in\N^n$ and $y\in R$. Then
\begin{align}
\sum_{i<n}s_\mu(g^{(i,y)})\prod_{k\neq i}\frac{g_k-yg_i}{g_k-g_i}&=s_\mu(g)\sum_{i<n}y^{\mu_i+n-1-i},\label{pf:ks1:eq-schury}\\
\sum_{i<n}s_\mu(g^{[i,y]})\prod_{k\neq i}\frac{g_i-yg_k}{g_i-g_k}&=s_\mu(g)\sum_{i<n}y^{\sum_{k\neq i}\mu_k+i}.\label{pf:ks1:eq-schuryoff}
\end{align}
\formallabel{UpperTheory.Gal.Schur.SchurY_holds}
\formallabel{UpperTheory.Gal.Schur.SchurYOff_holds}
\end{lemma}

\begin{proof}
Let $\ell_i=\mu_i+n-1-i$, $M(g)=(g_j^{\ell_i})_{i,j}$ and $d=(y^{\ell_i})_i$. Replacing $g_j$ by $yg_j$ multiplies the $j$-th column of
$M(g)$ by $\mathrm{diag}(d)$; so $M(g^{(i,y)})=M(g)^{\{i\}}$ and $M(g^{[i,y]})=M(g)^{\{i\}^c}$ in the notation of
\cref{pf:ks1:lem-mixdet}. Comparing the factors of $V$ which involve the index $i$ gives
\[V(g^{(i,y)})=V(g)\prod_{k\neq i}\frac{g_k-yg_i}{g_k-g_i},\qquad
V(g^{[i,y]})=y^{\binom{n-1}2}\,V(g)\prod_{k\neq i}\frac{g_i-yg_k}{g_i-g_k}.\] By \eqref{pf:ks1:eq-bialt} and \cref{pf:ks1:lem-mixdet},
$\sum_is_\mu(g^{(i,y)})V(g^{(i,y)})=\sum_i\det M(g)^{\{i\}}=(\sum_ky^{\ell_k})\,s_\mu(g)V(g)$; dividing by the unit $V(g)$ gives
\eqref{pf:ks1:eq-schury}. For \eqref{pf:ks1:eq-schuryoff} we first work in $R[Y]$ with $y=Y$. There
$\sum_is_\mu(g^{[i,Y]})V(g^{[i,Y]})=e_{n-1}(d)\,s_\mu(g)V(g)$ and
$e_{n-1}(d)=\sum_iY^{\sum_{k\neq i}\ell_k}=Y^{\binom{n-1}2}\sum_iY^{\sum_{k\neq i}\mu_k+i}$, since
$\sum_{k\neq i}(n-1-k)=\binom n2-(n-1-i)=\binom{n-1}2+i$. Dividing by $V(g)$ and cancelling the non-zero-divisor
$Y^{\binom{n-1}2}$ of $R[Y]$ gives \eqref{pf:ks1:eq-schuryoff} in $R[Y]$; evaluating at $Y=y$ gives it in $R$.
\end{proof}

\begin{lemma}[Twisted families]\label{pf:ks1:lem-fibre}
Let $I$, $t_0\neq t_1$, $\varepsilon$, $g$, $y$ be as above, $\lambda:I\to\N^n$, $m\in\N$, and put
$\Sigma_\varepsilon=\sum_{t\neq t_0,t_1}\varepsilon(t)|\lambda(t)|$.
\begin{enumerate}
\item If $\lambda(t_1)=(m,\dots,m)$ and the $g_{t_0,k}-g_{t_0,i}$ ($k\neq i$) are units, then
\[\sum_i\Bigl(\prod_ts_{\lambda(t)}(\mathrm{tw}_i(g)_t)\Bigr)\prod_{k\neq i}\frac{g_{t_0,k}-yg_{t_0,i}}{g_{t_0,k}-g_{t_0,i}}
=\prod_ts_{\lambda(t)}(g_t)\cdot y^{(n-1)m+\Sigma_\varepsilon}\sum_iy^{\lambda(t_0)_i+n-1-i}.\]
\item If $\lambda(t_0)=(m,\dots,m)$ and the $g_{t_1,k}-g_{t_1,i}$ ($k\neq i$) are units, then
\[\sum_i\Bigl(\prod_ts_{\lambda(t)}(\mathrm{tw}_i(g)_t)\Bigr)\prod_{k\neq i}\frac{g_{t_1,i}-yg_{t_1,k}}{g_{t_1,i}-g_{t_1,k}}
=\prod_ts_{\lambda(t)}(g_t)\cdot y^{m+\Sigma_\varepsilon}\sum_iy^{\sum_{k\neq i}\lambda(t_1)_k+i}.\]
\end{enumerate}
\formallabel{UpperTheory.Gal.Schur.Fibre_holds}
\formallabel{UpperTheory.Gal.Schur.FibreOff_holds}
\end{lemma}

\begin{proof}
For $t\neq t_0,t_1$, homogeneity gives $s_{\lambda(t)}(y^{\varepsilon(t)}g_t)=y^{\varepsilon(t)|\lambda(t)|}s_{\lambda(t)}(g_t)$. In case (1),
\[s_{\lambda(t_1)}(g_{t_1}^{[i,y]})=\Bigl(\prod_kg^{[i,y]}_{t_1,k}\Bigr)^m=y^{(n-1)m}s_{\lambda(t_1)}(g_{t_1}),\]
so the $i$-th summand is $y^{(n-1)m+\Sigma_\varepsilon}\prod_{t\neq t_0}s_{\lambda(t)}(g_t)$ times the $i$-th summand of
\eqref{pf:ks1:eq-schury} for $g_{t_0}$ and $\lambda(t_0)$, and \cref{pf:ks1:lem-schury} gives the claim. Case (2) is the same with
$s_{\lambda(t_0)}(g_{t_0}^{(i,y)})=y^ms_{\lambda(t_0)}(g_{t_0})$ and \eqref{pf:ks1:eq-schuryoff} for $g_{t_1}$.
\end{proof}

A subset $S\subseteq K^\sigma$ ($K$ a field, $\sigma$ finite) is \emph{Zariski dense} if the only polynomial in $K[X_s:s\in\sigma]$
vanishing on $S$ is $0$.

\begin{lemma}[Dense cosets]\label{pf:ks1:lem-dense}
Let $O$ be a ring, $K$ a field of characteristic $0$, $I$ a finite set and $(\tau_t)_{t\in I}$ pairwise distinct ring homomorphisms
$O\to K$. Let $N\in O$ with $\tau_t(N)\neq0$ for all $t$, $k\ge0$, $a\in O^k$ and $0\neq Q\in K[X_{t,s}:t\in I,s<k]$. Then the set of
points $x(\alpha)=(\tau_t(\alpha_s))_{t,s}$ with $\alpha\in O^k$, $\alpha\equiv a\bmod N$ and $Q(x(\alpha))\neq0$ is Zariski dense.
\formallabel{UpperTheory.Gal.Schur.DenseCoset_holds}
\end{lemma}

\begin{proof}
The vectors $(\tau_t(x))_t$, $x\in O$, span $K^I$: otherwise a nonzero linear form $\sum_tc_t\tau_t$ would vanish on $O$, contradicting the
linear independence of distinct monoid homomorphisms $(O,\cdot)\to(K,\cdot)$ (Dedekind--Artin). Choose $x_1,\dots,x_r\in O$
($r=|I|$) such that the vectors $(\tau_t(x_j))_t$ form a basis. Let $\Phi$ vanish at all $x(\alpha)$ with $\alpha\equiv a\bmod N$. For
$n=(n_{j,s})\in\Z^{r\times k}$ put $\alpha_s=a_s+N\sum_jn_{j,s}x_j$; then $\tau_t(\alpha_s)=\tau_t(a_s)+\tau_t(N)\sum_jn_{j,s}\tau_t(x_j)$, so
$\Phi\circ L$ vanishes on $\Z^{r\times k}$ for the invertible affine map $L:K^{r\times k}\to K^{I\times k}$ given by these formulas. A polynomial
over a field of characteristic $0$ vanishing on $\Z^{r\times k}$ is zero, hence $\Phi=0$: the coset is Zariski dense. Finally if $\Phi$
vanishes at the points of the coset with $Q(x(\alpha))\neq0$, then $\Phi Q$ vanishes on the whole coset, so $\Phi Q=0$ and $\Phi=0$.
\end{proof}

\begin{lemma}[Interpolation]\label{pf:ks1:lem-interp}
Let $O$, $K$, $I$, $(\tau_t)$ be as in \cref{pf:ks1:lem-dense}, $n\ge0$, and $P_t\in K[X_1,\dots,X_n]$ symmetric ($t\in I$). There is $k_0$
such that for every $k\ge k_0$, every $N\in O$ with all $\tau_t(N)\neq0$, every $a\in O^k$ and every $0\neq Q\in K[X_{t,s}]$, there is a
finitely supported $c:O^k\to K$, supported on $\alpha$ with $\alpha\equiv a\bmod N$ and $Q(x(\alpha))\neq0$, such that for every commutative
$K$-algebra $R$ and all $g=(g_t)_{t\in I}$, $g_t\in R^n$,
\begin{equation}\label{pf:ks1:eq-interp}
\sum_\alpha c_\alpha\prod_{t\in I}\prod_{s<k}\det\nolimits_1(\tau_t(\alpha_s),g_t)=\prod_{t\in I}P_t(g_t).
\end{equation}
\formallabel{UpperTheory.Gal.Schur.Interp_holds}
\formallabel{UpperTheory.Gal.Schur.interp_of}
\end{lemma}

\begin{proof}
(1) Every symmetric $P$ is a polynomial in the elementary symmetric polynomials $e_1,\dots,e_n$; writing each monomial as a product of
exactly $k$ factors $e_{m_1}\cdots e_{m_k}$ with $m_s\in\{0,\dots,n\}$ ($e_0=1$), which is possible as soon as $k$ is at least the degree,
we get, for all $k\ge k_0(P)$, $P=\sum_{m\in\{0..n\}^k}w(m)\prod_se_{m_s}$. Taking
$k_0=\max_tk_0(P_t)$ and multiplying over $t$, for $k\ge k_0$ there is $w:\{0..n\}^{I\times k}\to K$ with
$\prod_tP_t(g_t)=\sum_mw(m)\prod_{t,s}(-1)^{m_{t,s}}e_{m_{t,s}}(g_t)$ for all $R$ and $g$
(the signs are absorbed into $w$).

(2) Since $\det_1(x,g)=\sum_{j=0}^n(-x)^je_j(g)$, for $x\in K^{I\times k}$ we have
$\prod_{t,s}\det_1(x_{t,s},g_t)=\sum_mx^m\prod_{t,s}(-1)^{m_{t,s}}e_{m_{t,s}}(g_t)$, with $x^m=\prod_{t,s}x_{t,s}^{m_{t,s}}$.

(3) Let $S$ be the set of $\alpha$ in the statement; by \cref{pf:ks1:lem-dense} the points $x(\alpha)$, $\alpha\in S$, are Zariski dense. Then the
vectors $(x(\alpha)^m)_m\in K^{\{0..n\}^{I\times k}}$, $\alpha\in S$, span: a linear form $(u_m)\mapsto\sum_m\varphi_mu_m$ vanishing on them
gives the polynomial $\sum_m\varphi_mX^m$ (distinct exponent vectors give distinct monomials) vanishing on a dense set, so $\varphi=0$. Choose a finitely supported $c$ on $S$ with $\sum_\alpha c_\alpha x(\alpha)^m=w(m)$ for all $m$.

(4) By (2), (3) and (1), the left side of \eqref{pf:ks1:eq-interp} is
$\sum_m\bigl(\sum_\alpha c_\alpha x(\alpha)^m\bigr)\prod_{t,s}(-1)^{m_{t,s}}e_{m_{t,s}}(g_t)=\prod_tP_t(g_t)$.
\end{proof}

\begin{proposition}[Interpolation of Schur polynomials]\label{pf:ks1:prop-fibresum}
Let $E$ be a number field, $K$ a field of characteristic $0$ and $\lambda:\Hom(E,K)\to\N^n$. There is $k_0$ such that for every $k\ge k_0$,
every nonzero $N\in\OO_E$ and every finite $B\subseteq K$ there is a finitely supported $c:\OO_E^k\to K$, supported on $\alpha$ with
$\alpha_s\neq0$, $N\mid\alpha_s$ and $1-\tau(\alpha_s)b\neq0$ for all $s$, $\tau$ and $b\in B$, with the following property. For every
commutative $K$-algebra $R$, $t_0\neq t_1$ in $\Hom(E,K)$, $\varepsilon$, $m$, $g:\Hom(E,K)\to R^n$ and $y\in R$, the identities of
\cref{pf:ks1:lem-fibre} hold with $\prod_ts_{\lambda(t)}(\mathrm{tw}_i(g)_t)$ replaced by
$\sum_\alpha c_\alpha\prod_t\prod_s\det_1(t(\alpha_s),\mathrm{tw}_i(g)_t)$.
\formallabel{UpperTheory.Gal.Schur.FibreSum_holds}
\formallabel{UpperTheory.Gal.Schur.InterpNF_holds}
\end{proposition}

\begin{proof}
First we find $c$ with \eqref{pf:ks1:eq-interp} for $P_\tau=s_{\lambda(\tau)}$. If
$\Hom(E,K)=\emptyset$, take $c$ the indicator of $(N,\dots,N)$: both sides of \eqref{pf:ks1:eq-interp} are empty products. Otherwise apply
\cref{pf:ks1:lem-interp} with $O=\OO_E$, $I=\Hom(E,K)$, $\tau_t=t|_{\OO_E}$ (pairwise distinct since $\OO_E$ generates $E$), $P_t=s_{\lambda(t)}$,
$a=0$ and $Q=\prod_{t,s}X_{t,s}\prod_{b\in B}(1-X_{t,s}b)$, which is nonzero; for $\alpha$ in the support, $Q(x(\alpha))\neq0$ gives
$t(\alpha_s)\neq0$ (so $\alpha_s\neq0$) and $1-t(\alpha_s)b\neq0$, and $\alpha_s\equiv0\bmod N$. Now, exchanging the finite sums over $\alpha$ and $i$,
$\sum_\alpha c_\alpha\sum_i\prod_t\prod_s\det_1(t(\alpha_s),\mathrm{tw}_i(g)_t)C_i=\sum_i\prod_ts_{\lambda(t)}(\mathrm{tw}_i(g)_t)C_i$ for any
$C_i\in R$, by \eqref{pf:ks1:eq-interp} applied to the family $\mathrm{tw}_i(g)$; \cref{pf:ks1:lem-fibre} concludes.
\end{proof}

\begin{proposition}[Root sums of Lefschetz weights]\label{pf:ks1:prop-rootw}
Let $F$ be a number field, $t_0:F\to\C$ with $\bar t_0\neq t_0$, $\varepsilon:\Hom(F,\C)\to\N$, $\lambda:\Hom(F,\C)\to\N^n$, and $h$ a Hodge
exponent vector for $(t_0,\varepsilon,\lambda)$. There is $k_0$ such that for every $k\ge k_0$ there is a finitely supported
$c:\OO_F^k\to\C$, supported on $\alpha$ with every $\alpha_s$ nonzero and divisible by $2$, such that for every separable $\chi\in F[X]$ of degree
$n$ with $\chi(0)\neq0$ whose roots under $t_0$ all have the same absolute value, and every $y\in\C$,
\begin{equation}\label{pf:ks1:eq-rootw}
\sum_\alpha c_\alpha\sum_{\mu}\mathrm{loc}(\alpha,\chi,\mu,y)=\mathrm{S}_\lambda(\chi)\sum_{i<n}y^{h_i},
\end{equation}
the inner sum over the roots $\mu$ of $t_0(\chi)$ with multiplicity.
\formallabel{UpperTheory.Gal.LocalSmoothExtension.RootW_holds}
\formallabel{UpperTheory.Gal.LocalSmoothExtension.rootW_of}
\formallabel{UpperTheory.Gal.W2.RootSumW_holds}
\formallabel{UpperTheory.Gal.W2.RootSumWOff_holds}
\end{proposition}

\begin{proof}
Let $c$ be given by \cref{pf:ks1:prop-fibresum} for $E=F$, $K=\C$, $N=2$, $B=\emptyset$. Let $\chi$ be as in the statement and, for every $\tau$, let
$r_\tau=(r_{\tau,1},\dots,r_{\tau,n})$ list the roots of $\tau(\chi)$ with multiplicity (there are $n$ of them since $\C$ is
algebraically closed); they are pairwise distinct, and $\mathrm{S}_\lambda(\chi)=\prod_\tau s_{\lambda(\tau)}(r_\tau)$ (\lref{UpperTheory.Gal.LocalSmoothExtension.schWEq_of}).
Note $\bar t_0(x)=\overline{t_0(x)}$ and the roots of $\bar t_0(\chi)$ are the conjugates of those of $t_0(\chi)$.

\emph{First case: $\lambda(\bar t_0)=(m,\dots,m)$.} Put $t_1=\bar t_0$, $g_{t_1}=\bar r_{t_0}$
and $g_\tau=r_\tau$ for $\tau\neq t_1$. Since $g_{t_1}$ is a permutation of $r_{t_1}$ and Schur polynomials are symmetric,
$\prod_\tau s_{\lambda(\tau)}(g_\tau)=\mathrm{S}_\lambda(\chi)$. For the root $\mu=r_{t_0,i}$ we have
$\mathrm{ball}(\chi,\mu,y)=\prod_{k\neq i}(g_{t_0,k}-yg_{t_0,i})/(g_{t_0,k}-g_{t_0,i})$, and comparing factors,
$\mathrm{fib}(\alpha_s,\chi,\mu,y)=\prod_\tau\det_1(\tau(\alpha_s),\mathrm{tw}_i(g)_\tau)$: for $\tau=t_0$ the factor is
$(1-t_0(\alpha_s)y\mu)\prod_{k\neq i}(1-t_0(\alpha_s)r_{t_0,k})$, for $\tau=\bar t_0$ it is
$(1-\overline{t_0(\alpha_s)}\bar\mu)\prod_{k\neq i}(1-\overline{t_0(\alpha_s)}y\bar r_{t_0,k})$, and for the other $\tau$ it is
$\prod_r(1-\tau(\alpha_s)y^{\varepsilon(\tau)}r)$. Hence the left side of \eqref{pf:ks1:eq-rootw} is the left side of
\cref{pf:ks1:prop-fibresum}(1), and equals $\mathrm{S}_\lambda(\chi)\,y^{(n-1)m+\Sigma_\varepsilon}\sum_iy^{\lambda(t_0)_i+n-1-i}=
\mathrm{S}_\lambda(\chi)\sum_iy^{h_i}$ by the first case of the definition of $h$.

\emph{Second case: $\lambda(t_0)=(m,\dots,m)$.} Put $t_1=\bar t_0$, $g_{t_0}=\bar r_{t_1}$ and
$g_\tau=r_\tau$ for $\tau\neq t_0$; then $g_{t_0}$ lists the roots of $t_0(\chi)$, $g_{t_1,k}=\overline{g_{t_0,k}}$, and by hypothesis
$g_{t_0,k}g_{t_1,k}=q$ for a constant $q$, which is nonzero as $\chi(0)\neq0$. For $\mu=g_{t_0,i}$, substituting $g_{t_0,k}=q/g_{t_1,k}$,
\[\mathrm{ball}(\chi,\mu,y)=\prod_{k\neq i}\frac{q/g_{t_1,k}-yq/g_{t_1,i}}{q/g_{t_1,k}-q/g_{t_1,i}}=\prod_{k\neq i}\frac{g_{t_1,i}-yg_{t_1,k}}{g_{t_1,i}-g_{t_1,k}},\]
and $\mathrm{fib}(\alpha_s,\chi,\mu,y)=\prod_\tau\det_1(\tau(\alpha_s),\mathrm{tw}_i(g)_\tau)$ as in the
first case. Moreover $s_{\lambda(t_0)}(g_{t_0})=(\prod_kg_{t_0,k})^m=s_{\lambda(t_0)}(r_{t_0})$ since both products are the product of the roots of
$t_0(\chi)$; so $\prod_\tau s_{\lambda(\tau)}(g_\tau)=\mathrm{S}_\lambda(\chi)$. \Cref{pf:ks1:prop-fibresum}(2) and the second case of the definition of
$h$ give \eqref{pf:ks1:eq-rootw}.
\end{proof}

\begin{lemma}[Roots at the place of signature $(1,n-1)$]\label{pf:ks1:lem-roots}
Let $a=1$, $a_G=1$ and let $v$ be the place of signature $(1,n-1)$.
\begin{enumerate}
\item For a definite pair $P=(\beta,\gamma,c_P)$, $\chi_P(0)\neq0$ and every root $\mu$ of $v(\chi_P)$ satisfies $|\mu|^2=c_P$.
\item For $Q\in F_G[X]$ with $Q(0)\neq0$, the map $\mathrm{mk}$ carries the roots of $v(Q)$, with multiplicity, to the roots of
  $\tau_{i_0}(Q^\vee)$, with multiplicity.
\item If $Q\in F[X]$ has $Q(0)\neq0$, then $Q^\flat$ has the same degree, a nonzero constant term, and is separable if $Q$ is.
\end{enumerate}
\formallabel{UpperTheory.Gal.LocalSmoothExtension.dPairPol_of}
\formallabel{UpperTheory.Gal.LocalSmoothExtension.markRoots_of}
\formallabel{UpperTheory.Gal.LocalSmoothExtension.InvPolProps_holds}
\end{lemma}

\begin{proof}
(1) $\gamma$ is a unit, so $\chi_P(0)\neq0$. A root $\mu$ of $v(\chi_P)$ is an eigenvalue of $\psi_v(\gamma)$
(\lref{GlobalTheory.LKCx.PolyArch.roots_eig_any}: every root of $v(\chi_\gamma)$ is an eigenvalue of $\psi_v(\gamma)$); let $z\neq0$ be an eigenvector.
From $\gamma\beta\gamma^*=c_P\beta$ one gets $\langle\psi_v(\gamma)z,\psi_v(\gamma)z\rangle=c_P\langle z,z\rangle$ for the hermitian form
$\langle z,w\rangle=z^\dagger\mathrm{herm}_v(\beta)w$, which is definite, so $\langle z,z\rangle\neq0$ and $|\mu|^2=c_P$.
(2) By the link, $v$ is the place of an embedding $\phi:F_G\to\C$ with $\phi\circ e=\tau_{i_0}$, so the roots of $\tau_{i_0}(Q^\vee)$ are the inverses of the
roots of $\phi(Q)$ (\lref{UpperTheory.Gal.LocalSmoothExtension.InvRoots_holds}: the roots of $\tau(Q^\flat)$ are the inverses of those of $\tau(Q)$). If
$v$ is given by $\phi$, then $v(e(\omega))=\rho_c$ and $\mathrm{mk}(\mu)=\mu^{-1}$. Otherwise $v$ is given by $\bar\phi$, $v(e(\omega))=\bar\rho_c\neq\rho_c$
($\rho_c$ is not real since $\mathrm{tr}^2-4\,\mathrm{nm}<0$), the roots of $v(Q)$ are the conjugates of those of $\phi(Q)$, and $\mathrm{mk}(\mu)=\bar\mu^{-1}$.
(3) $Q^\flat$ is $Q(0)^{-1}$ times the reversed polynomial; reversal preserves the degree when $Q(0)\neq0$, its constant term is the leading coefficient
of $Q$, and its roots in an algebraic closure are the inverses of those of $Q$, so it is separable if $Q$ is.
\end{proof}

\begin{proposition}[The group side]\label{pf:ks1:prop-grouptop}
Let $D$, $G$, $L$ be as in \cref{pf:ks1:sec-groups}, $\lambda:\Hom(F,\C)\to\N^n$, $i_0<a$ and $h$ a Hodge exponent vector for
$(\tau_{i_0},\varepsilon_D,\lambda)$. There is $k_0$ such that for every $k\ge k_0$ there is a finitely supported $c:\OO_F^k\to\C$, supported on
$\alpha$ with every $\alpha_s$ nonzero and even, with the following property. If $a=1$, $a_G=1$, $n$ is odd, $v$ is the place of signature
$(1,n-1)$, $\xi$ has weight $\lambda$, $K$ is a level and $f$ is in the Hecke algebra of $K$ and good for $K$, then for all $y\in\C$
\begin{equation}\label{pf:ks1:eq-grouptop}
\Sigma\Bigl(\sum_\alpha c_\alpha\,c_W(\alpha,y);K,f;K_p,\mathbf 1_{K_p}\Bigr)=\mathrm{Tr}^{\mathrm{def}}(\xi;K,f;K_p,\mathbf 1_{K_p})\sum_{i<n}y^{h_i}.
\end{equation}
\formallabel{UpperTheory.Gal.LocalSmoothExtension.GroupTop_of}
\formallabel{UpperTheory.Gal.LocalSmoothExtension.groupTop_of}
\end{proposition}

\begin{proof}
Since $F$ is a CM field, $\tau_{i_0}$ is not real, so $\bar\tau_{i_0}\neq\tau_{i_0}$; let $k_0$ and $c$ be given by \cref{pf:ks1:prop-rootw} for
$t_0=\tau_{i_0}$, $\varepsilon=\varepsilon_D$. Assume the hypotheses of the second part. A function in the Hecke algebra of $K$ is a bi-$K$-invariant
test function for $K$, and $\mathbf 1_{K_p}$ is a bi-$K_p$-invariant test function for $K_p$. By goodness, $f$ is supported on regular elements
and the terms of marked pairs with reducible polynomial vanish, so \cref{pf:ks1:prop-h1} applies:
$\Sigma(\sum_\alpha c_\alpha c_W(\alpha,y))=\Sigma^{\mathrm{def}}(\mathrm{rs}(\sum_\alpha c_\alpha c_W(\alpha,y)))$. We claim that for every definite pair
$P$,
\begin{equation}\label{pf:ks1:eq-key}
\mathrm{rs}\Bigl(\sum_\alpha c_\alpha c_W(\alpha,y)\Bigr)(\chi_P)=\mathrm{S}_\lambda(\chi_P^\vee)\sum_iy^{h_i}.
\end{equation}
Let $\chi=\chi_P^\vee$. By \cref{pf:ks1:lem-roots}, $\chi$ has degree $n$, is separable, $\chi(0)\neq0$, and its roots under $\tau_{i_0}$ are the
$\mathrm{mk}(\mu)$, $\mu$ running over the roots of $v(\chi_P)$, so they all have absolute value $c_P^{-1/2}$. Hence \cref{pf:ks1:prop-rootw} applies
to $\chi$ and
\begin{align*}
\mathrm{rs}\Bigl(\sum_\alpha c_\alpha c_W(\alpha,y)\Bigr)(\chi_P)&=\sum_\alpha c_\alpha\sum_{\mu}\mathrm{loc}(\alpha,\chi,\mathrm{mk}(\mu),y)\\
&=\sum_\alpha c_\alpha\sum_{\mu'}\mathrm{loc}(\alpha,\chi,\mu',y)=\mathrm{S}_\lambda(\chi)\sum_iy^{h_i},
\end{align*}
with $\mu$ running over the roots of $v(\chi_P)$ and $\mu'$ over those of $\tau_{i_0}(\chi)$. This proves \eqref{pf:ks1:eq-key}. Since
$\Sigma^{\mathrm{def}}$ is a sum of $\mathrm{wt}(\chi_P)$ times terms, \eqref{pf:ks1:eq-key} gives
$\Sigma^{\mathrm{def}}(\mathrm{rs}(\dots))=\Sigma^{\mathrm{def}}(\mathrm{wt})\sum_iy^{h_i}$ with $\mathrm{wt}(Q)=\mathrm{S}_\lambda(Q^\vee)$. Finally,
since $\xi$ has weight $\lambda$, $\tr\xi(y')=\mathrm{wt}(\chi_{y'})$ for the relevant $y'$, and the orbital integrals of $f\otimes\mathbf 1_{K_p}$ at
elements which are not regular irreducible vanish by goodness; so \cref{pf:ks1:prop-h2} gives $\Sigma^{\mathrm{def}}(\mathrm{wt})=\mathrm{Tr}^{\mathrm{def}}(\xi;K,f;K_p,
\mathbf 1_{K_p})$.
\end{proof}


% ---------------------------------------------------------------------------------------------
\subsection{Proof of the Hodge occurrence theorem}\label{pf:ks1:sec-proof}

In this subsection $D$, $G$, $L$ satisfy the standing conditions $\mathrm{SC}(D,G)$ of \cref{pf:ks1:def-fixk}, and $T(f)$ denotes the
Hecke operator \eqref{pf:ks1:eq-hecke} of a Hecke function $f$ on the Hodge cohomology of a Kuga--Sato tower over $\C$. For a finitely
supported $c:\OO_F^k\to\C$ put $C_c=\sum_\alpha c_\alpha[\alpha]^*$.

\begin{lemma}[Projectivity and smoothness]\label{pf:ks1:lem-ksproj}
Let $g_b$ be a complex base, $T$ a Kuga--Sato tower for $(L,g_b,k)$ and $N$ a level. Then $Z_N\to\ksoneSp\C$ is projective and smooth of
some relative dimension.
\formallabel{UpperTheory.Gal.LocalSmoothExtension.KSProj_holds}
\end{lemma}

\begin{proof}
$Z_N$ is isomorphic over $\C$ to $\mathcal A_N^k\to M_N\to\ksoneSp\C$. The complex base defines a base for $D$ over $\C$ (the conditions at $p$
are automatic in characteristic $0$), so by \cref{pf:ks1:thm-models} there is a model of level $N$ over $\C$, and $M_N$ is an open and closed
subscheme of it (\lref{GlobalTheory.T.exists_open_closed_immersion_genModel_to_fibre}: the characteristic-$0$ fine moduli scheme is open
and closed in the corresponding fibre of the model). Since $D$ is of compact type, the model is projective
(\lref{GlobalTheory.UPel.HT.Model.isProjOver_of_compactB}: models of compact type over Noetherian domains of characteristic $0$ are
projective), hence so is $M_N$, and so is the fibre power $\mathcal A_N^k$ of the polarised abelian scheme $\mathcal A_N$
(\lref{GlobalTheory.KS.isProjOver_pow}: fibre powers of a polarised abelian scheme over a projective base are projective). For
smoothness, choose a prime $p'\nmid N$ at which the characteristic-$0$ datum underlying $D$ is a Harris--Taylor datum
(\lref{GlobalTheory.GoodPrime.exists_datum}); by \cref{pf:ks1:thm-models} for $p'$ and $R=\C$ the corresponding model is smooth, so its
open and closed subscheme $M_N$ is smooth; $\mathcal A_N\to M_N$ is smooth, and fibre powers of smooth morphisms of constant relative dimension
are smooth (\lref{GlobalTheory.FibPow.smoothOfRelativeDimension_pow}: the $k$-th fibre power of a smooth morphism of relative dimension $e$ is
smooth of relative dimension $ke$).
\end{proof}

\begin{proposition}[Geometric side]\label{pf:ks1:prop-geostep}
Assume $\sfname{HLF}$ and $\sfname{FixK}$. Let $g_b$ be a standard complex base, $i_0<a$, $v$ the place of signature $(1,n-1)$, $T$ a Kuga--Sato
tower for $(L,g_b,k)$, $N$ a level with $Z_N$ smooth of relative dimension $d$, $h$ an integral element such that $\mathbf 1_{K_Nh^{-1}K_N}$ is
supported on regular elements, and $\alpha\in\OO_F^k$ with every $\alpha_s$ nonzero and even. Then the finiteness condition
$\mathrm{Fin}(K_N,\mathbf 1_{K_Nh^{-1}K_N})$ holds and for all $y\in\C$
\[\mathcal L_d\bigl(T(\mathbf 1_{K_Nh^{-1}K_N})\circ[\alpha]^*\bigr)(y)=\Sigma\bigl(c_W(\alpha,y);K_N,\mathbf 1_{K_Nh^{-1}K_N};K_p,\mathbf 1_{K_p}\bigr).\]
\formallabel{UpperTheory.Gal.LocalSmoothExtension.geoStep_of}
\end{proposition}

\begin{proof}
By the tower axiom on deep levels applied to $(N,h)$ there is a level $N_1$ with $K_{N_1}\subseteq K_N$ and $hK_{N_1}h^{-1}\subseteq K_N$. Put
$N'=NN_1$; it is $\ge3$ and prime to $p$, $N\mid N'$, and $K_{N'}\subseteq K_{N_1}$ since principal congruence subgroups decrease along
divisibility. Hence $1$ and $h$ are admissible for $(N',N)$. Let $c_1=\mathrm{tr}(N',N,1)$, which is finite \'etale, and
$c_2=[\alpha]\circ\mathrm{tr}(N',N,h)$. $\sfname{FixK}$ gives the finiteness condition, the finite \'etaleness of $c_2$, the fixed point set
$\mathcal F$, tangent matrices $A_x$ with $\det(1-A_x)\neq0$, and \eqref{pf:ks1:eq-fixk}. By \cref{pf:ks1:lem-ksproj}, $Z_N$ is projective, so
$\sfname{HLF}$ applies to $(c_1,c_2)$ and gives, for every $a$,
$\sum_b(-1)^b\tr((c_1)_*c_2^*\mid\ksoneH{Z_N}{a}{b})=\sum_{x\in\mathcal F}[X^a]\det(1+XA_x)/\det(1-A_x)$. By \cref{pf:ks1:lem-hopop},
$T(\mathbf 1_{K_Nh^{-1}K_N})=\delta^{-1}(c_1)_*\mathrm{tr}(N',N,h)^*$ with $\delta=[\mathrm{lev}(K_N,h):K_{N'}]\neq0$, so by functoriality of pull-backs
$T(\mathbf 1_{K_Nh^{-1}K_N})\circ[\alpha]^*=\delta^{-1}(c_1)_*c_2^*$ on every $\ksoneH{Z_N}{a}{b}$. Multiplying by $(-y)^a$, summing over $a\le d$ and
using \cref{pf:ks1:lem-lefsum},
\[\mathcal L_d\bigl(T(\mathbf 1_{K_Nh^{-1}K_N})\circ[\alpha]^*\bigr)(y)=\delta^{-1}\sum_{x\in\mathcal F}\mathrm{lef}(A_x,y),\]
and \eqref{pf:ks1:eq-fixk} gives the assertion.
\end{proof}

\begin{proposition}[Hodge traces]\label{pf:ks1:prop-hodgetr}
Assume $\sfname{HLF}$ and $\sfname{FixK}$. Let $\lambda:\Hom(F,\C)\to\N^n$, $i_0<a$ and $h$ a Hodge exponent vector for
$(\tau_{i_0},\varepsilon_D,\lambda)$. There is $k_0$ such that for every $k\ge k_0$ there is a finitely supported $c:\OO_F^k\to\C$ with the
following property: for every standard complex base $g_b$, Kuga--Sato tower $T$ for $(L,g_b,k)$, level $N$ with $Z_N$ smooth of relative dimension
$d$, representation $\xi$ of weight $\lambda$, and Hecke function $f$ at level $N$ which is good for $K_N$,
\[\mathcal L_d\bigl(T(f)\circ C_c\bigr)(y)=\mathrm{Tr}^{\mathrm{def}}(\xi;K_N,f;K_p,\mathbf 1_{K_p})\sum_{i<n}y^{h_i}\qquad(y\in\C).\]
\formallabel{UpperTheory.Gal.LocalSmoothExtension.hodgeTr_of}
\end{proposition}

\begin{proof}
Let $k_0$ and $c$ be given by \cref{pf:ks1:prop-grouptop}. Let $v$ be the place of signature $(1,n-1)$; $K_N$ is a level
(\lref{GlobalTheory.LKGeo.KNLevel.exists_level_eq_KN}: principal congruence subgroups are compact open). The function $f$ is bi-$K_N$-invariant
with finitely many cosets in its support, so it is a finite linear combination $f=\sum_lc_l\mathbf 1_{K_Nh_l^{-1}K_N}$ of the indicator functions of the
double cosets in its support, with $h_l$ integral because $g^{-1}$ is integral for $g$ in the support. Since $f$ is supported on regular elements, so is
each $\mathbf 1_{K_Nh_l^{-1}K_N}$, and \cref{pf:ks1:prop-geostep} applies to every $h_l$ and every $\alpha$ in the support of $c$; with $\alpha=(2,\dots,2)$
it gives the finiteness conditions $\mathrm{Fin}(K_N,\mathbf 1_{K_Nh_l^{-1}K_N})$. The Hecke operator is linear in $f$ (\eqref{pf:ks1:eq-hecke} is linear
and the operator is unique), and $\mathcal L_d$ is linear in its argument; hence
\begin{align*}
\mathcal L_d(T(f)\circ C_c)(y)&=\sum_lc_l\sum_\alpha c_\alpha\,\mathcal L_d\bigl(T(\mathbf 1_{K_Nh_l^{-1}K_N})\circ[\alpha]^*\bigr)(y)\\
&=\sum_lc_l\sum_\alpha c_\alpha\,\Sigma\bigl(c_W(\alpha,y);K_N,\mathbf 1_{K_Nh_l^{-1}K_N}\bigr).
\end{align*}
Under the finiteness conditions all the sums defining $\Sigma(\,\cdot\,;K_N,\mathbf 1_{K_Nh_l^{-1}K_N})$ and $\Sigma(\,\cdot\,;K_N,f)$ run over subsets of one
finite set of triples, so $\Sigma$ is linear in the weight and in the test function
(\lref{UpperTheory.Gal.LocalSmoothExtension.CxLin_holds}), and the right side equals $\Sigma(\sum_\alpha c_\alpha c_W(\alpha,y);K_N,f)$. Now $f$ is in
the Hecke algebra of $K_N$ and good, so \cref{pf:ks1:prop-grouptop} gives the assertion.
\end{proof}

\begin{lemma}\label{pf:ks1:lem-hexp}
Let $t_0$, $\varepsilon$ be as in \cref{pf:ks1:def-hexp}.
\begin{enumerate}
\item If every $\lambda(\tau)$ is non-increasing, every Hodge exponent vector for $(t_0,\varepsilon,\lambda)$ is injective.
\item For every finite set $I$, maps $t:I\to\Hom(F,\C)$ and every function $\kappa$ of $(\lambda',i,h)$ with values in $\N$, there is $k$ with
  $\kappa(\sigma\lambda,i,h)\le k$ for all $\sigma\in\Aut(\C)$, $i\in I$ and Hodge exponent vectors $h$ for $(t(i),\varepsilon,\sigma\lambda)$.
\item If $\sum_{a=0}^d(-y)^at_a=c\sum_{i<n}y^{h_i}$ for all $y\in\C$, with $c\neq0$ and $h$ injective, then $t_{h_i}\neq0$ (and $h_i\le d$) for all $i$.
\end{enumerate}
\formallabel{UpperTheory.Gal.LocalSmoothExtension.hExpInj_of}
\formallabel{UpperTheory.Gal.LocalSmoothExtension.unifBound_of}
\formallabel{UpperTheory.Gal.LocalSmoothExtension.coeffOcc_of}
\end{lemma}

\begin{proof}
(1) For $i<k$: in the first case $h_i-h_k=\lambda(t_0)_i-\lambda(t_0)_k+(k-i)>0$; in the second case $h_i-h_k=\lambda(\bar t_0)_k-\lambda(\bar t_0)_i+(i-k)<0$.
(2) $\sigma\lambda=\lambda\circ\pi$ for the map $\pi=\sigma^{-1}\circ(\cdot)$ of the finite set $\Hom(F,\C)$ to itself, so $\sigma\lambda$ takes finitely many
values; for fixed $\lambda'$ and $i$ there are finitely many Hodge exponent vectors (the two formulas). Take the maximum of $\kappa$ over this finite
set. (3) Both sides are polynomials in $y$ which agree on $\C$; the coefficient of $y^{h_i}$ on the right is $c\neq0$ since $h$ is injective, and on the
left it is $(-1)^{h_i}t_{h_i}$ if $h_i\le d$ and $0$ otherwise.
\end{proof}

\begin{lemma}[Common eigenvectors]\label{pf:ks1:lem-eigof}
Let $B$ be a finite set, $W_b$ ($b\in B$) finite-dimensional complex vector spaces, and $T_{b,j}$ ($j\in J$), $u_b$ endomorphisms of $W_b$ such that
$T_{b,j}T_{b,j'}T_{b,w}u_b=T_{b,j'}T_{b,j}T_{b,w}u_b$ for all $j,j'$ and all words $w$ (where $T_{b,w}=T_{b,j_1}\cdots T_{b,j_r}$). For every
$\theta:J\to\C$ there is a noncommutative polynomial $s$ with $s(\theta)=1$ such that for every $b$ with $s(T_b)u_b\neq0$ there is $0\neq x\in W_b$
with $T_{b,j}x=\theta_jx$ for all $j$.
\formallabel{UpperTheory.Gal.LocalSmoothExtension.EigOf_holds}
\formallabel{UpperTheory.Gal.LocalSmoothExtension.eigOf_of}
\end{lemma}

\begin{proof}
Let $U_b\subseteq W_b$ be the span of the images of the $T_{b,w}u_b$. It is stable under every $T_{b,j}$, and by hypothesis the restrictions
$S_{b,j}=T_{b,j}|_{U_b}$ commute. The elements $S_j=(S_{b,j})_b$ of the finite-dimensional algebra $\prod_b\End(U_b)$ commute, so there is a
noncommutative polynomial $s$ with $s(\theta)=1$, $s(S)$ idempotent and $(S_j-\theta_j)^Ms(S)=0$ for all $j$ and some $M$
(\lref{UpperTheory.Gal.CertificateRingHomMap.pProjOne_holds}: such ``certificate'' polynomials exist for commuting elements of a finite-dimensional
algebra). If $s(T_b)u_bz\neq0$, then $y=s(S_b)(u_bz)\in U_b$ is nonzero and lies in the generalised $\theta_j$-eigenspace of every $S_{b,j}$. The
intersection of these generalised eigenspaces is a nonzero subspace stable under the commuting $S_{b,j}$, hence contains a common eigenvector $x$; its
eigenvalue for $S_{b,j}$ is $\theta_j$, since $x$ lies in the generalised $\theta_j$-eigenspace.
\end{proof}

\begin{proposition}[Hodge eigenvectors]\label{pf:ks1:prop-hodgeeig}
Assume $\sfname{HLF}$ and $\sfname{FixK}$. Let $\lambda:\Hom(F,\C)\to\N^n$ with each $\lambda(\tau)$ non-increasing. There is $k$ such that for every
standard complex base $g_b$, Kuga--Sato tower $T$ for $(L,g_b,k)$, level $N$, $\sigma\in\Aut(\C)$, $i_0<a$, Hodge exponent vector $h$ for
$(\tau_{i_0},\varepsilon_D,\sigma\lambda)$, family $\psi=(\psi_j)_{j\in J}$ of Hecke functions at level $N$, $\theta:J\to\C$ and Hecke function $f$ at level $N$
satisfying the spectral condition $\mathcal S(L,\sigma\lambda,K_N,\psi,\theta,f)$, and every $i<n$, there are $b$ and $0\neq x\in\ksoneH{Z_N}{h_i}{b}$ with
$T(\psi_j)x=\theta_jx$ for all $j$.
\formallabel{UpperTheory.Gal.LocalSmoothExtension.hodgeEig_of}
\end{proposition}

\begin{proof}
For $(\lambda',i_0,h)$ with $h$ a Hodge exponent vector for $(\tau_{i_0},\varepsilon_D,\lambda')$ let $\kappa(\lambda',i_0,h)$ be the $k_0$ of
\cref{pf:ks1:prop-hodgetr}; by \cref{pf:ks1:lem-hexp}(2) there is $k$ with $\kappa(\sigma\lambda,i_0,h)\le k$ for all $\sigma,i_0,h$. Fix the data and let
$c$ be given by \cref{pf:ks1:prop-hodgetr} for $(\sigma\lambda,i_0,h)$ and this $k$. By \cref{pf:ks1:lem-ksproj}, $Z_N$ is projective and smooth of some relative
dimension $d$, and by \cref{pf:ks1:lem-hodgefin} the spaces $W_b=\ksoneH{Z_N}{h_i}{b}$ are finite-dimensional. The Hecke operators are multiplicative
for convolution: $T(\psi_w*f)=T(\psi_{j_1})\cdots T(\psi_{j_r})T(f)$ and $T(s(\psi)*f)=s(T(\psi))T(f)$, all the functions involved being Hecke functions
(\lref{UpperTheory.Gal.LocalSmoothExtension.HopWord_holds} and \lref{UpperTheory.Gal.G3.pFnRing}: Hecke functions are stable under sums and convolution, and
the Hecke operator of a convolution is the composite). The first part of the spectral condition therefore gives the commutation hypothesis of
\cref{pf:ks1:lem-eigof} for $T_{b,j}=T(\psi_j)|_{W_b}$, $u_b=T(f)|_{W_b}$, $b\le d$; let $s$ be the resulting polynomial. As $s(\theta)=1\neq0$, the second part
of the spectral condition gives $\xi$ of weight $\sigma\lambda$ and $a$ such that $f'=(as)(\psi)*f$ is good for $K_N$ and
$\mathrm{Tr}^{\mathrm{def}}(\xi;K_N,f';K_p,\mathbf 1_{K_p})\neq0$. By \cref{pf:ks1:prop-hodgetr},
$\mathcal L_d(T(f')\circ C_c)(y)=\mathrm{Tr}^{\mathrm{def}}(\xi;K_N,f')\sum_iy^{h_i}$ for all $y$. Since $\sigma\lambda$ is non-increasing, $h$ is injective
(\cref{pf:ks1:lem-hexp}(1)), and \cref{pf:ks1:lem-hexp}(3) gives $\sum_b(-1)^b\tr(T(f')\circ C_c\mid W_b)\neq0$. Hence $T(f')\neq0$ on some $W_b$, $b\le d$.
As $T(f')=a(T(\psi))\,s(T(\psi))\,T(f)$, also $s(T_b)u_b\neq0$, and \cref{pf:ks1:lem-eigof} gives $x$.
\end{proof}

\begin{proposition}[Hodge occurrence]\label{pf:ks1:prop-kshodge}
The conclusion of \cref{pf:ks1:prop-hodgeeig} implies $\sfname{KSHodge}$.
\formallabel{UpperTheory.Gal.LocalSmoothExtension.kSHodge_of}
\end{proposition}

\begin{proof}
Let $\lambda$ be given, with each $\lambda(\tau)$ non-increasing, and let $k$ be given by \cref{pf:ks1:prop-hodgeeig}. Let $R$, $g_b$, $T$, $N$, $(g_j)$,
$\theta$, $f$, $\tau$, $\sigma$, $i_0$, $h$, $i$ be as in \cref{pf:ks1:def-kshodge}. Let $g_b'$ be the standard complex base $(\rho_c,(\tau_i)_i)$ (comaximality is
vacuous since $a=1$). We use four facts, proved in the development by routine verifications:
(a) the base change $T'$ of $T$ along $\tau$ is a Kuga--Sato tower for $(L,g_b',k)$ with the same levels, since $\tau(\rho)=\rho_c$ and $\tau\circ t_i=\tau_i$
(\lref{UpperTheory.Gal.LocalSmoothExtension.KSBc_holds}: Kuga--Sato towers are stable under base change of the field, with the base changed accordingly);
(b) for $\sigma\in\Aut(\C)$, $\mathcal S(L,\lambda,K,\psi,\theta,f)$ implies $\mathcal S(L,\sigma\lambda,K,\sigma\psi,\sigma\theta,\sigma f)$
(\lref{UpperTheory.Gal.LocalSmoothExtension.SpecConj_holds}: the $\sigma$-conjugate $\xi^\sigma$ of a representation of weight $\lambda$ has weight $\sigma\lambda$
since $\sigma(\mathrm{S}_\lambda(\chi))=\mathrm{S}_{\sigma\lambda}(\chi)$, traces of Hecke operators on forms satisfy
$\mathrm{Tr}^{\mathrm{def}}(\xi^\sigma;\sigma F)=\sigma(\mathrm{Tr}^{\mathrm{def}}(\xi;F))$, goodness is preserved and convolution commutes with $\sigma$);
(c) each $f_{g_j}$ is a Hecke function at level $N$, each $s_j$ is finite \'etale, and $T(f_{g_j})$ is the operator of the correspondence $(s_j,t_j)$
(\lref{UpperTheory.Gal.LocalSmoothExtension.MemOp_holds}: Hecke data define Hecke functions whose operators are the correspondences of the Hecke tower);
(d) a nonzero $x\in\ksoneH{Z'_N}{a}{b}$ with $T(f_{g_j})x=\vartheta_jx$ for all $j$ gives $\mathrm{HO}(\mathrm{bc}_\tau\mathcal H_T(N,g),\vartheta,a)$
(\lref{UpperTheory.Gal.LocalSmoothExtension.OccOf_holds}: by (c) and \lref{UpperTheory.Gal.G3.pOpTr}, $T(f_{g_j})=(s_j)_*t_j^*$ after base change).
The $f_{g_j}$ take integral values, so $\sigma\circ f_{g_j}=f_{g_j}$, and (b) gives $\mathcal S(L,\sigma\lambda,K_N,(f_{g_j}),\sigma\theta,\sigma f)$. The functions $f$ and
$\sigma f$ are in the Hecke algebra of $K_N$ and their supports consist of elements with integral inverse, so they are Hecke functions at level $N$ (the
monoid of the tower is the monoid of integral elements). \Cref{pf:ks1:prop-hodgeeig} applied to $T'$, $N$, $\sigma$, $i_0$, $h$, $\psi=(f_{g_j})$, $\sigma\theta$ and
$\sigma f$ gives $b$ and $0\neq x\in\ksoneH{Z'_N}{h_i}{b}$ with $T(f_{g_j})x=\sigma(\theta_j)x$, and (d) gives
$\mathrm{HO}(\mathrm{bc}_\tau\mathcal H_T(N,g),\sigma\circ\theta,h_i)$.
\end{proof}

\begin{proposition}[Assembly]\label{pf:ks1:prop-assembly}
$\sfname{HLF}$, \cref{pf:ks1:lem-ksproj}, $\sfname{FixK}$ and the identities of \cref{pf:ks1:prop-h1,pf:ks1:prop-h2} imply the conclusions of
\cref{pf:ks1:prop-hodgetr,pf:ks1:prop-hodgeeig} and $\sfname{KSHodge}$.
\formallabel{UpperTheory.Gal.LocalSmoothExtension.hodgeTrace_of_lefschetz_and_groupSide}
\formallabel{UpperTheory.Gal.LocalSmoothExtension.ksHodge_of_inputs}
\end{proposition}

\begin{proof}
\Cref{pf:ks1:prop-grouptop} uses \cref{pf:ks1:prop-h1,pf:ks1:prop-h2}; \cref{pf:ks1:prop-geostep} uses $\sfname{HLF}$, $\sfname{FixK}$ and
\cref{pf:ks1:lem-ksproj}; \cref{pf:ks1:prop-hodgetr} combines them; \cref{pf:ks1:prop-hodgeeig} uses in addition $\sfname{HLF}$ through
\cref{pf:ks1:lem-hodgefin}; and \cref{pf:ks1:prop-kshodge} concludes.
\end{proof}

\begin{proof}[Proof of \cref{pf:ks1:thm-hodge}]
Apply \cref{pf:ks1:prop-assembly}: \cref{pf:ks1:lem-ksproj,pf:ks1:prop-h1,pf:ks1:prop-h2} are proved above.
\end{proof}

% ---------------------------------------------------------------------------------------------
\subsection{Remarks on the formal statements}\label{pf:ks1:sec-remarks}

\begin{remark}[Conventions for sums, traces and operators]\label{pf:ks1:rem-conventions}
Several definitions use conventions of the formal development: a sum $\ksonefs$ with infinite support is $0$; $|S|^{-1}=0$ for an infinite set $S$; the trace
$\mathrm{htr}$ of \cref{pf:ks1:def-dtrace} is $0$ if no endomorphism with the stated formula exists; the Hecke operator $T(f)$ of \eqref{pf:ks1:eq-hecke} is $0$ if
no endomorphism with the stated property exists; and $\mathrm{S}_\lambda(\chi)$ is $0$ if the roots of $\chi$ cannot be listed (which does not happen for $\chi$ of degree
$n$). Consequently the fixed point sums $\Sigma$, the definite traces $\mathrm{Tr}^{\mathrm{def}}$ and the sums over pairs are defined for all test functions, and the
identities of \cref{pf:ks1:sec-counting,pf:ks1:sec-proof} need the finiteness statements cited there (and the finiteness condition $\mathrm{Fin}$ in
$\sfname{FixK}$) to be meaningful. The sums also depend on choices of representatives (of classes of forms, of points, of pairs) and of trivialisations
$(x,t)$; the statements hold for the choices made in the definitions.
\end{remark}

\begin{remark}[Shape of the Hodge occurrence statement]\label{pf:ks1:rem-kshodge}
(1) \Cref{pf:ks1:thm-hodge} is an implication: $\sfname{HLF}$ and $\sfname{FixK}$ are hypotheses, proved in
\lref{UpperTheory.Gal.Int.Geo.weight_one_geometric_realization} and \lref{UpperTheory.Gal.FixZ.FixK_holds}. (2) In $\sfname{KSHodge}$ the exponent $k$
depends only on $(p,D,G,L,\lambda)$, not on the tower, the level, the Hecke data, the test function, $\tau$ or $\sigma$. (3) The spectral condition is assumed
for $\lambda$, while the conclusion concerns the Hodge exponent vectors of every conjugate $\sigma\lambda$ and the conjugate eigensystem $\sigma\circ\theta$. (4) The
standing conditions require $a=m=a_G=1$, $n$ odd and compact type; in the proofs above $n$ odd enters only through
\cref{pf:ks1:lem-move} and compact type through \lref{GlobalTheory.UPel.HT.Model.isProjOver_of_compactB}.
The monotonicity of the $\lambda(\tau)$ is used only to make the Hodge exponents distinct. (5) If $\lambda$ is constant both at $t_0$ and at $\bar t_0$, both
alternatives of \cref{pf:ks1:def-hexp} apply and give two (possibly different) Hodge exponent vectors; the conclusion holds for each of them.
\end{remark}

\begin{remark}[The hypotheses on test functions]\label{pf:ks1:rem-good}
Goodness (\cref{pf:ks1:def-spec}) contains, besides regularity of the support, the vanishing of the terms of marked pairs with reducible characteristic
polynomial and of the orbital integrals at elements of the definite groups which are not regular irreducible. These are assumptions on the test function,
and they enter the spectral condition only through the functions $(as)(\psi)*f$. Similarly a representation of weight $\lambda$ is only constrained on
elements whose minimal polynomial has degree $n$ and is separable.
\end{remark}

\begin{remark}[Integral models]\label{pf:ks1:rem-models}
\Cref{pf:ks1:thm-models} makes no compactness assumption and allows any Noetherian base ring with base data; smoothness of relative dimension $a(n-1)$ is part
of the definition of a model. Kottwitz's condition is formulated through the traces in (D6) and, for objects, as the splitting and rank conditions (O3).
In (O4) and in the Siegel level structures the Weil pairing of the polarisation of the restriction of $P$ to the diagonal is compared with
$\zeta^{2E_\Lambda}$; the factor $2$ accounts for the fact that, for the pull-back $P=(1\times\lambda)^*\mathcal P$ of the Poincar\'e bundle by a
polarisation $\lambda$, the restriction of $P$ to the diagonal defines the polarisation $2\lambda$ \cite{Mumford1970}.
\end{remark}
